\section{Límites proyectivos de preesquemas}
\label{section:IV.8}

\subsection{Introducción}
\label{subsection:IV.8.1}

\oldpage[IV-3]{5}

\begin{env}[8.1.1]
\label{IV.8.1.1}
En este párrafo estudiaremos sistemáticamente la siguiente situación.
Sea $\mathrm{I}$ un conjunto creciente preordenado filtrado, $(\mathrm{A}_\lambda,\varphi_{\mu\lambda})$ un sistema inductivo de anillos cuyo conjunto de índices es $\mathrm{I}$, y sea $\mathrm{A}=\varinjlim\mathrm{A}_\lambda$ su límite inductivo.
Para cada $\alpha\in\mathrm{I}$ y cada $\mathrm{A}_\alpha$-preesquema $\mathrm{X}_\alpha$, consideremos los $\mathrm{A}_\lambda$-preesquemas $\mathrm{X}_\lambda=\mathrm{X}_\alpha\otimes_{\mathrm{A}_\alpha}\mathrm{A}_\lambda$ para $\lambda\geqslant\alpha$, y el $\mathrm{A}$-preesquema $\mathrm{X}=\mathrm{X}_\alpha\otimes_{\mathrm{A}_\alpha}\mathrm{A}$; está claro que los preesquemas $\mathrm{X}_\lambda$ (para $\lambda\geqslant\alpha$) forman un \emph{sistema proyectivo}, y veremos en \sref{IV.8.2.5} que $\mathrm{X}$ es un \emph{límite proyectivo} de este sistema en la categoría de preesquemas.
Nos proponemos encontrar condiciones sobre $\mathrm{X}_\alpha$ o sobre los $\mathrm{A}_\lambda$ que permitan demostrar propiedades del tipo siguiente:
para que $\mathrm{X}$ tenga una propiedad $\mathbf{P}$ (por ejemplo, ser propio sobre $\mathrm{S}=\operatorname{Spec}(\mathrm{A})$, o irreducible, o conexo, etc.), es necesario y suficiente que exista un índice $\mu\geqslant\alpha$ tal que, para todo $\lambda\geqslant\mu$, $\mathrm{X}_\lambda$ tenga (con respecto a $\mathrm{S}_\lambda=\operatorname{Spec}(\mathrm{A}_\lambda)$, según el caso) la misma propiedad $\mathbf{P}$.
Obtendremos enunciados análogos para propiedades de $\sh{O}_\mathrm{X}$-Módulos, de $\mathrm{A}$-morfismos de preesquemas $\mathrm{A}$, etc.
También mostraremos en \sref{IV.8.9.1} que el dato de un preesquema $\mathrm{A}$ de presentación finita \sref{IV.1.6.1} es esencialmente equivalente al dato de un preesquema $\mathrm{A}_\lambda$ de presentación finita $\mathrm{X}_\lambda$ para algún $\lambda$ suficientemente grande, siendo entonces $\mathrm{X}$ \emph{isomorfo a} $\mathrm{X}_\lambda\otimes_{\mathrm{A}_\lambda}\mathrm{A}$.
También existen enunciados análogos para $\sh{O}_\mathrm{X}$-Módulos de presentación finita, sus homomorfismos, los $\mathrm{A}$-morfismos de preesquemas $\mathrm{A}$ de presentación finita, etc.
\end{env}

\begin{env}[8.1.2]
\label{IV.8.1.2}
La utilidad de estos resultados aparecerá, por ejemplo, en las preguntas siguientes:
\begin{enumerate}[label=\emph{\alph*)}]
  \item
  Sea $\mathrm{Y}$ un preesquema, $y$ un punto de $\mathrm{Y}$, y $(\mathrm{U}_\lambda)$ el sistema proyectivo de vecindades afines abiertas filtradas decrecientes de $y$ en $\mathrm{Y}$; si $\mathrm{A}_\lambda$ es el anillo de $\mathrm{U}_\lambda$, los $\mathrm{A}_\lambda$ forman un sistema inductivo creciente filtrado, cuyo límite inductivo $\mathrm{A}$ es el anillo local $\sh{O}_y$; además, si denotamos por $\mathfrak{p}_\lambda$ el ideal primo $\mathfrak{j}_y$ en el anillo $\mathrm{A}_\lambda$, el sistema inductivo $(\mathrm{A}_\lambda)$ es cofinal en todo sistema inductivo $(\mathrm{A}_\alpha)_f$, donde $f$ recorre $\mathrm{A}_\alpha-\mathfrak{p}_\alpha$
\oldpage[IV-3]{6}
(para un $\alpha$ fijo), pues los $\mathrm{D}(f)$ forman un sistema fundamental de vecindades de $y$ en $\mathrm{U}_\alpha$ y, por tanto, en $\mathrm{Y}$.
Por consiguiente, los resultados de este párrafo implican que la geometría algebraica de los preesquemas $\sh{O}_y$ de presentación finita (y la teoría de los Módulos de presentación finita sobre estos preesquemas) es esencialmente equivalente a la geometría algebraica de los preesquemas de presentación finita en vecindades abiertas «suficientemente pequeñas» de $y$.
Así, el enunciado \sref{IV.8.10.5},~(xii) implica que, si un morfismo de presentación finita $f:\mathrm{X}\to\mathrm{Y}$ es tal que $\mathrm{X}\times_\mathrm{Y}\operatorname{Spec}(\sh{O}_y)$ es propio sobre $\operatorname{Spec}(\sh{O}_y)$, es necesario y suficiente que exista una vecindad abierta $\mathrm{U}$ de $y$ en $\mathrm{Y}$ tal que $f^{-1}(\mathrm{U})$ sea propio sobre $\mathrm{U}$.

Un caso particular importante, clásico en cierto sentido, es aquel en que $\mathrm{Y}$ es íntegro y $y$ es su punto genérico, de modo que $\sh{O}_y$ no es otra cosa que \emph{el cuerpo} $\mathrm{R}(\mathrm{Y})=\mathrm{K}$ \emph{de funciones racionales sobre} $\mathrm{Y}$.
Los resultados de este párrafo equivalen entonces a interpretar la geometría algebraica sobre $\mathrm{K}$ en términos de geometría algebraica sobre subconjuntos abiertos no vacíos «suficientemente pequeños» de $\mathrm{Y}$, es decir, intuitivamente, en términos de «familias» de objetos geométricos indexadas por los puntos de tal abierto.
Además, este punto de vista se ha utilizado habitualmente desde hace mucho tiempo, no solo en geometría algebraica sobre cuerpos algebraicamente cerrados, sino también en el estudio aritmético de variedades definidas sobre un \emph{cuerpo de números} $\mathrm{K}$ (una extensión finita de $\mathbf{Q}$), considerando este último como el cuerpo de fracciones de su anillo de enteros $\mathrm{A}$ (la «teoría de la \emph{reducción módulo} $\mathfrak{p}$» ($\mathfrak{p}$ un ideal primo de $\mathrm{A}$); cf.\ \sref[I]{I.3.7}).
Los resultados de las secciones \textsection\textsection\,8 y~9 proporcionan así, entre otras cosas, algunos fundamentos del lenguaje de la «reducción módulo $\mathfrak{p}$» en aritmética.

Observamos que, en el ejemplo considerado aquí, los morfismos $\mathrm{S}_\mu\to\mathrm{S}_\lambda$ (para $\lambda\leqslant\mu$) son las \emph{inmersiones abiertas} canónicas $\mathrm{U}_\mu\hookrightarrow\mathrm{U}_\lambda$ y, \emph{a fortiori}, son morfismos \emph{planos} (aunque en general no fielmente planos), lo que explica el interés de los enunciados que recurren a tal restricción.

  \item
  Supongamos que los $\mathrm{A}_\lambda$ son \emph{cuerpos}, de modo que $\mathrm{A}=\varinjlim\mathrm{A}_\lambda$ también es un cuerpo.
Este caso aparece generalmente cuando se parte de datos geométricos sobre un cuerpo arbitrario $\mathrm{K}$, que se considera como extensión de un cuerpo $k$ (por ejemplo, el subcuerpo primo de $\mathrm{K}$).
Resulta entonces útil considerar $\mathrm{K}$ como límite inductivo de sus subextensiones que son \emph{de tipo finito sobre $k$}, lo que permite reducir muchas cuestiones al caso en que $\mathrm{K}$ es una extensión de tipo finito de $k$.
Usando también el método esbozado en \emph{a)}, se puede reducir en general al caso de un anillo de base $\mathrm{A}$ que sea una \emph{álgebra íntegra de tipo finito sobre $k$}.

Observamos que, en este ejemplo, los morfismos $\mathrm{S}_\mu\to\mathrm{S}_\lambda$ son \emph{fielmente planos}.

  \item
  Supongamos que nos interesan propiedades, locales sobre $\mathrm{Y}$, de preesquemas de presentación finita sobre un preesquema arbitrario $\mathrm{Y}$, que en adelante podemos suponer un esquema afín de anillo $\mathrm{A}$.
Resulta entonces útil considerar $\mathrm{A}$ como límite inductivo de sus subanillos que son \emph{$\mathbf{Z}$-álgebras de tipo finito}, lo que permite reducir muchas cuestiones al caso en que $\mathrm{Y}$ es el espectro de una de tales álgebras.
Esta es la explicación del «punto de vista kroneckeriano», según el cual la geometría algebraica se reduce a la
\oldpage[IV-3]{7}
geometría algebraica de preesquemas de tipo finito sobre $\mathbf{Z}$ (que a veces se califica de «geometría algebraica absoluta»).
Este ejemplo muestra, en particular, que en la mayoría de las cuestiones «relativas» a un preesquema de base $\mathrm{Y}$ se puede reducir al caso en que $\mathrm{Y}$ es \emph{noetheriano}.

Observamos que, en este ejemplo, a diferencia de los anteriores, los morfismos $\mathrm{S}_\mu\to\mathrm{S}_\lambda$ no tienen en general ninguna propiedad particular de regularidad.
\end{enumerate}

En lo que sigue, al aplicar los resultados que siguen a cualquiera de las tres situaciones particulares que acabamos de describir, no entraremos en detalles sobre el procedimiento que permite realizar estas aplicaciones, y nos limitaremos a remitirnos a lo anterior.
\end{env}

\begin{env}[8.1.3]
\label{IV.8.1.3}
En el ejemplo \emph{a)} de \sref{IV.8.1.2} vimos que, si $\mathrm{Y}$ es un preesquema íntegro con punto genérico $y$, $f:\mathrm{X}\to\mathrm{Y}$ un morfismo de presentación finita, y la fibra genérica $f^{-1}(y)$ es propia sobre $\boldsymbol{k}(y)$, entonces existe una vecindad abierta $\mathrm{U}$ de $y$ en $\mathrm{Y}$ tal que $f^{-1}(\mathrm{U})$ es propia sobre $\mathrm{U}$; se deduce \emph{a fortiori} que, para todo $s\in\mathrm{U}$, $f^{-1}(s)$ es propia sobre $\boldsymbol{k}(s)$.
Resulta que se necesita un recíproco que asegure que, si $f^{-1}(s)$ es propia sobre $\boldsymbol{k}(s)$ para «suficientemente muchos» puntos $s\in\mathrm{U}$, entonces $f^{-1}(y)$ es propia sobre $\boldsymbol{k}(y)$.
Por ejemplo, supongamos que $\mathrm{X}$ y $\mathrm{Y}$ son preesquemas algebraicos sobre un cuerpo algebraicamente cerrado $k$ (para fijar ideas, se puede tomar como $k$ el cuerpo $\mathbf{C}$ de los números complejos); a veces se necesita saber que, si para todo $s\in\mathrm{Y}$, \emph{racional sobre $k$}, la fibra $f^{-1}(s)$ es propia sobre $\boldsymbol{k}(s)$, entonces $f^{-1}(y)$ es propia sobre $\boldsymbol{k}(y)$ y, por consiguiente, $f^{-1}(\mathrm{U})$ es propia sobre $\mathrm{U}$ para una vecindad $\mathrm{U}$ de $y$\footnote{Observamos que tal enunciado es, en último término, de naturaleza puramente geométrica, pues solo recurre a puntos racionales sobre $k$, y no a puntos genéricos; por ejemplo, cuando $k=\mathbf{C}$, este enunciado tiene un significado topológico evidente para el analista, interpretando «propio» en el sentido topológico del término para los espacios analíticos subyacentes formados por los puntos de $\mathrm{X}$ y $\mathrm{Y}$ racionales sobre $\mathbf{C}$.}.
Ahora bien, este enunciado se deduce fácilmente de lo siguiente: el conjunto $\mathrm{E}$ de puntos $s\in\mathrm{Y}$ tales que $f^{-1}(s)$ es propio sobre $\boldsymbol{k}(s)$ es \emph{constructible} (y, por tanto, coincide con todo $\mathrm{Y}$ si contiene los puntos cerrados de $\mathrm{Y}$, por el teorema de los ceros de Hilbert \sref[I]{I.10.4.8}); esto equivale una vez más a decir que, si $f^{-1}(y)$ \emph{no es propio} sobre $\boldsymbol{k}(y)$, existe una vecindad abierta $\mathrm{U}$ de $y$ tal que $f^{-1}(s)$ no es propio sobre $\boldsymbol{k}(s)$ para todo $s\in\mathrm{U}$ (cf.\ \sref{IV.9.6.1},~(iv))).
Este ejemplo ilustra el interés de desarrollar sistemáticamente \emph{criterios de constructibilidad}: esto es lo que se hará en la \textsection\,9.
\end{env}

\subsection{Límites proyectivos de preesquemas}
\label{subsection:IV.8.2}

\begin{env}[8.2.1]
\label{IV.8.2.1}
Sea $\mathrm{S}_0$ un espacio anillado, $\mathrm{L}$ un conjunto creciente preordenado filtrado, $(\sh{A}_\lambda,\varphi_{\mu\lambda})$ un sistema inductivo de $\sh{O}_{\mathrm{S}_0}$-Álgebras (no necesariamente conmutativas), cuyo conjunto de índices es $\mathrm{L}$.
Sabemos que, considerado como sistema inductivo de $\sh{O}_{\mathrm{S}_0}$-Módulos, $(\sh{A}_\lambda,\varphi_{\mu\lambda})$ admite un límite inductivo $\sh{A}$; denotemos por $\varphi_\lambda:\sh{A}_\lambda\to\sh{A}$ el homomorfismo canónico de $\sh{O}_{\mathrm{S}_0}$-Módulos.
Sea $m_\lambda:\sh{A}_\lambda\otimes\sh{A}_\lambda\to\sh{A}_\lambda$ el homomorfismo de $\sh{O}_{\mathrm{S}_0}$-Módulos que define la multiplicación en el $\sh{O}_{\mathrm{S}_0}$-Álgebra $\sh{A}_\lambda$; la hipótesis sobre los $\varphi_{\mu\lambda}$ implica que los $m_\lambda$ forman un sistema inductivo de homomorfismos,
\oldpage[IV-3]{8}
y, puesto que el funtor $\varinjlim$ conmuta con los productos tensoriales, $m=\varinjlim m_\lambda$ es un homomorfismo $\sh{A}\otimes\sh{A}\to\sh{A}$ de $\sh{O}_{\mathrm{S}_0}$-Módulos; pasando al límite en los diagramas conmutativos que expresan la asociatividad de $m_\lambda$ y la existencia de la sección identidad en $\sh{A}_\lambda$, se ve que $m$ dota a $\sh{A}$ de una estructura de $\sh{O}_{\mathrm{S}_0}$-Álgebra y que $\varphi_\lambda$ es un homomorfismo de $\sh{O}_{\mathrm{S}_0}$-Álgebras para todo $\lambda\in\mathrm{L}$.
Además, $\sh{A}$ es el límite inductivo del sistema $(\sh{A}_\lambda,\varphi_{\mu\lambda})$ \emph{en la categoría de $\sh{O}_{\mathrm{S}_0}$-Álgebras}; en otras palabras, para todo $\sh{O}_{\mathrm{S}_0}$-Álgebra $\sh{B}$, la aplicación canónica
\begin{equation}
  \operatorname{Hom}_{\sh{O}_{\mathrm{S}_0}\text{-}\mathrm{Alg.}}(\sh{A},\sh{B})\to\varprojlim\operatorname{Hom}_{\sh{O}_{\mathrm{S}_0}\text{-}\mathrm{Alg.}}(\sh{A}_\lambda,\sh{B})
  \tag{8.2.1.1}
\label{IV.8.2.1.1}
\end{equation}
que asocia a cada homomorfismo $f:\sh{A}\to\sh{B}$ de $\sh{O}_{\mathrm{S}_0}$-Álgebras la familia $(f\circ\varphi_\lambda)$ es \emph{biyectiva}.
En efecto, ya sabemos que es inyectiva e identifica $\operatorname{Hom}_{\sh{O}_{\mathrm{S}_0}\text{-}\mathrm{Alg.}}(\sh{A},\sh{B})$ con un subconjunto de $\varprojlim\operatorname{Hom}_{\sh{O}_{\mathrm{S}_0}\text{-}\mathrm{Mod.}}(\sh{A}_\lambda,\sh{B})$; todo se reduce a ver que, si $(f_\lambda)$ es un sistema inductivo de homomorfismos de $\sh{O}_{\mathrm{S}_0}$-Álgebras, $f_\lambda:\sh{A}_\lambda\to\sh{B}$, su límite inductivo $f:\sh{A}\to\sh{B}$, que por definición es un homomorfismo de $\sh{O}_{\mathrm{S}_0}$-\emph{Módulos}, es también un homomorfismo de $\sh{O}_{\mathrm{S}_0}$-\emph{Álgebras}; pero esto resulta de pasar al límite inductivo en el diagrama conmutativo de $\sh{O}_{\mathrm{S}_0}$-Módulos
\[
\xymatrix{
  \sh{A}_\lambda\otimes\sh{A}_\lambda \ar[r]^{f_\lambda\otimes f_\lambda} \ar[d]_{m_\lambda} & \sh{B}\otimes\sh{B} \ar[d] \\
  \sh{A}_\lambda \ar[r]_{f_\lambda} & \sh{B}
}
\]
y del hecho de que el funtor $\varinjlim$ conmuta con los productos tensoriales.
Observamos finalmente que, si los $\sh{A}_\lambda$ son $\sh{O}_{\mathrm{S}_0}$-Álgebras conmutativas, lo mismo vale para $\sh{A}$.
\end{env}

\begin{env}[8.2.2]
\label{IV.8.2.2}
Supongamos ahora que $\mathrm{S}_0$ es un \emph{preesquema} y que las $\sh{A}_\lambda$ son $\sh{O}_{\mathrm{S}_0}$-Álgebras (conmutativas) \emph{cuasi-coherentes}; sabemos entonces que $\sh{A}=\varinjlim\sh{A}_\lambda$ es un $\sh{O}_{\mathrm{S}_0}$-Álgebra cuasi-coherente \sref[I]{I.4.1.1}.
Denotemos por $\mathrm{S}_\lambda$ (resp. $\mathrm{S}$) el espectro del $\sh{O}_{\mathrm{S}_0}$-Álgebra $\sh{A}_\lambda$ (resp. $\sh{A}$) \sref[II]{II.1.3.1}, y sean $u_{\lambda\mu}:\mathrm{S}_\mu\to\mathrm{S}_\lambda$ (para $\lambda\leqslant\mu$) y $u_\lambda:\mathrm{S}\to\mathrm{S}_\lambda$ los $\mathrm{S}_0$-morfismos correspondientes a los homomorfismos $\varphi_{\mu\lambda}$ y $\varphi_\lambda$, respectivamente \sref[II]{II.1.2.7}; está claro que $(\mathrm{S}_\lambda,u_{\lambda\mu})$ es un \emph{sistema proyectivo} en la categoría de preesquemas $\mathrm{S}_0$.
Observamos que los $u_{\lambda\mu}$ y $u_\lambda$ son morfismos \emph{afines} \sref[II]{II.1.6.2}, por tanto \emph{cuasi-compactos} y \emph{separados}.
\end{env}

\begin{proposition}[8.2.3]
\label{IV.8.2.3}
Con las notaciones de \sref{IV.8.2.2}, los morfismos $u_\lambda:\mathrm{S}\to\mathrm{S}_\lambda$ hacen de $\mathrm{S}$ un límite proyectivo del sistema proyectivo $(\mathrm{S}_\lambda,u_{\lambda\mu})$ en la categoría de preesquemas.
Además, si $h:\mathrm{S}_0\to\mathrm{T}$ es un morfismo que convierte todo preesquema $\mathrm{S}_0$ en un preesquema $\mathrm{T}$, $\mathrm{S}$ es también un límite proyectivo del sistema $(\mathrm{S}_\lambda,u_{\lambda\mu})$ en la categoría de preesquemas $\mathrm{T}$.
\end{proposition}

\begin{proof}
Demostremos primero la segunda afirmación del enunciado en el caso $\mathrm{T}=\mathrm{S}_0$.
\oldpage[IV-3]{9}
Todo se reduce a demostrar que, si $\mathrm{X}$ es un preesquema $\mathrm{S}_0$ cualquiera, la aplicación canónica
\begin{equation}
  \operatorname{Hom}_{\mathrm{S}_0}(\mathrm{X},\mathrm{S})\to\varprojlim\operatorname{Hom}_{\mathrm{S}_0}(\mathrm{X},\mathrm{S}_\lambda)
  \tag{8.2.3.1}
\label{IV.8.2.3.1}
\end{equation}
que asocia a cada morfismo $\mathrm{S}_0$ $v:\mathrm{X}\to\mathrm{S}$ la familia $(u_\lambda\circ v)$ es \emph{biyectiva}.
Ahora bien, si $g:\mathrm{X}\to\mathrm{S}_0$ es el morfismo estructural y ponemos $\sh{B}=g_*(\sh{O}_\mathrm{X})$, que es un $\sh{O}_{\mathrm{S}_0}$-Álgebra, la aplicación \sref{IV.8.2.3.1} se identifica canónicamente con \sref{IV.8.2.1.1} \sref[II]{II.1.2.7}, y la conclusión se sigue de lo visto en \sref{IV.8.2.1}.

Las demás afirmaciones de \sref{IV.8.2.3} son consecuencias del siguiente lema general:
\end{proof}

\begin{lemma}[8.2.4]
\label{IV.8.2.4}
Sea $\mathbf{C}$ una categoría, $\mathrm{T}$ un objeto de $\mathbf{C}$ y $\mathbf{C}_\mathrm{T}$ la subcategoría de objetos de $\mathbf{C}$ sobre $\mathrm{T}$.
Sea $(\mathrm{S}_\lambda,u_{\lambda\mu})$ un sistema proyectivo en $\mathbf{C}_\mathrm{T}$; entonces todo límite proyectivo de este sistema en $\mathbf{C}_\mathrm{T}$ es también un límite proyectivo en $\mathbf{C}$, y recíprocamente.
\end{lemma}

\begin{proof}
Sea $f_\lambda:\mathrm{S}_\lambda\to\mathrm{T}$ el morfismo estructural.
Supongamos que $\mathrm{S}$ es un límite proyectivo de $(\mathrm{S}_\lambda,u_{\lambda\mu})$ en $\mathbf{C}$, y denotemos por $u_\lambda:\mathrm{S}\to\mathrm{S}_\lambda$ los morfismos canónicos correspondientes.
Consideremos un sistema proyectivo de morfismos $\mathrm{T}$ (en $\mathbf{C}$) $w_\lambda:\mathrm{Y}\to\mathrm{S}_\lambda$, donde $\mathrm{Y}\in\mathbf{C}_\mathrm{T}$.
Por hipótesis existe un único morfismo (en $\mathbf{C}$) $w:\mathrm{Y}\to\mathrm{S}$ tal que $w_\lambda=u_\lambda\circ w$ para todo $\lambda$.
La hipótesis de que los $u_\lambda$ son morfismos $\mathrm{T}$ implica que los morfismos $f_\lambda\circ u_\lambda:\mathrm{S}\to\mathrm{T}$ son todos iguales, y este compuesto $f$ dota a $\mathrm{S}$ de una estructura de objeto $\mathrm{T}$.
Si $g:\mathrm{Y}\to\mathrm{T}$ es el morfismo estructural de $\mathrm{Y}$, se tiene $f\circ w=f_\lambda\circ u_\lambda\circ w=f_\lambda\circ w_\lambda=g$ para todo $\lambda$, lo que demuestra que $w$ es un morfismo $\mathrm{T}$.
Recíprocamente, supongamos (con las mismas notaciones) que $\mathrm{S}$ es un límite proyectivo de $(\mathrm{S}_\lambda,u_{\lambda\mu})$ en $\mathbf{C}_\mathrm{T}$, y consideremos ahora un sistema proyectivo de morfismos (en $\mathbf{C}$) $w_\lambda:\mathrm{Y}\to\mathrm{S}_\lambda$.
Los compuestos $f_\lambda\circ w_\lambda:\mathrm{Y}\to\mathrm{T}$ son entonces todos iguales: en efecto, para dos índices cualesquiera $\lambda$, $\mu$, existe un índice $\nu$ tal que $\lambda\leqslant\nu$ y $\mu\leqslant\nu$, de donde $f_\nu=f_\lambda\circ u_{\lambda\nu}=f_\mu\circ u_{\mu\nu}$; como $w_\lambda=u_{\lambda\nu}\circ w_\nu$ y $w_\mu=u_{\mu\nu}\circ w_\nu$, se tiene $f_\lambda\circ w_\lambda=f_\lambda\circ u_{\lambda\nu}\circ w_\nu=f_\nu\circ w_\nu$ y $f_\mu\circ w_\mu=f_\mu\circ u_{\mu\nu}\circ w_\nu=f_\nu\circ w_\nu$.
Si $g:\mathrm{Y}\to\mathrm{T}$ es el único morfismo así definido, $g$ dota a $\mathrm{Y}$ de una estructura de objeto $\mathrm{T}$ y los $w_\lambda$ son morfismos $\mathrm{T}$; por tanto tienen un límite proyectivo $w:\mathrm{Y}\to\mathrm{S}$ que es un morfismo $\mathrm{T}$ y, \emph{a fortiori}, un morfismo de $\mathbf{C}$; además, la primera parte del argumento muestra que todo límite proyectivo $w'$ (en $\mathbf{C}$) del sistema proyectivo $(w_\lambda)$ es también un morfismo $\mathrm{T}$ y, por consiguiente, igual a $w$, lo que concluye la demostración del lema.
\end{proof}

\begin{proposition}[8.2.5]
\label{IV.8.2.5}
Bajo las condiciones de \sref{IV.8.2.2}, sea $\alpha$ un elemento de $\mathrm{L}$ y $\mathrm{X}_\alpha$ un preesquema $\mathrm{S}_\alpha$.
Para todo $\lambda\geqslant\alpha$, pongamos $\mathrm{X}_\lambda=\mathrm{X}_\alpha\times_{\mathrm{S}_\alpha}\mathrm{S}_\lambda$ y, para $\alpha\leqslant\lambda\leqslant\mu$, pongamos $v_{\lambda\mu}=1_{\mathrm{X}_\alpha}\times u_{\lambda\mu}$, de modo que $(\mathrm{X}_\lambda,v_{\lambda\mu})$ sea un sistema proyectivo de preesquemas $\mathrm{X}_\alpha$, cuyo conjunto de índices está formado por los $\lambda\geqslant\alpha$ de $\mathrm{L}$.
Pongamos asimismo $\mathrm{X}=\mathrm{X}_\alpha\times_{\mathrm{S}_\alpha}\mathrm{S}$ y $v_\lambda=1_{\mathrm{X}_\alpha}\times u_\lambda$.
Entonces los morfismos $\mathrm{X}_\alpha$ $v_\lambda:\mathrm{X}\to\mathrm{X}_\lambda$ hacen de $\mathrm{X}$ un límite proyectivo del sistema proyectivo $(\mathrm{X}_\lambda,v_{\lambda\mu})$ en la categoría de preesquemas $\mathrm{X}_\alpha$, o en la categoría de todos los preesquemas.
\end{proposition}

\begin{proof}
Esto también se deducirá del siguiente lema general:
\end{proof}

\begin{lemma}[8.2.6]
\label{IV.8.2.6}
Sea $\mathbf{C}$ una categoría en la que existen productos fibrados, $q:\mathrm{T}'\to\mathrm{T}$ un morfismo de $\mathbf{C}$, y $\mathbf{C}_\mathrm{T}$ (resp. $\mathbf{C}_{\mathrm{T}'}$) la categoría de objetos de $\mathbf{C}$ sobre
\oldpage[IV-3]{10}
$\mathrm{T}$ (resp. $\mathrm{T}'$).
Sea $(\mathrm{S}_\lambda,u_{\lambda\mu})$ un sistema proyectivo (no necesariamente filtrado) en $\mathbf{C}_\mathrm{T}$, y pongamos $\mathrm{S}'_\lambda=\mathrm{S}_\lambda\times_\mathrm{T}\mathrm{T}'$, $u'_{\lambda\mu}=u_{\lambda\mu}\times 1_{\mathrm{T}'}$, de modo que $(\mathrm{S}'_\lambda,u'_{\lambda\mu})$ sea un sistema proyectivo en $\mathbf{C}_{\mathrm{T}'}$.
Entonces, si $(\mathrm{S},u_\lambda)$ es un límite proyectivo de $(\mathrm{S}_\lambda,u_{\lambda\mu})$ en $\mathbf{C}_\mathrm{T}$, $(\mathrm{S}\times_\mathrm{T}\mathrm{T}',u_\lambda\times 1_{\mathrm{T}'})$ es un límite proyectivo de $(\mathrm{S}'_\lambda,u'_{\lambda\mu})$ en $\mathbf{C}_{\mathrm{T}'}$.
\end{lemma}

\begin{proof}
Por hipótesis, para todo $\lambda$ se tiene un diagrama conmutativo
\[
\xymatrix{
  \mathrm{S}' \ar[r]^{u'_\lambda} \ar[d]_p & \mathrm{S}'_\lambda \ar[r]^{h'_\lambda} \ar[d]_{p_\lambda} & \mathrm{T}' \ar[d]^q \\
  \mathrm{S} \ar[r]_{u_\lambda} & \mathrm{S}_\lambda \ar[r]_{h_\lambda} & \mathrm{T}
}
\]
donde hemos puesto $\mathrm{S}'=\mathrm{S}\times_\mathrm{T}\mathrm{T}'$, $u'_\lambda=u_\lambda\times 1_{\mathrm{T}'}$, $h'_\lambda=h_\lambda\times 1_{\mathrm{T}'}$.
Sea $\mathrm{Y}$ un objeto $\mathrm{T}'$, $g':\mathrm{Y}\to\mathrm{T}'$, y consideremos un sistema proyectivo de morfismos $\mathrm{T}'$ $w'_\lambda:\mathrm{Y}\to\mathrm{S}'_\lambda$.
Entonces $\mathrm{Y}$ es un objeto $\mathrm{T}$ para el morfismo $g=q\circ g'$, y los $w_\lambda=p_\lambda\circ w'_\lambda$ son morfismos $\mathrm{T}$, pues $h_\lambda\circ w_\lambda=h_\lambda\circ p_\lambda\circ w'_\lambda=q\circ h'_\lambda\circ w'_\lambda=q\circ g'$ por hipótesis.
Además, se verifica inmediatamente que $(w_\lambda)$ es un sistema proyectivo.
Por hipótesis existe entonces un único morfismo $\mathrm{T}$ $w:\mathrm{Y}\to\mathrm{S}$ tal que $u_\lambda\circ w=w_\lambda$ para todo $\lambda$.
Por la definición del producto fibrado, existe un único morfismo $\mathrm{T}'$ $w':\mathrm{Y}\to\mathrm{S}'$ tal que $p\circ w'=w$.
Entonces $u_\lambda\circ p\circ w'=u_\lambda\circ w=w_\lambda=p_\lambda\circ w'_\lambda$, lo que también puede escribirse $p_\lambda\circ u'_\lambda\circ w'=p_\lambda\circ w'_\lambda$; por otra parte, como $w'$ es un morfismo $\mathrm{T}'$, $h'_\lambda\circ u'_\lambda\circ w'=g'=h'_\lambda\circ w'_\lambda$. La definición de $\mathrm{S}'_\lambda$ como producto fibrado $\mathrm{S}_\lambda\times_\mathrm{T}\mathrm{T}'$ da por tanto $u'_\lambda\circ w'=w'_\lambda$, y es inmediato que $w'$ es el único morfismo $\mathrm{T}'$ que satisface estas relaciones, de donde se sigue el lema.
\end{proof}

\begin{remark}[8.2.7]
\label{IV.8.2.7}
Dado un espacio anillado $\mathrm{S}$, los límites inductivos respecto de un conjunto preordenado $\mathrm{L}$ (no necesariamente filtrado) existen en la categoría de $\sh{O}_\mathrm{S}$-Álgebras conmutativas, pues el límite inductivo filtrado existe por \sref{IV.8.2.1} y, por otra parte, para dos homomorfismos cualesquiera de $\sh{O}_\mathrm{S}$-Álgebras $\sh{A}\to\sh{B}$, $\sh{A}\to\sh{C}$, el producto tensorial $\sh{B}\otimes_{\sh{A}}\sh{C}$ es la «suma amalgamada» correspondiente en esta categoría.

Cuando $\mathrm{S}$ es un preesquema, sabemos que el producto tensorial $\sh{B}\otimes_{\sh{A}}\sh{C}$ es un $\sh{O}_\mathrm{S}$-Álgebra cuasi-coherente cuando esto vale para $\sh{A}$, $\sh{B}$ y $\sh{C}$ \sref[I]{I.1.3.13}; de ello se concluye que, en la categoría de $\sh{O}_\mathrm{S}$-Álgebras cuasi-coherentes, existen siempre límites inductivos para cualquier conjunto de índices preordenado.
Esto permite generalizar la definición de límite proyectivo de preesquemas y las proposiciones \sref{IV.8.2.3} y \sref{IV.8.2.5} al caso en que el conjunto preordenado $\mathrm{L}$ no sea necesariamente filtrado.
\end{remark}

\begin{env}[8.2.8]
\label{IV.8.2.8}
Con las notaciones de \sref{IV.8.2.2}, pongamos $u_{\lambda\mu}=(\psi_{\lambda\mu},\theta_{\lambda\mu})$ y $u_\lambda=(\psi_\lambda,\theta_\lambda)$; $(\mathrm{S}_\lambda,\psi_{\lambda\mu})$ es así un sistema proyectivo de espacios topológicos y $(\psi_\lambda)$ un sistema proyectivo de aplicaciones continuas $\mathrm{S}\to\mathrm{S}_\lambda$ entre los espacios subyacentes, respectivamente, de los preesquemas $\mathrm{S}$ y $\mathrm{S}_\lambda$.
\end{env}

\begin{proposition}[8.2.9]
\label{IV.8.2.9}
Con las notaciones de \sref{IV.8.2.8}, el límite proyectivo del sistema proyectivo $(\psi_\lambda)$ de aplicaciones continuas es un homeomorfismo del espacio subyacente de $\mathrm{S}$ sobre el límite proyectivo del sistema proyectivo $(\mathrm{S}_\lambda,\psi_{\lambda\mu})$ de espacios topológicos.
\end{proposition}

\begin{proof}
\oldpage[IV-3]{11}
Sea $\mathrm{T}$ el espacio límite topológico del sistema $(\mathrm{S}_\lambda,\psi_{\lambda\mu})$ y pongamos $\psi=\varprojlim\psi_\lambda:\mathrm{S}\to\mathrm{T}$.
Podemos restringirnos al caso en que $\mathrm{S}_0=\mathrm{S}_\alpha$ para algún $\alpha\in\mathrm{L}$ y $\lambda\geqslant\alpha$.

\medskip
\noindent
(i) Mostremos primero que la topología de $\mathrm{S}$ es la imagen inversa por $\psi$ de la topología de $\mathrm{T}$; en otras palabras, si $\pi_\lambda:\mathrm{T}\to\mathrm{S}_\lambda$ es la aplicación canónica, se trata de demostrar que todo abierto de $\mathrm{S}$ es unión de abiertos de la forma $\psi^{-1}(\pi_\lambda^{-1}(\mathrm{U}_\lambda))$, donde $\mathrm{U}_\lambda$ es abierto en $\mathrm{S}_\lambda$.
Ahora bien, por definición todo abierto de $\mathrm{S}$ es unión de abiertos obtenidos del modo siguiente: se considera un abierto afín $\mathrm{U}_0$ de $\mathrm{S}_0$, de anillo $\mathrm{A}_0$, de modo que $\psi_\alpha^{-1}(\mathrm{U}_0)$ sea el abierto afín de $\mathrm{S}$ de anillo $\mathrm{A}=\Gamma(\mathrm{U}_0,\sh{A})$; después se toma un elemento $f\in\mathrm{A}$ y se considera en $\operatorname{Spec}(\mathrm{A})$, identificado con $\psi_\alpha^{-1}(\mathrm{U}_0)$, el abierto $\mathrm{D}(f)$.
Estos son los abiertos $\mathrm{D}(f)$ que forman una base de la topología de $\mathrm{S}$ \sref[II]{II.1.3.1}.
Ahora, si ponemos $\mathrm{A}_\lambda=\Gamma(\mathrm{U}_0,\sh{A}_\lambda)$, se tiene $\mathrm{A}=\varinjlim\mathrm{A}_\lambda$ \sref[I]{I.1.3.9}, por lo que existe un índice $\lambda$ tal que $f$ es la imagen canónica de un elemento $f_\lambda\in\mathrm{A}_\lambda$; entonces $\mathrm{D}(f)=\psi_\lambda^{-1}(\mathrm{D}(f_\lambda))$ \sref[I]{I.1.2.2}, y como $\psi_\lambda=\pi_\lambda\circ\psi$, queda demostrada nuestra afirmación.

\medskip
\noindent
(ii) Mostremos ahora que $\psi$ es \emph{biyectiva}, lo que concluirá la demostración.
Como $\mathrm{S}$ es un espacio de Kolmogoroff, de (i) se deduce ya que $\psi$ es inyectiva, y queda por demostrar que $\psi$ es sobreyectiva.
Evidentemente podemos sustituir $\mathrm{T}$ por un abierto $\pi_\alpha^{-1}(\mathrm{U}_0)$, donde $\mathrm{U}_0$ es un abierto afín de $\mathrm{S}_\alpha=\mathrm{S}_0$; por tanto podemos restringirnos al caso en que los $\mathrm{S}_\lambda$ y $\mathrm{S}$ sean afines, es decir, $\sh{A}_\lambda$ sea el haz de álgebras asociado a un $\mathrm{A}_0$-álgebra $\mathrm{A}_\lambda$, y $\sh{A}$ el haz de álgebras asociado a $\mathrm{A}=\varinjlim\mathrm{A}_\lambda$; seguiremos denotando por $\varphi_{\mu\lambda}:\mathrm{A}_\lambda\to\mathrm{A}_\mu$ y $\varphi_\lambda:\mathrm{A}_\lambda\to\mathrm{A}$ los homomorfismos canónicos.
Por definición, un elemento de $\mathrm{T}$ es una familia $(\mathfrak{p}_\lambda)_{\lambda\in\mathrm{L}}$, donde $\mathfrak{p}_\lambda$ es un ideal primo de $\mathrm{A}_\lambda$ y donde $\mathfrak{p}_\lambda=\varphi_{\mu\lambda}^{-1}(\mathfrak{p}_\mu)$ para $\lambda\leqslant\mu$.
Sabemos (\sref{IV.5.13.3} y \sref{IV.5.13.1}) que existe un ideal primo $\mathfrak{p}$ de $\mathrm{A}$ tal que $\mathfrak{p}_\lambda=\varphi_\lambda^{-1}(\mathfrak{p})$ para todo $\lambda\in\mathrm{L}$, lo que concluye la demostración.
\end{proof}

En particular, hemos demostrado así el siguiente

\begin{corollary}[8.2.10]
\label{IV.8.2.10}
Sea $(\mathrm{A}_\lambda)_{\lambda\in\mathrm{L}}$ un sistema inductivo filtrado de anillos, y sean $\mathrm{A}=\varinjlim\mathrm{A}_\lambda$, $\varphi_\lambda:\mathrm{A}_\lambda\to\mathrm{A}$ los homomorfismos canónicos.
La aplicación canónica $\mathfrak{p}\mapsto(\varphi_\lambda^{-1}(\mathfrak{p}))$ es un homeomorfismo de $\operatorname{Spec}(\mathrm{A})$ sobre el espacio topológico $\varprojlim\operatorname{Spec}(\mathrm{A}_\lambda)$.
\end{corollary}

\begin{corollary}[8.2.11]
\label{IV.8.2.11}
Con las notaciones de \sref{IV.8.2.8}, para todo abierto cuasi-compacto $\mathrm{U}$ de $\mathrm{S}$ existen un índice $\lambda$ y un abierto cuasi-compacto $\mathrm{U}_\lambda$ de $\mathrm{S}_\lambda$ tales que $\mathrm{U}=\psi_\lambda^{-1}(\mathrm{U}_\lambda)$.
\end{corollary}

\begin{proof}
Esto se deduce del hecho de que, por definición de la topología del límite proyectivo, los $\psi_\lambda^{-1}(\mathrm{U}_\lambda)$ ($\mathrm{U}_\lambda$ abierto cuasi-compacto de $\mathrm{S}_\lambda$) forman una base de la topología de $\mathrm{S}$, y del hecho de que el conjunto de índices $\mathrm{L}$ es filtrado.
\end{proof}

\begin{corollary}[8.2.12]
\label{IV.8.2.12}
Con las notaciones de \sref{IV.8.2.8}, el límite inductivo del sistema inductivo de homomorfismos $\theta_\lambda^\sharp:\psi_\lambda^*(\sh{O}_{\mathrm{S}_\lambda})\to\sh{O}_\mathrm{S}$ de haces de anillos sobre $\mathrm{S}$ es un isomorfismo
\begin{equation}
  \varinjlim_\lambda\psi_\lambda^*(\sh{O}_{\mathrm{S}_\lambda})\xrightarrow{\sim}\sh{O}_\mathrm{S}.
  \label{IV.8.2.12.1}
\tag{8.2.12.1}
\end{equation}
\end{corollary}

\begin{proof}
Podemos suponer evidentemente que los $\mathrm{S}_\lambda$ son afines; con las notaciones de la demostración de \sref{IV.8.2.9}, todo se reduce a comprobar que el límite inductivo del sistema inductivo de aplicaciones canónicas $(\mathrm{A}_\lambda)_{\mathfrak{p}_\lambda}\to\mathrm{A}_\mathfrak{p}$ es un isomorfismo, lo cual no es otra cosa que \sref{IV.5.13.3},~(ii).
\end{proof}

\begin{proposition}[8.2.13]
\label{IV.8.2.13}
Supongamos que los morfismos $u_{\lambda\mu}:\mathrm{S}_\mu\to\mathrm{S}_\lambda$ son inmersiones abiertas, de modo que $\mathrm{S}_\mu$ se identifica con el subesquema inducido en un abierto de $\mathrm{S}_\lambda$ para $\lambda\leqslant\mu$.
Entonces, para todo $\alpha\in\mathrm{L}$, el espacio subyacente del preesquema $\mathrm{S}$ se identifica con el subespacio de $\mathrm{S}_\alpha$ dado por la intersección de los $\mathrm{S}_\lambda$ para $\lambda\geqslant\alpha$, y el haz estructural $\sh{O}_\mathrm{S}$ se identifica con el haz inducido \textnormal{(}$\mathrm{G}$, II, $1.5$\textnormal{)} por $\sh{O}_{\mathrm{S}_\alpha}$ sobre esta intersección; más generalmente, para todo $\sh{O}_{\mathrm{S}_\alpha}$-Módulo $\sh{F}_\alpha$, $u_\alpha^*(\sh{F}_\alpha)$ se identifica con el $\sh{O}_\mathrm{S}$-Módulo inducido por $\sh{F}_\alpha$ sobre $\mathrm{S}$.
\end{proposition}

\begin{proof}
\oldpage[IV-3]{12}
La primera afirmación se deduce de \sref{IV.8.2.9}, teniendo en cuenta la definición de límite proyectivo de espacios topológicos; además, todos los $\psi_\lambda^*(\sh{O}_{\mathrm{S}_\lambda})$ son iguales al haz inducido por $\sh{O}_{\mathrm{S}_\alpha}$ sobre $\mathrm{S}$ por definición \sref[$0_{\mathrm{I}}$]{0.3.7.1}, y, con las notaciones de la demostración de \sref{IV.8.2.9}, los homomorfismos $(\mathrm{A}_\lambda)_{\mathfrak{p}_\lambda}\to(\mathrm{A}_\mu)_{\mathfrak{p}_\mu}$ son biyectivos para un sistema $(\mathfrak{p}_\lambda)$ de ideales primos correspondiente a un único punto de $\mathrm{S}$; la afirmación relativa a $\sh{O}_\mathrm{S}$ se deduce entonces de \sref{IV.8.2.12}.
La última afirmación se deduce entonces de la definición de $u_\alpha^*(\sh{F}_\alpha)$ \sref[$0_{\mathrm{I}}$]{0.4.3.1}.
\end{proof}

\begin{remark}[8.2.14]
\label{IV.8.2.14}
Los resultados de \sref{IV.8.2.9} y \sref{IV.8.2.12} muestran que $\mathrm{S}$ es un límite proyectivo del sistema proyectivo $(\mathrm{S}_\lambda,u_{\lambda\mu})$ en la categoría de \emph{todos} los espacios anillados (o de todos los espacios localmente anillados).
En efecto, sea $\mathrm{Y}$ un espacio anillado y consideremos un sistema proyectivo de morfismos de espacios anillados $w_\lambda:\mathrm{Y}\to\mathrm{S}_\lambda$.
Si se escribe $w_\lambda=(\rho_\lambda,\omega_\lambda)$, los $\rho_\lambda$ forman un sistema proyectivo de aplicaciones continuas y, por \sref{IV.8.2.9}, su límite proyectivo se identifica con una aplicación continua $\rho:\mathrm{Y}\to\mathrm{S}$ tal que $\rho_\lambda=\psi_\lambda\circ\rho$. Por otra parte, los homomorfismos
\[
 \omega_\lambda^\sharp:\rho_\lambda^*(\sh{O}_{\mathrm{S}_\lambda})\longrightarrow\sh{O}_\mathrm{Y}
\]
forman un sistema inductivo de homomorfismos de haces de anillos. Puesto que
\[
 \rho_\lambda^*(\sh{O}_{\mathrm{S}_\lambda})
 =\rho^*\bigl(\psi_\lambda^*(\sh{O}_{\mathrm{S}_\lambda})\bigr)
\]
y el funtor $\rho^*$ es exacto, el límite inductivo de los $\rho_\lambda^*(\sh{O}_{\mathrm{S}_\lambda})$ es $\rho^*(\sh{O}_\mathrm{S})$ por \sref{IV.8.2.12}. Existe por tanto un único homomorfismo
\[
 \omega^\sharp:\rho^*(\sh{O}_\mathrm{S})\longrightarrow\sh{O}_\mathrm{Y}
\]
tal que $\omega_\lambda^\sharp=\omega^\sharp\circ\rho^*(\theta_\lambda^\sharp)$ para todo $\lambda$, lo que demuestra la afirmación.
\end{remark}

\subsection{Subconjuntos constructibles de un límite proyectivo de preesquemas}
\label{subsection:IV.8.3}

\begin{env}[8.3.1]
\label{IV.8.3.1}
En todo lo que sigue en este párrafo suponemos satisfechas las condiciones de \sref{IV.8.2.2} y conservamos las notaciones.
\end{env}

\begin{theorem}[8.3.2]
\label{IV.8.3.2}
Para todo $\lambda$, sean $\mathrm{E}_\lambda$, $\mathrm{F}_\lambda$ dos subconjuntos de $\mathrm{S}_\lambda$.
Pongamos
\begin{equation}
  \mathrm{E}=\bigcap_\lambda u_\lambda^{-1}(\mathrm{E}_\lambda),\qquad \mathrm{F}=\bigcup_\lambda u_\lambda^{-1}(\mathrm{F}_\lambda).
  \tag{8.3.2.1}
\label{IV.8.3.2.1}
\end{equation}
Supongamos satisfechas las condiciones siguientes:
\begin{enumerate}[label=(\roman*)]
   \item Para todo $\lambda$, $\mathrm{E}_\lambda$ es pro-constructible y $\mathrm{F}_\lambda$ es ind-constructible \sref{IV.1.9.4}.
   \item Para $\lambda\leqslant\mu$, se tiene
  \[
    \mathrm{E}_\mu\subset u_{\lambda\mu}^{-1}(\mathrm{E}_\lambda),\qquad \mathrm{F}_\mu\supset u_{\lambda\mu}^{-1}(\mathrm{F}_\lambda).
  \]
   \item Existe $\alpha\in\mathrm{L}$ tal que $\mathrm{S}_\alpha$ es cuasi-compacto (lo que implica que $\mathrm{S}_\lambda$ es cuasi-compacto para $\lambda\geqslant\alpha$).
\end{enumerate}
Entonces las propiedades siguientes son equivalentes:
\begin{enumerate}[label=\emph{\alph*)}]
   \item $\mathrm{E}\subset\mathrm{F}$.
   \item Existe $\lambda\geqslant\alpha$ tal que $u_\lambda^{-1}(\mathrm{E}_\lambda)\subset u_\lambda^{-1}(\mathrm{F}_\lambda)$ (y entonces también se tiene $u_\mu^{-1}(\mathrm{E}_\mu)\subset u_\mu^{-1}(\mathrm{F}_\mu)$ para
\oldpage[IV-3]{13}
$\mu\geqslant\lambda$).
   \item Existe $\lambda\geqslant\alpha$ tal que $\mathrm{E}_\lambda\subset\mathrm{F}_\lambda$ (y entonces también se tiene $\mathrm{E}_\mu\subset\mathrm{F}_\mu$ para $\mu\geqslant\lambda$).
\end{enumerate}
Las observaciones entre paréntesis de \emph{b)} y \emph{c)} se deducen de (ii).
Pongamos
\[
  \mathrm{G}_\lambda=\mathrm{E}_\lambda\cap(\mathrm{S}_\lambda-\mathrm{F}_\lambda),\qquad \mathrm{G}=\mathrm{E}\cap(\mathrm{S}-\mathrm{F}).
\]
Entonces $\mathrm{G}_\lambda$ es un subconjunto pro-constructible de $\mathrm{S}_\lambda$ \sref{IV.1.9.5},~(ii), y, en virtud de \sref{IV.8.3.2.1} y de la condición (ii), se tiene
\[
  \mathrm{G}_\mu\subset u_{\lambda\mu}^{-1}(\mathrm{G}_\lambda)\qquad\text{para }\lambda\leqslant\mu
\]
\[
  \mathrm{G}=\bigcap_\lambda u_\lambda^{-1}(\mathrm{G}_\lambda).
\]
Quedamos así reducidos a demostrar el caso particular de \sref{IV.8.3.2} correspondiente a $\mathrm{F}_\lambda=\varnothing$ para todo $\lambda$:
\end{theorem}

\begin{corollary}[8.3.3]
\label{IV.8.3.3}
Para todo $\lambda$, sea $\mathrm{E}_\lambda$ un subconjunto pro-constructible de $\mathrm{S}_\lambda$ tal que, para $\lambda\leqslant\mu$, $\mathrm{E}_\mu\subset u_{\lambda\mu}^{-1}(\mathrm{E}_\lambda)$.
Supongamos que existe $\alpha\in\mathrm{L}$ tal que $\mathrm{S}_\alpha$ es cuasi-compacto.
Entonces las condiciones siguientes son equivalentes:
\begin{enumerate}[label=\emph{\alph*)}]
   \item $\mathrm{E}=\bigcap_\lambda u_\lambda^{-1}(\mathrm{E}_\lambda)=\varnothing$.
   \item Existe $\lambda$ tal que $u_\lambda^{-1}(\mathrm{E}_\lambda)=\varnothing$ (y entonces $u_\mu^{-1}(\mathrm{E}_\mu)=\varnothing$ para $\mu\geqslant\lambda$).
   \item Existe $\lambda$ tal que $\mathrm{E}_\lambda=\varnothing$ (y entonces $\mathrm{E}_\mu=\varnothing$ para $\mu\geqslant\lambda$).
\end{enumerate}
\end{corollary}

\begin{proof}
Es claro que \emph{c)} implica \emph{a)}.
Demostremos que \emph{a)} implica \emph{b)}: como $\mathrm{S}_\lambda$ es cuasi-compacto, lo mismo vale para $\mathrm{S}$ \sref{IV.8.2.2}; como $\mathrm{E}_\lambda$ es pro-constructible, lo mismo vale para $u_\lambda^{-1}(\mathrm{E}_\lambda)$ \sref{IV.1.9.5},~(vi); por tanto, la familia filtrada decreciente de subconjuntos pro-constructibles $u_\lambda^{-1}(\mathrm{E}_\lambda)$ tiene intersección vacía, y por \sref{IV.1.9.9} uno de ellos es vacío.

Finalmente, demostremos que \emph{b)} implica \emph{c)}.
Como $\mathrm{S}_\alpha$ es cuasi-compacto y $\mathrm{L}$ es filtrado, podemos sustituir $\mathrm{S}_\alpha$ por un afín abierto; supongamos entonces \sref{IV.8.2.2} que $\mathrm{S}$ y los $\mathrm{S}_\lambda$ para $\lambda\geqslant\alpha$ son afines; se tiene entonces \sref{IV.1.9.2.1}, para $\lambda\geqslant\alpha$,
\[
  u_\lambda(\mathrm{S})=\bigcap_{\mu\geqslant\lambda}u_{\lambda\mu}(\mathrm{S}_\mu)
\]
de donde
\[
  \mathrm{E}_\lambda\cap u_\lambda(\mathrm{S})=\bigcap_{\mu\geqslant\lambda}(\mathrm{E}_\lambda\cap u_{\lambda\mu}(\mathrm{S}_\mu)).
\]
Ahora bien, como $u_\lambda$ y los $u_{\lambda\mu}$ son cuasi-compactos, $u_\lambda(\mathrm{S})$ y $u_{\lambda\mu}(\mathrm{S}_\mu)$ son pro-constructibles en $\mathrm{S}_\lambda$ \sref{IV.1.9.5},~(viii); por tanto, los conjuntos $\mathrm{E}_\lambda\cap u_{\lambda\mu}(\mathrm{S}_\mu)$ para $\mu\geqslant\lambda$ forman una familia filtrada decreciente de subconjuntos pro-constructibles de $\mathrm{S}_\lambda$.
Como $\mathrm{S}_\lambda$ es cuasi-compacto, la hipótesis \emph{b)} implica que la intersección de esta familia es vacía y, por \sref{IV.1.9.9}, uno de los conjuntos $\mathrm{E}_\lambda\cap u_{\lambda\mu}(\mathrm{S}_\mu)$ es vacío; por tanto, $\mathrm{E}_\mu\subset u_{\lambda\mu}^{-1}(\mathrm{E}_\lambda)$ es vacío.
\end{proof}

\begin{corollary}[8.3.4]
\label{IV.8.3.4}
Para todo $\lambda$, sea $\mathrm{F}_\lambda$ un subconjunto ind-constructible de $\mathrm{S}_\lambda$ tal que, para $\lambda\leqslant\mu$, se tenga $u_{\lambda\mu}^{-1}(\mathrm{F}_\lambda)\subset\mathrm{F}_\mu$.
Supongamos que existe $\alpha\in\mathrm{L}$ tal que $\mathrm{S}_\alpha$ es cuasi-compacto.
Entonces las condiciones siguientes son equivalentes:
\begin{enumerate}[label=\emph{\alph*)}]
   \item El conjunto $\mathrm{F}=\bigcup_\lambda u_\lambda^{-1}(\mathrm{F}_\lambda)$ es igual a $\mathrm{S}$.
   \item Existe $\lambda$ tal que $u_\lambda^{-1}(\mathrm{F}_\lambda)=\mathrm{S}$ (y entonces $u_\mu^{-1}(\mathrm{F}_\mu)=\mathrm{S}$ para $\mu\geqslant\lambda$).
\oldpage[IV-3]{14}
   \item Existe $\lambda$ tal que $\mathrm{F}_\lambda=\mathrm{S}_\lambda$ (y entonces $\mathrm{F}_\mu=\mathrm{S}_\mu$ para $\mu\geqslant\lambda$).
\end{enumerate}
Esto se deduce inmediatamente de \sref{IV.8.3.3} pasando a los complementarios.
\end{corollary}

\begin{corollary}[8.3.5]
\label{IV.8.3.5}
Para todo $\lambda$, sean $\mathrm{E}_\lambda$, $\mathrm{F}_\lambda$ dos subconjuntos constructibles de $\mathrm{S}_\lambda$ tales que, para $\mu\geqslant\lambda$, se tenga $\mathrm{E}_\mu\subset u_{\lambda\mu}^{-1}(\mathrm{E}_\lambda)$ y $\mathrm{F}_\mu\supset u_{\lambda\mu}^{-1}(\mathrm{F}_\lambda)$.
Supongamos que existe $\alpha$ tal que $\mathrm{S}_\alpha$ es cuasi-compacto.
Entonces, para que $\mathrm{E}\subset\mathrm{F}$ \textnormal{(}resp. $\mathrm{E}=\mathrm{F}$\textnormal{)}, es necesario y suficiente que exista $\lambda$ tal que $\mathrm{E}_\lambda\subset\mathrm{F}_\lambda$ \textnormal{(}resp. $\mathrm{E}_\lambda=\mathrm{F}_\lambda$\textnormal{)}, en cuyo caso también se tiene $\mathrm{E}_\mu\subset\mathrm{F}_\mu$ \textnormal{(}resp. $\mathrm{E}_\mu=\mathrm{F}_\mu$\textnormal{)} para $\mu\geqslant\lambda$.

Este es un caso particular de \sref{IV.8.3.2}.
\end{corollary}

\begin{corollary}[8.3.6]
\label{IV.8.3.6}
Supongamos que existe $\alpha$ tal que $\mathrm{S}_\alpha$ es cuasi-compacto.
Para que $\mathrm{S}=\varnothing$, es necesario y suficiente que exista $\lambda$ tal que $\mathrm{S}_\lambda=\varnothing$.
\end{corollary}

\begin{corollary}[8.3.7]
\label{IV.8.3.7}
Para todo $\lambda$, se tiene
\begin{equation}
  u_\lambda(\mathrm{S})=\bigcap_{\mu\geqslant\lambda}u_{\lambda\mu}(\mathrm{S}_\mu).
  \tag{8.3.7.1}
\label{IV.8.3.7.1}
\end{equation}
\end{corollary}

\begin{proof}
Es claro que el primer miembro de \sref{IV.8.3.7.1} está contenido en el segundo.
Sea $s$ un punto de $\mathrm{S}_\lambda$ y pongamos $\mathrm{X}_\lambda=\operatorname{Spec}(\boldsymbol{k}(s))$; consideremos el sistema proyectivo $(\mathrm{X}_\mu,v_{\mu\nu})$ donde $\mathrm{X}_\mu=\mathrm{X}_\lambda\times_{\mathrm{S}_\lambda}\mathrm{S}_\mu$ y $v_{\mu\nu}=1\times u_{\mu\nu}$ para $\lambda\leqslant\mu\leqslant\nu$; su límite proyectivo es $\mathrm{X}=\mathrm{X}_\lambda\times_{\mathrm{S}_\lambda}\mathrm{S}$, y $v_\lambda=1\times u_\lambda$ es la aplicación canónica $\mathrm{X}\to\mathrm{X}_\lambda$ \sref{IV.8.2.5}.
Si $s\in\bigcap_{\mu\geqslant\lambda}u_{\lambda\mu}(\mathrm{S}_\mu)$, esto implica que $\mathrm{X}_\mu\neq\varnothing$ para todo $\mu\geqslant\lambda$ \sref[I]{I.3.4.8}; se sigue entonces de \sref{IV.8.3.6} que $\mathrm{X}\neq\varnothing$, y por tanto $s\in u_\lambda(\mathrm{S})$ por \sref[I]{I.3.4.8}.
\end{proof}

\begin{proposition}[8.3.8]
\label{IV.8.3.8}
\begin{enumerate}[label=(\roman*)]
   \item Para que el morfismo $u_\lambda:\mathrm{S}\to\mathrm{S}_\lambda$ sea dominante (resp. sobreyectivo), es necesario y suficiente que, para $\mu\geqslant\lambda$, el morfismo $u_{\lambda\mu}:\mathrm{S}_\mu\to\mathrm{S}_\lambda$ sea dominante (resp. sobreyectivo).
   \item Si, para algún índice $\lambda$, los morfismos $u_{\lambda\mu}:\mathrm{S}_\mu\to\mathrm{S}_\lambda$ son planos (resp. fielmente planos) para todo $\mu\geqslant\lambda$, entonces el morfismo $u_\lambda:\mathrm{S}\to\mathrm{S}_\lambda$ es plano (resp. fielmente plano).
   \item Supongamos que los morfismos $u_\mu:\mathrm{S}\to\mathrm{S}_\mu$ son sobreyectivos para $\mu\geqslant\lambda$.
Para que $u_\lambda$ sea un morfismo abierto (resp. universalmente abierto), es necesario y suficiente que, para todo $\mu\geqslant\lambda$, $u_{\lambda\mu}$ sea un morfismo abierto (resp. universalmente abierto).
\end{enumerate}
\end{proposition}

\begin{proof}
\begin{enumerate}[label=(\roman*)]
   \item Como $u_\lambda(\mathrm{S})\subset u_{\lambda\mu}(\mathrm{S}_\mu)$ para todo $\mu\geqslant\lambda$, la necesidad de las condiciones es trivial, y se deduce inmediatamente de \sref{IV.8.3.7.1} que, si los $u_{\lambda\mu}$ son sobreyectivos, también lo es $u_\lambda$.
Supongamos ahora que los $u_{\lambda\mu}$ son dominantes para $\mu\geqslant\lambda$, y consideremos en $\mathrm{S}_\lambda$ un subconjunto abierto, cuasi-compacto y no vacío $\mathrm{U}_\lambda$; entonces los $\mathrm{U}_\mu=u_{\lambda\mu}^{-1}(\mathrm{U}_\lambda)$ para $\mu\geqslant\lambda$ forman un sistema proyectivo cuyo límite proyectivo es $\mathrm{U}=u_\lambda^{-1}(\mathrm{U}_\lambda)$ \sref{IV.8.2.5}.
Por hipótesis los $\mathrm{U}_\mu$ son todos no vacíos, y lo mismo vale para $\mathrm{U}$ por \sref{IV.8.3.6}, lo que demuestra que $u_\lambda$ es dominante.

   \item En virtud de (i), basta considerar el caso en que los $u_{\lambda\mu}$ son planos; entonces se puede restringir al caso en que $\mathrm{S}_\lambda$ es afín, y por tanto también lo son los $\mathrm{S}_\mu$ para $\mu\geqslant\lambda$ y $\mathrm{S}$, y la afirmación se deduce de \sref{IV.2.1.2} y \sref[0]{0.6.2.3}.

   \item En virtud de \sref{IV.8.2.5} y de \sref[I]{I.3.5.2}, basta tratar el caso de morfismos abiertos.
Como $u_\lambda=u_{\lambda\mu}\circ u_\mu$ y $u_\mu$ es sobreyectivo, sabemos que, si $u_\lambda$ es abierto, también lo es $u_{\lambda\mu}$ para $\mu\geqslant\lambda$ (Bourbaki, \emph{Top.\ g\'en.}, chap.~I\textsuperscript{er}, 3\textsuperscript{e}~ed., \textsection\,5, n\textsuperscript{o}\,1, prop.~1).
\oldpage[IV-3]{15}
Recíprocamente, para demostrar que $u_\lambda$ es abierto cuando todos los $u_{\lambda\mu}$ son abiertos para $\mu\geqslant\lambda$, basta comprobar que, para todo abierto cuasi-compacto $\mathrm{V}$ de $\mathrm{S}$, $u_\lambda(\mathrm{V})$ es abierto en $\mathrm{S}_\lambda$; entonces existe $\mu\geqslant\lambda$ tal que $\mathrm{V}=u_\mu^{-1}(\mathrm{V}_\mu)$, donde $\mathrm{V}_\mu$ es abierto en $\mathrm{S}_\mu$ \sref{IV.8.2.11}; como $u_\mu$ es sobreyectivo, se tiene $\mathrm{V}_\mu=u_\mu(\mathrm{V})$ y $u_\lambda(\mathrm{V})=u_{\lambda\mu}(u_\mu(\mathrm{V}))$ es, por tanto, abierto por hipótesis.
\end{enumerate}

Observamos que puede ocurrir que todos los $u_{\lambda\mu}$ sean abiertos sin que $u_\lambda$ sea abierto, cuando los $u_\mu$ no son sobreyectivos.
En efecto, se obtiene un ejemplo considerando un anillo íntegro $\mathrm{A}$ que no sea un cuerpo, y su cuerpo de fracciones $\mathrm{K}$, que es el límite inductivo de los $\mathrm{A}_f$, donde $f$ recorre $\mathrm{A}-\{0\}$; tomando $\mathrm{S}_1=\operatorname{Spec}(\mathrm{A})$, $\mathrm{S}_f=\operatorname{Spec}(\mathrm{A}_f)$ y $\mathrm{S}=\varprojlim\mathrm{S}_f=\operatorname{Spec}(\mathrm{K})$, el morfismo $\mathrm{S}\to\mathrm{S}_1$ no es abierto, aunque los morfismos $\mathrm{S}_f\to\mathrm{S}_1$ son abiertos.
\end{proof}

\begin{env}[8.3.9]
\label{IV.8.3.9}
Para todo preesquema $\mathrm{X}$, denotaremos como de costumbre por $\mathfrak{P}(\mathrm{X})$ el conjunto de todos los subconjuntos del conjunto subyacente de $\mathrm{X}$, y por $\mathfrak{C}(\mathrm{X})$ (resp. $\mathfrak{D}_c(\mathrm{X})$, $\mathfrak{F}_c(\mathrm{X})$, $\mathfrak{LF}_c(\mathrm{X})$) el conjunto de subconjuntos constructibles (resp. constructibles y abiertos, constructibles y cerrados, y constructibles y localmente cerrados) de $\mathrm{X}$.
Es claro que $(\mathfrak{P}(\mathrm{S}_\lambda),u_{\lambda\mu}^{-1})$ es un \emph{sistema inductivo de conjuntos} y que las aplicaciones $u_\lambda^{-1}:\mathfrak{P}(\mathrm{S}_\lambda)\to\mathfrak{P}(\mathrm{S})$ forman un sistema inductivo de aplicaciones; por tanto, pasando al límite inductivo, se obtiene una aplicación canónica
\begin{equation}
  u_\mathfrak{P}:\varinjlim\mathfrak{P}(\mathrm{S}_\lambda)\to\mathfrak{P}(\mathrm{S}).
  \tag{8.3.9.1}
\label{IV.8.3.9.1}
\end{equation}

Además, se deduce de \sref{IV.1.8.2} que $u_{\lambda\mu}^{-1}$ envía $\mathfrak{C}(\mathrm{S}_\lambda)$ (resp. $\mathfrak{D}_c(\mathrm{S}_\lambda)$, $\mathfrak{F}_c(\mathrm{S}_\lambda)$, $\mathfrak{LF}_c(\mathrm{S}_\lambda)$) en $\mathfrak{C}(\mathrm{S}_\mu)$ (resp. $\mathfrak{D}_c(\mathrm{S}_\mu)$, $\mathfrak{F}_c(\mathrm{S}_\mu)$, $\mathfrak{LF}_c(\mathrm{S}_\mu)$), y que $u_\lambda^{-1}$ envía $\mathfrak{C}(\mathrm{S}_\lambda)$ (resp. $\mathfrak{D}_c(\mathrm{S}_\lambda)$, $\mathfrak{F}_c(\mathrm{S}_\lambda)$, $\mathfrak{LF}_c(\mathrm{S}_\lambda)$) en $\mathfrak{C}(\mathrm{S})$ (resp. $\mathfrak{D}_c(\mathrm{S})$, $\mathfrak{F}_c(\mathrm{S})$, $\mathfrak{LF}_c(\mathrm{S})$).
Tenemos, pues, por restricción de \sref{IV.8.3.9.1}, las aplicaciones canónicas
\begin{align}
  &\varinjlim\mathfrak{C}(\mathrm{S}_\lambda)\to\mathfrak{C}(\mathrm{S}) \tag{8.3.9.2}\label{IV.8.3.9.2}\\
  &\varinjlim\mathfrak{D}_c(\mathrm{S}_\lambda)\to\mathfrak{D}_c(\mathrm{S}) \tag{8.3.9.3}\label{IV.8.3.9.3}\\
  &\varinjlim\mathfrak{F}_c(\mathrm{S}_\lambda)\to\mathfrak{F}_c(\mathrm{S}) \tag{8.3.9.4}\label{IV.8.3.9.4}\\
  &\varinjlim\mathfrak{LF}_c(\mathrm{S}_\lambda)\to\mathfrak{LF}_c(\mathrm{S}). \tag{8.3.9.5}
\label{IV.8.3.9.5}
\end{align}
\end{env}

\begin{env}[8.3.10]
\label{IV.8.3.10}
Sea $g_\alpha:\mathrm{X}_\alpha\to\mathrm{S}_\alpha$ un morfismo; con las notaciones de \sref{IV.8.2.5} se tiene, por una parte, una aplicación canónica $v_\mathfrak{P}:\varinjlim\mathfrak{P}(\mathrm{X}_\lambda)\to\mathfrak{P}(\mathrm{X})$; por otra parte, se tienen morfismos de proyección $g_\lambda:\mathrm{X}_\lambda\to\mathrm{S}_\lambda$ para todo $\lambda\geqslant\alpha$ y un morfismo de proyección $g:\mathrm{X}\to\mathrm{S}$.
Es claro que los $g_\lambda^{-1}:\mathfrak{P}(\mathrm{S}_\lambda)\to\mathfrak{P}(\mathrm{X}_\lambda)$ forman un sistema inductivo de aplicaciones, y que los diagramas
\[
\xymatrix{
  \mathfrak{P}(\mathrm{S}_\lambda) \ar[r]^{u_\lambda^{-1}} \ar[d]_{g_\lambda^{-1}} & \mathfrak{P}(\mathrm{S}) \ar[d]^{g^{-1}} \\
  \mathfrak{P}(\mathrm{X}_\lambda) \ar[r]_{v_\lambda^{-1}} & \mathfrak{P}(\mathrm{X})
}
\]
\oldpage[IV-3]{16}
son conmutativos; pasando al límite inductivo se deduce de ello un diagrama conmutativo
\begin{equation}
\xymatrix{
  \varinjlim\mathfrak{P}(\mathrm{S}_\lambda) \ar[r]^{u_\mathfrak{P}} \ar[d]_{\varinjlim g_\lambda^{-1}} & \mathfrak{P}(\mathrm{S}) \ar[d]^{g^{-1}} \\
  \varinjlim\mathfrak{P}(\mathrm{X}_\lambda) \ar[r]_{v_\mathfrak{P}} & \mathfrak{P}(\mathrm{X})
}
\tag{8.3.10.1}
\label{IV.8.3.10.1}
\end{equation}
y se deduce de \sref{IV.1.8.2} que se tienen diagramas análogos sustituyendo $\mathfrak{P}$ por $\mathfrak{C}$, $\mathfrak{D}_c$, $\mathfrak{F}_c$ o $\mathfrak{LF}_c$.
\end{env}

Se sigue de \sref{IV.8.3.5} que, bajo la hipótesis de que para algún $\alpha\in\mathrm{L}$, $\mathrm{S}_\alpha$ es \emph{cuasi-compacto}, la aplicación canónica \sref{IV.8.3.9.2} es \emph{inyectiva}.
Furthermore:

\begin{theorem}[8.3.11]
\label{IV.8.3.11}
Supongamos que existe $\alpha\in\mathrm{L}$ tal que $\mathrm{S}_\alpha$ es cuasi-compacto y cuasi-separado.
Entonces las aplicaciones canónicas \sref{IV.8.3.9.2}, \sref{IV.8.3.9.3}, \sref{IV.8.3.9.4} y \sref{IV.8.3.9.5} son \emph{biyectivas}.
\end{theorem}

\begin{proof}
En virtud de la observación precedente, queda por demostrar que estas aplicaciones son \emph{sobreyectivas}; como todo subconjunto constructible de $\mathrm{S}$ es unión finita de conjuntos de la forma $\mathrm{U}\cap\complement\mathrm{V}$, donde $\mathrm{U}$ y $\mathrm{V}$ son abiertos y constructibles, bastará demostrar que \sref{IV.8.3.9.3} es sobreyectiva para que lo sea también \sref{IV.8.3.9.2} (y también \sref{IV.8.3.9.4}, pasando a los complementarios).
Ahora bien, como los morfismos $\mathrm{S}_\lambda\to\mathrm{S}_\alpha$ y $\mathrm{S}\to\mathrm{S}_\alpha$ son afines y, por tanto, separados, los $\mathrm{S}_\lambda$ para $\lambda\geqslant\alpha$ y $\mathrm{S}$ son cuasi-compactos y cuasi-separados \sref{IV.1.2.2}, y se sabe que los subconjuntos abiertos constructibles de tal preesquema no son otros que los abiertos cuasi-compactos \sref{IV.1.8.1}.
La conclusión se sigue entonces de \sref{IV.8.2.11}, salvo para la aplicación \sref{IV.8.3.9.5}.
Para demostrar que esta última aplicación es sobreyectiva, consideremos un subconjunto localmente cerrado y constructible $\mathrm{Z}$ de $\mathrm{S}$; por tanto, $\mathrm{Z}$ es cuasi-compacto ($\mathbf{0}_{\mathrm{III}}$,~9.1.4).
Como, por hipótesis, todo punto $x\in\mathrm{Z}$ admite una vecindad abierta cuasi-compacta $\mathrm{V}_x$ en $\mathrm{S}$ tal que $\mathrm{Z}\cap\mathrm{V}_x$ es cerrado en $\mathrm{V}_x$, podemos cubrir $\mathrm{Z}$ con un número finito de los $\mathrm{V}_x$; en otras palabras, existe un abierto cuasi-compacto $\mathrm{U}$ que contiene a $\mathrm{Z}$ tal que $\mathrm{Z}$ es cerrado en $\mathrm{U}$; como $\mathrm{Z}$ es constructible en $\mathrm{S}$, también es constructible en $\mathrm{U}$ ($\mathbf{0}_{\mathrm{III}}$,~9.1.8).
Sabemos por \sref{IV.8.2.11} que existen $\lambda$ y un abierto cuasi-compacto $\mathrm{U}_\lambda$ de $\mathrm{S}_\lambda$ tales que $\mathrm{U}=u_\lambda^{-1}(\mathrm{U}_\lambda)$.
Aplicando a $\mathrm{U}$ (que es el límite proyectivo de los $\mathrm{U}_\mu=u_{\lambda\mu}^{-1}(\mathrm{U}_\lambda)$ para $\mu\geqslant\lambda$) el hecho de que la aplicación \sref{IV.8.3.9.4} es sobreyectiva, se ve que existen $\mu\geqslant\lambda$ y un subconjunto cerrado constructible $\mathrm{Z}_\mu$ en $\mathrm{U}_\mu$ tales que $\mathrm{Z}=u_\mu^{-1}(\mathrm{Z}_\mu)$.
Pero, como la inmersión canónica $\mathrm{U}_\mu\to\mathrm{S}_\mu$ es cuasi-compacta por hipótesis \sref{IV.1.2.7}, es de presentación finita \sref{IV.1.6.2},~(i), y $\mathrm{Z}_\mu$ es también un subconjunto constructible \emph{de} $\mathrm{S}_\mu$ en virtud de \sref{IV.1.8.4} y \sref{IV.1.8.1}; como $\mathrm{Z}_\mu$ es evidentemente localmente cerrado en $\mathrm{S}_\mu$, queda demostrada la afirmación.
\end{proof}

\begin{corollary}[8.3.12]
\label{IV.8.3.12}
Supongamos que existe $\alpha$ tal que $\mathrm{S}_\alpha$ es cuasi-compacto, y sea, para todo $\lambda$, $\mathrm{Z}_\lambda$ un subconjunto constructible de $\mathrm{S}_\lambda$ tal que $\mathrm{Z}_\mu=u_{\lambda\mu}^{-1}(\mathrm{Z}_\lambda)$ para $\mu\geqslant\lambda$.
Si $\mathrm{Z}=u_\lambda^{-1}(\mathrm{Z}_\lambda)$,
\oldpage[IV-3]{17}
entonces, para que $\mathrm{Z}$ sea abierto (resp. cerrado, resp. localmente cerrado) en $\mathrm{S}$, es necesario y suficiente que exista $\lambda\geqslant\alpha$ tal que $\mathrm{Z}_\lambda$ lo sea en $\mathrm{S}_\lambda$.
\end{corollary}

\begin{proof}
Cubramos $\mathrm{S}_\alpha$ con un número finito de abiertos afines $\mathrm{U}_\alpha^{(i)}$; entonces los $\mathrm{U}^{(i)}=u_\alpha^{-1}(\mathrm{U}_\alpha^{(i)})$ forman un recubrimiento afín abierto de $\mathrm{S}$, y para que $\mathrm{Z}$ sea abierto (resp. cerrado, resp. localmente cerrado) en $\mathrm{S}$ es necesario y suficiente que cada $\mathrm{Z}\cap\mathrm{U}^{(i)}$ lo sea en $\mathrm{U}^{(i)}$.
Como $\mathrm{L}$ es filtrado y cada $\mathrm{Z}\cap\mathrm{U}^{(i)}$ es constructible en $\mathrm{U}^{(i)}$ ($\mathbf{0}_{\mathrm{III}}$,~9.1.8), podemos restringirnos a demostrar el corolario cuando $\mathrm{S}_\alpha$ es afín, y por tanto cuasi-compacto y cuasi-separado; pero entonces se sigue de \sref{IV.8.3.11}.
\end{proof}

\begin{proposition}[8.3.13]
\label{IV.8.3.13}
Supongamos que los morfismos $u_{\lambda\mu}:\mathrm{S}_\mu\to\mathrm{S}_\lambda$ son planos para $\lambda\leqslant\mu$, y que existe $\alpha$ tal que $\mathrm{S}_\alpha$ es cuasi-compacto.
Para todo $\lambda$, sean $\mathrm{Z}'_\lambda$, $\mathrm{Z}''_\lambda$ dos subconjuntos pro-constructibles de $\mathrm{S}_\lambda$ tales que $\mathrm{Z}'_\mu=u_{\lambda\mu}^{-1}(\mathrm{Z}'_\lambda)$ y $\mathrm{Z}''_\mu=u_{\lambda\mu}^{-1}(\mathrm{Z}''_\lambda)$ para $\mu\geqslant\lambda$; supongamos además que $\overline{\mathrm{Z}'_\alpha}$ es constructible en $\mathrm{S}_\alpha$.
Pongamos $\mathrm{Z}'=u_\lambda^{-1}(\mathrm{Z}'_\lambda)$, $\mathrm{Z}''=u_\lambda^{-1}(\mathrm{Z}''_\lambda)$; para que $\mathrm{Z}''\subset\overline{\mathrm{Z}'}$, es necesario y suficiente que exista $\lambda\geqslant\alpha$ tal que $\mathrm{Z}''_\lambda\subset\overline{\mathrm{Z}'_\lambda}$.
\end{proposition}

\begin{proof}
En efecto, se sabe que $u_\lambda:\mathrm{S}\to\mathrm{S}_\lambda$ es también un morfismo plano para todo $\lambda$ \sref{IV.8.3.8}; como $\mathrm{Z}'_\lambda$ es pro-constructible, la clausura de $\mathrm{Z}'_\mu$ para $\mu\geqslant\lambda$ (resp. de $\mathrm{Z}'$) en $\mathrm{S}_\mu$ (resp. $\mathrm{S}$) es igual a $u_{\lambda\mu}^{-1}(\overline{\mathrm{Z}'_\lambda})$ (resp. $u_\lambda^{-1}(\overline{\mathrm{Z}'_\lambda})$) \sref{IV.2.3.10}.
Como $u_{\lambda\mu}^{-1}(\overline{\mathrm{Z}'_\lambda})$ y $u_\lambda^{-1}(\overline{\mathrm{Z}'_\lambda})$ son constructibles \sref{IV.1.8.2}, la conclusión se sigue de \sref{IV.8.3.2}.
\end{proof}

\subsection{Criterios de irreducibilidad y conexión para los límites proyectivos de preesquemas}
\label{subsection:IV.8.4}

\begin{proposition}[8.4.1]
\label{IV.8.4.1}
Supongamos que existe un índice $\alpha$ tal que $\mathrm{S}_\alpha$ es cuasi-compacto.
\begin{enumerate}[label=(\roman*)]
  \item Si $\mathrm{S}$ no es irreducible y si, además, el espacio subyacente de $\mathrm{S}$ es noetheriano y $\mathrm{S}_\alpha$ es cuasi-separado, existe $\lambda\geqslant\alpha$ tal que, para $\mu\geqslant\lambda$, $\mathrm{S}_\mu$ no es irreducible.
  \item Si $\mathrm{S}$ no es conexo, existe $\lambda\geqslant\alpha$ tal que, para $\mu\geqslant\lambda$, $\mathrm{S}_\mu$ no es conexo.
\end{enumerate}
\end{proposition}

\begin{proof}
Supongamos que $\mathrm{S}$ es la unión de dos subconjuntos cerrados distintos $\mathrm{S}'$, $\mathrm{S}''$ de $\mathrm{S}$ (resp.\ de dos subconjuntos cerrados no vacíos y disjuntos).
En el caso (i), $\mathrm{S}'$ y $\mathrm{S}''$ son constructibles, puesto que el espacio $\mathrm{S}$ es noetheriano.
En virtud de \sref{IV.8.3.11}, existen $\lambda\geqslant\alpha$ y dos subconjuntos cerrados constructibles $\mathrm{S}'_\lambda$, $\mathrm{S}''_\lambda$ de $\mathrm{S}_\lambda$ tales que $\mathrm{S}'=u_\lambda^{-1}(\mathrm{S}'_\lambda)$, $\mathrm{S}''=u_\lambda^{-1}(\mathrm{S}''_\lambda)$; como $\mathrm{S}=\mathrm{S}'\cup\mathrm{S}''$, se deduce de \sref{IV.8.3.11} que se puede suponer $\mathrm{S}_\lambda=\mathrm{S}'_\lambda\cup\mathrm{S}''_\lambda$; como $\mathrm{S}'_\lambda$ y $\mathrm{S}''_\lambda$ son distintos de $\mathrm{S}_\lambda$, esto prueba que $\mathrm{S}_\lambda$ no es irreducible.

En el caso (ii), $\mathrm{S}'$ y $\mathrm{S}''$ son abiertos y cuasi-compactos; por tanto, en virtud de \sref{IV.8.2.11}, existen $\lambda\geqslant\alpha$ y dos subconjuntos abiertos cuasi-compactos $\mathrm{S}'_\lambda$, $\mathrm{S}''_\lambda$ de $\mathrm{S}_\lambda$ tales que $\mathrm{S}'=u_\lambda^{-1}(\mathrm{S}'_\lambda)$, $\mathrm{S}''=u_\lambda^{-1}(\mathrm{S}''_\lambda)$.
Además, como $\mathrm{S}'$ y $\mathrm{S}''$ son abiertos y cerrados en $\mathrm{S}$, ambos son proconstructibles e indconstructibles \sref{IV.1.9.6}, y por tanto constructibles \sref{IV.1.9.11}; se deduce entonces de \sref{IV.8.3.5} que se puede suponer $\lambda$ elegido de modo que $\mathrm{S}_\lambda=\mathrm{S}'_\lambda\cup\mathrm{S}''_\lambda$ y $\mathrm{S}'_\lambda\cap\mathrm{S}''_\lambda=\varnothing$, lo que muestra que $\mathrm{S}_\lambda$ no es conexo.
\end{proof}

\begin{proposition}[8.4.2]
\label{IV.8.4.2}
Supongamos que el espacio subyacente de $\mathrm{S}$ es noetheriano y que se satisface una de las dos condiciones siguientes:
\begin{enumerate}[label=\emph{\alph*)}]
  \item Para $\lambda\leqslant\mu$, $u_{\lambda\mu}:\mathrm{S}_\mu\to\mathrm{S}_\lambda$ es dominante, y existe $\alpha$ tal que $\mathrm{S}_\alpha$ es cuasi-compacto.
\oldpage[IV-3]{18}
  \item Existe $\alpha$ tal que el espacio subyacente de $\mathrm{S}_\alpha$ es noetheriano y, para $\mu\geqslant\lambda$, $u_{\lambda\mu}$ es un homeomorfismo de $\mathrm{S}_\mu$ sobre un subespacio de $\mathrm{S}_\lambda$.
\end{enumerate}
En estas condiciones, existe $\lambda$ tal que, para todo $\mu\geqslant\lambda$:
\begin{enumerate}[label=(\roman*)]
  \item Para toda componente irreducible $\mathrm{Y}_i$ de $\mathrm{S}$ ($1\leqslant i\leqslant m$), $\overline{u_\mu(\mathrm{Y}_i)}$ es una componente irreducible de $\mathrm{S}_\mu$, y la aplicación $\mathrm{Y}_i\mapsto\overline{u_\mu(\mathrm{Y}_i)}$ es una biyección del conjunto de las componentes irreducibles de $\mathrm{S}$ sobre el conjunto de las componentes irreducibles de $\mathrm{S}_\mu$.
  \item Para toda componente conexa $\mathrm{C}_j$ de $\mathrm{S}$ ($1\leqslant j\leqslant n$), $\overline{u_\mu(\mathrm{C}_j)}$ es una componente conexa de $\mathrm{S}_\mu$, y la aplicación $\mathrm{C}_j\mapsto\overline{u_\mu(\mathrm{C}_j)}$ es una biyección del conjunto de las componentes conexas de $\mathrm{S}$ sobre el conjunto de las componentes conexas de $\mathrm{S}_\mu$.
\end{enumerate}
\end{proposition}

\begin{proof}
Demostremos primero el

\begin{lemma}[8.4.2.1]
\label{IV.8.4.2.1}
Bajo la condición \emph{a)} o \emph{b)} de \sref{IV.8.4.2}, existe $\lambda$ tal que, para $\mu\geqslant\lambda$, $u_\mu:\mathrm{S}\to\mathrm{S}_\mu$ es dominante.
\end{lemma}

En el caso \emph{a)}, esto ya se ha demostrado sin suponer noetheriano el espacio $\mathrm{S}$ \sref{IV.8.3.8},~(i).
En el caso \emph{b)}, pongamos $\mathrm{Z}_\alpha=\overline{u_\alpha(\mathrm{S})}$; como subconjunto cerrado del espacio noetheriano $\mathrm{S}_\alpha$, $\mathrm{Z}_\alpha$ es constructible y, como $u_\alpha^{-1}(\mathrm{Z}_\alpha)=\mathrm{S}$, se deduce de \sref{IV.8.3.5} que existe $\lambda\geqslant\alpha$ tal que, para $\mu\geqslant\lambda$, $\mathrm{Z}_\mu=u_{\alpha\mu}^{-1}(\mathrm{Z}_\alpha)=\mathrm{S}_\mu$.
Pero, como $u_{\alpha\mu}$ es un homeomorfismo de $\mathrm{S}_\mu$ sobre un subespacio de $\mathrm{S}_\alpha$ y la aplicación compuesta $\mathrm{S}\to\mathrm{Z}_\mu\to\mathrm{Z}_\alpha$ es dominante, lo mismo vale para $\mathrm{S}\to\mathrm{Z}_\mu=\mathrm{S}_\mu$.

Demostrado este lema, se puede suponer que, para todo $\lambda$, $u_\lambda$ es un morfismo dominante.
\begin{enumerate}[label=(\roman*)]
  \item Cada $\mathrm{S}_\lambda$ es la unión de los $\overline{u_\lambda(\mathrm{Y}_i)}$, que son irreducibles.
Por otro lado, si $\mathrm{U}_i$ es el abierto de $\mathrm{S}$ complementario de la unión de los $\mathrm{Y}_j$ de índice $j\neq i$ ($1\leqslant i\leqslant m$), los $\mathrm{U}_i$ son dos a dos disjuntos y $\mathrm{Y}_i=\overline{\mathrm{U}_i}$ \sref[0]{0.2.1.6}.
Como el espacio subyacente de $\mathrm{S}$ es noetheriano, los $\mathrm{U}_i$ son cuasi-compactos; por tanto, existen un índice $\lambda$ y abiertos $\mathrm{U}_{i\lambda}$ de $\mathrm{S}_\lambda$ tales que $\mathrm{U}_i=u_\lambda^{-1}(\mathrm{U}_{i\lambda})$ para $1\leqslant i\leqslant m$ \sref{IV.8.2.11}.
De esto se deduce que, si se pone $\mathrm{U}_{i\mu}=u_{\lambda\mu}^{-1}(\mathrm{U}_{i\lambda})$ para $\lambda\leqslant\mu$, los $\mathrm{U}_{i\mu}$ son dos a dos disjuntos, pues los $\mathrm{U}_i=u_\mu^{-1}(\mathrm{U}_{i\mu})$ lo son y $u_\mu$ es dominante.
En consecuencia, ninguna de las clausuras $\overline{\mathrm{U}_{i\mu}}$ está contenida en otra, y $u_\mu(\mathrm{U}_i)$ es denso en $\mathrm{U}_{i\mu}$ puesto que $u_\mu$ es dominante; de ahí $\overline{\mathrm{U}_{i\mu}}=\overline{u_\mu(\mathrm{Y}_i)}$, lo que prueba que los $\overline{\mathrm{U}_{i\mu}}$ son las componentes irreducibles de $\mathrm{S}_\mu$ \sref[0]{0.2.1.7} y concluye la demostración.

  \item Como el espacio subyacente $\mathrm{S}$ es noetheriano, los $\mathrm{C}_j$ son abiertos y cerrados en $\mathrm{S}$ y cuasi-compactos; el mismo razonamiento que en (i) muestra que existen $\lambda$ y abiertos $\mathrm{V}_{j\lambda}$ de $\mathrm{S}_\lambda$ tales que $\mathrm{C}_j=u_\lambda^{-1}(\mathrm{V}_{j\lambda})$ para $1\leqslant j\leqslant n$.
Se ve asimismo, como en (i), que si se pone $\mathrm{V}_{j\mu}=u_{\lambda\mu}^{-1}(\mathrm{V}_{j\lambda})$ para $\mu\geqslant\lambda$, los $\mathrm{V}_{j\mu}$ son dos a dos disjuntos y $u_\mu(\mathrm{C}_j)$ es denso en $\mathrm{V}_{j\mu}$; esto implica que $\mathrm{V}_{j\mu}$ es conexo.
Además, se deduce de \sref{IV.8.3.4} que, para $\mu$ suficientemente grande, la unión de los $\mathrm{V}_{j\mu}$ es $\mathrm{S}_\mu$, ya que todo abierto de un preesquema es indconstructible \sref{IV.1.9.6}.
Los $\mathrm{V}_{j\mu}$ son, pues, las componentes conexas de $\mathrm{S}_\mu$, y esto concluye la demostración.
\end{enumerate}
\end{proof}

Observemos que, si los morfismos $u_{\lambda\mu}$ son \emph{inmersiones}, satisfacen en particular la condición \emph{b)} de \sref{IV.8.4.2}.

\oldpage[IV-3]{19}
\begin{corollary}[8.4.3]
\label{IV.8.4.3}
Supongamos que se satisface una de las condiciones \emph{a)}, \emph{b)} de \sref{IV.8.4.2}, siendo noetheriano el espacio subyacente de $\mathrm{S}$; entonces, para que $\mathrm{S}$ sea irreducible, es necesario y suficiente que exista $\lambda$ tal que $\mathrm{S}_\mu$ lo sea para todo $\mu\geqslant\lambda$.
\end{corollary}

\begin{proposition}[8.4.4]
\label{IV.8.4.4}
Supongamos que existe $\alpha$ tal que $\mathrm{S}_\alpha$ es cuasi-compacto y que, para $\lambda\leqslant\mu$, $u_{\lambda\mu}:\mathrm{S}_\mu\to\mathrm{S}_\lambda$ es dominante.
Entonces, para que $\mathrm{S}$ sea conexo, es necesario y suficiente que exista $\lambda$ tal que $\mathrm{S}_\mu$ lo sea para todo $\mu\geqslant\lambda$.
\end{proposition}

\begin{proof}
La condición es suficiente en virtud de \sref{IV.8.4.1}; por otro lado, hemos visto en \sref{IV.8.3.8},~(i), que $u_\lambda:\mathrm{S}\to\mathrm{S}_\lambda$ es dominante para $\lambda$ suficientemente grande; por tanto, si $\mathrm{S}$ es conexo, lo mismo ocurre con $\mathrm{S}_\lambda$, ya que $u_\lambda(\mathrm{S})$ es denso en $\mathrm{S}_\lambda$ y es conexo.
\end{proof}

\begin{corollary}[8.4.5]
\label{IV.8.4.5}
Sea $k$ un cuerpo y $\mathrm{X}$ un preesquema $k$ cuasi-compacto.
Para que $\mathrm{X}$ sea geométricamente conexo \sref{IV.4.5.2}, es necesario y suficiente que, para toda extensión finita y separable $\mathrm{K}$ de $k$, $\mathrm{X}\otimes_k\mathrm{K}$ sea conexo.
\end{corollary}

\begin{proof}
La condición es trivialmente necesaria.
Para ver que es suficiente, debemos demostrar que si $\Omega$ es una clausura algebraica de $k$, $\mathrm{X}\otimes_k\Omega$ es conexo \sref{IV.4.5.1}.
Ahora bien, $\Omega$ es un límite inductivo filtrante de las subextensiones finitas $k'$ de $k$, y para $k\subset k'\subset k''\subset\Omega$, el morfismo $\mathrm{X}\otimes_k k''\to\mathrm{X}\otimes_k k'$ es sobreyectivo.
Así pues, en virtud de \sref{IV.8.4.4}, basta demostrar que $\mathrm{X}\otimes_k k'$ es conexo para toda extensión finita $k'$ de $k$.
Pero si $\mathrm{K}$ es la mayor extensión separable contenida en $k'$, el morfismo $\mathrm{X}\otimes_k k'\to\mathrm{X}\otimes_k\mathrm{K}$ es finito, sobreyectivo y radicial, y por tanto \sref{IV.2.4.5} un homeomorfismo; como $\mathrm{X}\otimes_k\mathrm{K}$ es conexo por hipótesis, también lo es $\mathrm{X}\otimes_k k'$.
\end{proof}

\begin{remark}[8.4.6]
\label{IV.8.4.6}
\begin{enumerate}[label=(\roman*)]
  \item La demostración de \sref{IV.8.4.2} muestra que la conclusión de esta proposición sigue siendo válida si se supone que el espacio subyacente de $\mathrm{S}$ es noetheriano, que existe $\alpha$ tal que $\mathrm{S}_\alpha$ es cuasi-compacto y, finalmente, que los $u_\lambda:\mathrm{S}\to\mathrm{S}_\lambda$ son dominantes.
  \item Por otro lado, la conclusión de \sref{IV.8.4.2} puede fallar cuando los $u_\lambda$ no son dominantes para $\lambda$ suficientemente grande, incluso cuando los $\mathrm{S}_\lambda$ y $\mathrm{S}$ son noetherianos, como muestra el ejemplo siguiente.
Tómese $\mathbf{N}$ como conjunto de índices, todos los $\mathrm{S}_n$ iguales a $\operatorname{Spec}(\mathrm{A}\times\mathrm{K})=\operatorname{Spec}(\mathrm{A})\amalg\operatorname{Spec}(\mathrm{K})$, donde $\mathrm{K}$ es un cuerpo, $\mathrm{A}$ cualquier $\mathrm{K}$-álgebra, y todos los morfismos $u_{n,n+1}$ iguales al mismo morfismo correspondiente al homomorfismo $(x,y)\mapsto(j(y),y)$ de $\mathrm{A}\times\mathrm{K}$ en sí misma, donde $j:\mathrm{K}\to\mathrm{A}$ es el homomorfismo canónico.
Se verifica fácilmente que el límite inductivo de este sistema de anillos es $\mathrm{K}$, siendo el homomorfismo canónico $\mathrm{A}\times\mathrm{K}\to\mathrm{K}$ la segunda proyección.
Así se ve que $\mathrm{S}=\operatorname{Spec}(\mathrm{K})$ es irreducible aunque ninguno de los $\mathrm{S}_n$ sea conexo.
\end{enumerate}
\end{remark}

\subsection{Módulos de presentación finita sobre un límite proyectivo de preesquemas}
\label{subsection:IV.8.5}

\begin{env}[8.5.1]
\label{IV.8.5.1}
Conservamos siempre las notaciones de \sref{IV.8.2.2}; además nos restringiremos al caso en que $\mathrm{S}_0$ sea uno de los $\mathrm{S}_\lambda$, al que siempre se puede reducir.

\emph{Siempre que, en esta sección, consideremos una familia $(\sh{F}_\lambda)$, donde, para cada $\lambda\in\mathrm{L}$, $\sh{F}_\lambda$ es un Módulo sobre $\sh{O}_{\mathrm{S}_\lambda}$, se entenderá que esta familia satisface la condición}
\begin{equation}
  \sh{F}_\mu=u_{\lambda\mu}^*(\sh{F}_\lambda) \qquad\text{para }\lambda\leqslant\mu.
  \tag{8.5.1.1}
\label{IV.8.5.1.1}
\end{equation}
\oldpage[IV-3]{20}
Pondremos entonces
\begin{equation}
  \sh{F}=u_\lambda^*(\sh{F}_\lambda)
  \label{IV.8.5.1.2}
\tag{8.5.1.2}
\end{equation}
que es un Módulo sobre $\sh{O}_\mathrm{S}$ que no depende del índice $\lambda\in\mathrm{L}$, en virtud de la hipótesis \sref{IV.8.5.1.1}.

Sean $(\sh{F}_\lambda)$ y $(\sh{G}_\lambda)$ dos familias de ese tipo de Módulos sobre $\sh{O}_{\mathrm{S}_\lambda}$.
Es claro que las aplicaciones $f_\lambda\mapsto u_{\lambda\mu}^*(f_\lambda)$ de $\operatorname{Hom}_{\sh{O}_{\mathrm{S}_\lambda}}(\sh{F}_\lambda,\sh{G}_\lambda)$ en $\operatorname{Hom}_{\sh{O}_{\mathrm{S}_\mu}}(\sh{F}_\mu,\sh{G}_\mu)$ definen un \emph{sistema inductivo} de grupos abelianos $(\operatorname{Hom}_{\sh{O}_{\mathrm{S}_\lambda}}(\sh{F}_\lambda,\sh{G}_\lambda),u_{\lambda\mu}^*)$, y que las aplicaciones $f_\lambda\mapsto u_\lambda^*(f_\lambda)$ forman un sistema inductivo de homomorfismos de grupos abelianos
\[
  u_\lambda^*:\operatorname{Hom}_{\sh{O}_{\mathrm{S}_\lambda}}(\sh{F}_\lambda,\sh{G}_\lambda)\to\operatorname{Hom}_\mathrm{S}(\sh{F},\sh{G});
\]
de donde, pasando al límite inductivo, se obtiene un homomorfismo canónico de grupos abelianos
\begin{equation}
  u_{\sh{F},\sh{G}}^*:\varinjlim\operatorname{Hom}_{\sh{O}_{\mathrm{S}_\lambda}}(\sh{F}_\lambda,\sh{G}_\lambda)\to\operatorname{Hom}_\mathrm{S}(\sh{F},\sh{G}).
  \tag{8.5.1.3}
\label{IV.8.5.1.3}
\end{equation}

Obsérvese que cuando $\sh{F}_\lambda=\sh{O}_{\mathrm{S}_\lambda}$ se satisface la condición \sref{IV.8.5.1.1}, y se tiene $\sh{F}=\sh{O}_\mathrm{S}$; por tanto, el homomorfismo \sref{IV.8.5.1.3} proporciona \sref[0]{0.5.1.1} un homomorfismo canónico de grupos abelianos
\begin{equation}
  u_{\sh{G}}:\varinjlim\Gamma(\mathrm{S}_\lambda,\sh{G}_\lambda)\to\Gamma(\mathrm{S},\sh{G}).
  \tag{8.5.1.4}
\label{IV.8.5.1.4}
\end{equation}
\end{env}

\begin{theorem}[8.5.2]
\label{IV.8.5.2}
\begin{enumerate}[label=(\roman*)]
  \item Supóngase que $\mathrm{S}_0$ es cuasi-compacto (resp.\ cuasi-compacto y cuasi-separado) y que, para algún $\lambda\in\mathrm{L}$, $\sh{F}_\lambda$ es cuasi-coherente y de tipo finito (resp.\ de presentación finita) y $\sh{G}_\lambda$ cuasi-coherente.
Entonces el homomorfismo $u_{\sh{F},\sh{G}}^*$ es inyectivo (resp.\ biyectivo).
  \item Supóngase que $\mathrm{S}_0$ es cuasi-compacto y cuasi-separado.
Para todo Módulo sobre $\sh{O}_\mathrm{S}$ cuasi-coherente y de presentación finita $\sh{F}$, existen $\lambda\in\mathrm{L}$ y un Módulo sobre $\sh{O}_{\mathrm{S}_\lambda}$ cuasi-coherente y de presentación finita $\sh{F}_\lambda$ tales que $\sh{F}$ es isomorfo a $u_\lambda^*(\sh{F}_\lambda)$.
\end{enumerate}
\end{theorem}

\begin{proof}
\begin{enumerate}[label=(\roman*)]
  \item Podemos evidentemente restringirnos al caso $\mathrm{S}_0=\mathrm{S}_\lambda$, ya que los morfismos $u_{0\lambda}:\mathrm{S}_\lambda\to\mathrm{S}_0$ son afines, y por tanto cuasi-compactos y separados.
Consideremos primero el caso en que $\mathrm{S}_0=\operatorname{Spec}(\mathrm{A}_0)$ es \emph{afín}.
Entonces la afirmación (i) es equivalente al

\begin{lemma}[8.5.2.1]
\label{IV.8.5.2.1}
Sea $\mathrm{A}_0$ un anillo, $(\mathrm{A}_\lambda)_{\lambda\in\mathrm{L}}$ un sistema inductivo de $\mathrm{A}_0$-álgebras, $\mathrm{A}=\varinjlim\mathrm{A}_\lambda$; sean $\mathrm{M}_0$, $\mathrm{N}_0$ dos $\mathrm{A}_0$-módulos, y pongamos $\mathrm{M}_\lambda=\mathrm{M}_0\otimes_{\mathrm{A}_0}\mathrm{A}_\lambda$, $\mathrm{N}_\lambda=\mathrm{N}_0\otimes_{\mathrm{A}_0}\mathrm{A}_\lambda$, $\mathrm{M}=\mathrm{M}_0\otimes_{\mathrm{A}_0}\mathrm{A}=\varinjlim\mathrm{M}_\lambda$, $\mathrm{N}=\mathrm{N}_0\otimes_{\mathrm{A}_0}\mathrm{A}=\varinjlim\mathrm{N}_\lambda$.
Si $\mathrm{M}_0$ es de tipo finito (resp.\ de presentación finita), el homomorfismo canónico
\begin{equation}
  \varinjlim\operatorname{Hom}_{\mathrm{A}_\lambda}(\mathrm{M}_\lambda,\mathrm{N}_\lambda)\to\operatorname{Hom}_\mathrm{A}(\mathrm{M},\mathrm{N})
  \tag{8.5.2.2}
\label{IV.8.5.2.2}
\end{equation}
es inyectivo (resp.\ biyectivo).
\end{lemma}

Sabemos, en efecto (Bourbaki, \emph{Alg.}, cap.~II, 3\textsuperscript{a}~ed., \textsection\,5, n\textsuperscript{o}\,1), que existen isomorfismos canónicos functoriales
\[
  \operatorname{Hom}_{\mathrm{A}_\lambda}(\mathrm{M}_\lambda,\mathrm{N}_\lambda)\xrightarrow{\sim}\operatorname{Hom}_{\mathrm{A}_0}(\mathrm{M}_0,\mathrm{N}_\lambda),\qquad\operatorname{Hom}_\mathrm{A}(\mathrm{M},\mathrm{N})\xrightarrow{\sim}\operatorname{Hom}_{\mathrm{A}_0}(\mathrm{M}_0,\mathrm{N})
\]
de modo que el homomorfismo \sref{IV.8.5.2.2} no es otro, salvo isomorfismos canónicos, que el homomorfismo canónico
\begin{equation}
  \varinjlim\operatorname{Hom}_{\mathrm{A}_0}(\mathrm{M}_0,\mathrm{N}_\lambda)\to\operatorname{Hom}_{\mathrm{A}_0}(\mathrm{M}_0,\varinjlim\mathrm{N}_\lambda),
  \tag{8.5.2.3}
\label{IV.8.5.2.3}
\end{equation}
\oldpage[IV-3]{21}
que asocia a cada sistema inductivo de homomorfismos de $\mathrm{A}_0$-módulos $\theta_\lambda:\mathrm{M}_0\to\mathrm{N}_\lambda$ su límite inductivo.

Ahora bien, si $\mathrm{M}_0$ es de tipo finito (resp.\ de presentación finita), existe una sucesión exacta $\mathrm{A}_0^m\to\mathrm{M}_0\to 0$ (resp.\ $\mathrm{A}_0^n\to\mathrm{A}_0^m\to\mathrm{M}_0\to 0$); como es claro que \sref{IV.8.5.2.3} es biyectivo cuando $\mathrm{M}_0$ es de la forma $\mathrm{A}_0^q$, basta utilizar la exactitud por la izquierda del functor $\mathrm{M}_0\mapsto\operatorname{Hom}_{\mathrm{A}_0}(\mathrm{M}_0,\mathrm{P})$ y la exactitud del functor $\varinjlim$ (en la categoría de grupos abelianos) para concluir.

\medskip
Pasemos al caso en que $\mathrm{S}_0$ es cuasi-compacto, y sea $(\mathrm{U}_i)$ un recubrimiento finito de $\mathrm{S}_0$ por abiertos afines; para cada $\lambda$, los $\mathrm{U}_{i\lambda}=u_{0\lambda}^{-1}(\mathrm{U}_i)$ forman un recubrimiento abierto afín de $\mathrm{S}_\lambda$, y los $\mathrm{V}_i=u_0^{-1}(\mathrm{U}_i)$ un recubrimiento abierto afín de $\mathrm{S}$.
Para demostrar que $u_{\sh{F},\sh{G}}^*$ es inyectivo, basta demostrar que si $f_\lambda:\sh{F}_\lambda\to\sh{G}_\lambda$ satisface $f=u_\lambda^*(f_\lambda)=0$, entonces existe $\mu\geqslant\lambda$ tal que $u_{\lambda\mu}^*(f_\lambda)=0$.
En virtud del lema \sref{IV.8.5.2.1}, para cada $i$ existe un índice $\lambda_i$ tal que $f_\mu|_{\mathrm{U}_{i\mu}}=0$ para $\mu\geqslant\lambda_i$.
Por tanto basta tomar $\mu$ mayor que todos los $\lambda_i$.

Supongamos además que $\mathrm{S}_0$ es cuasi-separado y que $\sh{F}_0$ es de presentación finita, y sea $f:\sh{F}\to\sh{G}$ un homomorfismo de Módulos sobre $\sh{O}_\mathrm{S}$.
En virtud del lema \sref{IV.8.5.2.1}, para cada $i$ existen un índice $\lambda_i$ y un homomorfismo $f_{\lambda_i}^{(i)}:\sh{F}_{\lambda_i}|\mathrm{U}_{i,\lambda_i}\to\sh{G}_{\lambda_i}|\mathrm{U}_{i,\lambda_i}$ tales que $u_{\lambda_i}^*(f_{\lambda_i}^{(i)})=f|\mathrm{V}_i$.
Como $\mathrm{L}$ es filtrado, podemos suponer además que todos los $\lambda_i$ son iguales a un único $\lambda$.
Ahora $\mathrm{S}_\lambda$ es cuasi-separado \sref{IV.1.2.3}, y $\sh{F}_\lambda$ es un Módulo sobre $\sh{O}_{\mathrm{S}_\lambda}$ cuasi-coherente y de presentación finita \sref[0]{0.5.2.5}; puesto que, para cada par de índices $i,j$, $\mathrm{U}_{ij\lambda}=\mathrm{U}_{i\lambda}\cap\mathrm{U}_{j\lambda}$ es cuasi-compacto y, por definición,
\[
  u_\lambda^*(f_\lambda^{(i)}|\mathrm{U}_{ij\lambda})=u_\lambda^*(f_\lambda^{(j)}|\mathrm{U}_{ij\lambda})=f|(\mathrm{V}_i\cap\mathrm{V}_j),
\]
se deduce de lo demostrado anteriormente que existe un índice $\lambda_{ij}$ tal que
\[
  u_{\lambda\mu}^*(f_\lambda^{(i)}|\mathrm{U}_{ij\lambda})=u_{\lambda\mu}^*(f_\lambda^{(j)}|\mathrm{U}_{ij\lambda})\quad\text{para }\mu\geqslant\lambda_{ij}.
\]
Tomando $\mu$ mayor que todos los $\lambda_{ij}$, se ve que $u_{\lambda\mu}^*(f_\lambda^{(i)})$ y $u_{\lambda\mu}^*(f_\lambda^{(j)})$ coinciden en $\mathrm{U}_{i\mu}\cap\mathrm{U}_{j\mu}$ para cada par $(i,j)$, y por tanto definen un homomorfismo $f_\mu:\sh{F}_\mu\to\sh{G}_\mu$ tal que $f=u_\mu^*(f_\mu)$.

\end{enumerate}
\end{proof}

Antes de pasar a la demostración de (ii), observemos los corolarios siguientes de (i):

\begin{corollary}[8.5.2.4]
\label{IV.8.5.2.4}
Supongamos que $\mathrm{S}_0$ es cuasi-compacto, $\sh{F}_\lambda$ cuasi-coherente de tipo finito y $\sh{G}_\lambda$ cuasi-coherente de presentación finita.
Sea $f_\lambda:\sh{F}_\lambda\to\sh{G}_\lambda$ un homomorfismo.
Para que $f=u_\lambda^*(f_\lambda):\sh{F}\to\sh{G}$ sea un isomorfismo, es necesario y suficiente que exista $\mu\geqslant\lambda$ tal que $f_\mu=u_{\lambda\mu}^*(f_\lambda):\sh{F}_\mu\to\sh{G}_\mu$ lo sea.
\end{corollary}

\begin{proof}
Siempre podemos suponer $\mathrm{S}_\lambda=\mathrm{S}_0$; como la cuestión es local en $\mathrm{S}_0$, podemos además (puesto que $\mathrm{S}_0$ es cuasi-compacto y $\mathrm{L}$ es filtrado) reducirnos al caso en que $\mathrm{S}_0$ es afín, y por tanto cuasi-separado.
La condición es trivialmente suficiente; queda demostrar que es necesaria: por hipótesis existe un homomorfismo de $\sh{O}_\mathrm{S}$-Módulos $g:\sh{G}\to\sh{F}$ tal que $g\circ f=1_\sh{F}$ y $f\circ g=1_\sh{G}$.
Como $\sh{G}$ es de presentación finita, existen un índice $\nu\geqslant\lambda$ y un homomorfismo $g_\nu:\sh{G}_\nu\to\sh{F}_\nu$ tales que $g=u_\nu^*(g_\nu)$ en virtud de \sref{IV.8.5.2},~(i); por tanto se tiene
\[
  u_\nu^*(g_\nu\circ f_\nu)=1_\sh{F}
\]
\oldpage[IV-3]{22}
y $u_\nu^*(f_\nu\circ g_\nu)=1_\sh{G}$; como $\sh{F}_\nu$ y $\sh{G}_\nu$ son de tipo finito, se deduce de \sref{IV.8.5.2},~(i), que existe $\mu\geqslant\nu$ tal que $g_\mu\circ f_\mu=1_{\sh{F}_\mu}$ y $f_\mu\circ g_\mu=1_{\sh{G}_\mu}$, de donde se obtiene el corolario.
\end{proof}

\begin{corollary}[8.5.2.5]
\label{IV.8.5.2.5}
Supongamos que $\mathrm{S}_0$ es cuasi-compacto y cuasi-separado.
Supongamos que $\sh{F}_\lambda$ y $\sh{G}_\lambda$ son Módulos sobre $\sh{O}_{\mathrm{S}_\lambda}$ cuasi-coherentes y de presentación finita.
Para que $\sh{F}$ y $\sh{G}$ sean isomorfos, es necesario y suficiente que exista $\mu\geqslant\lambda$ tal que $\sh{F}_\mu$ y $\sh{G}_\mu$ sean isomorfos.
Además, para todo isomorfismo $f:\sh{F}\xrightarrow{\sim}\sh{G}$, existen $\nu\geqslant\mu$ y un isomorfismo $f_\nu:\sh{F}_\nu\xrightarrow{\sim}\sh{G}_\nu$ tales que $f=u_\nu^*(f_\nu)$.
\end{corollary}

\begin{proof}
Esto se deduce de \sref{IV.8.5.2.4} y de \sref{IV.8.5.2},~(i), ya que todo homomorfismo $f:\sh{F}\to\sh{G}$ es de la forma $u_\mu^*(f_\mu)$ para algún $\mu\geqslant\lambda$ y algún homomorfismo $f_\mu:\sh{F}_\mu\to\sh{G}_\mu$.
\end{proof}

\begin{proof}[Demostración de \sref{IV.8.5.2},~(ii)]
Consideremos de nuevo primero el caso en que $\mathrm{S}_0=\operatorname{Spec}(\mathrm{A}_0)$ es afín.
Entonces la afirmación es equivalente al lema \sref{IV.5.13.7.1}.

En el caso general, partiendo de un recubrimiento abierto afín finito $(\mathrm{U}_i)$ de $\mathrm{S}_0$, se deduce de \sref{IV.5.13.7.1} que para cada $i$ existen un índice $\lambda(i)$ y un Módulo sobre $\sh{O}_{\mathrm{U}_{i,\lambda(i)}}$ cuasi-coherente y de presentación finita $\sh{F}_i$ tales que $u_{\lambda(i)}^*(\sh{F}_i)=\sh{F}|_{\mathrm{V}_i}$ (con las notaciones de (i)).
Además, como $\mathrm{L}$ es filtrado, podemos suponer que todos los $\lambda_i$ son iguales a un único $\lambda$.
Como $\mathrm{U}_{ij\lambda}=\mathrm{U}_{i\lambda}\cap\mathrm{U}_{j\lambda}$ es cuasi-compacto y cuasi-separado \sref{IV.1.2.7}, se deduce de \sref{IV.8.5.2.5} que para cada par $(i,j)$ existen un índice $\lambda_{ij}\geqslant\lambda$ y un isomorfismo
\[
  \theta_{ij}:u_{\lambda,\lambda_{ij}}^*(\sh{F}_i|\mathrm{U}_{ij\lambda})\xrightarrow{\sim}u_{\lambda,\lambda_{ij}}^*(\sh{F}_j|\mathrm{U}_{ij\lambda})
\]
tal que $u_{\lambda_{ij}}^*(\theta_{ij})$ es el automorfismo identidad de $\sh{F}|(\mathrm{V}_i\cap\mathrm{V}_j)$; podemos suponer además que todos los $\lambda_{ij}$ son iguales a $\lambda$.
Cambiando las notaciones, podemos suponer por tanto que para cada par $(i,j)$ existe un isomorfismo $\theta_{ij}:\sh{F}_i|_{\mathrm{U}_{ij\lambda}}\xrightarrow{\sim}\sh{F}_j|_{\mathrm{U}_{ij\lambda}}$ tal que $u_\lambda^*(\theta_{ij})$ es el automorfismo identidad de $\sh{F}|_{(\mathrm{V}_i\cap\mathrm{V}_j)}$.
Finalmente, para cualesquiera tres índices $i$, $j$, $k$, si ponemos $\mathrm{U}_{ijk,\lambda}=\mathrm{U}_{i\lambda}\cap\mathrm{U}_{j\lambda}\cap\mathrm{U}_{k\lambda}$, $\mathrm{U}_{ijk,\lambda}$ es cuasi-compacto, y si $\theta'_{ij}$, $\theta'_{jk}$ y $\theta'_{ik}$ denotan las restricciones de $\theta_{ij}$, $\theta_{jk}$ y $\theta_{ik}$ a $\mathrm{U}_{ijk,\lambda}$, se tiene $u_\lambda^*(\theta'_{ij}\circ\theta'_{jk})=u_\lambda^*(\theta'_{ik})$.
Por tanto, en virtud de (i), existe un índice $\mu\geqslant\lambda$ tal que $u_{\lambda\mu}^*(\theta'_{ik})=u_{\lambda\mu}^*(\theta'_{ij}\circ\theta'_{jk})$; en consecuencia, los isomorfismos $u_{\lambda\mu}^*(\theta'_{ij}):u_{\lambda\mu}^*(\sh{F}_i)|_{\mathrm{U}_{ij\mu}}\xrightarrow{\sim}u_{\lambda\mu}^*(\sh{F}_j)|_{\mathrm{U}_{ij\mu}}$ satisfacen la condición de pegado y definen en $\mathrm{S}_\mu$ un Módulo sobre $\sh{O}_{\mathrm{S}_\mu}$ cuasi-coherente y de presentación finita $\sh{F}_\mu$ \sref[0]{0.3.3.1} tal que $\sh{F}$ es isomorfo a $u_\mu^*(\sh{F}_\mu)$.
\end{proof}

\begin{env}[8.5.3]
\label{IV.8.5.3}
El resultado de \sref{IV.8.5.2} también puede expresarse diciendo que, si $\mathrm{S}_0$ es cuasi-compacto y cuasi-separado, la categoría de Módulos sobre $\sh{O}_\mathrm{S}$ cuasi-coherentes y de presentación finita queda determinada, salvo equivalencia, por los datos de las categorías de Módulos sobre $\sh{O}_{\mathrm{S}_\lambda}$ cuasi-coherentes y de presentación finita, los funtores $u_{\lambda\mu}^*$ entre estas categorías y los isomorfismos de transición $u_{\mu\nu}^*\circ u_{\lambda\mu}^*\xrightarrow{\sim}u_{\lambda\nu}^*$.
De manera figurada, puede decirse que el dato de un Módulo sobre $\sh{O}_\mathrm{S}$ cuasi-coherente y de presentación finita $\sh{F}$ es ``funcionalmente'' equivalente al dato de un Módulo sobre $\sh{O}_{\mathrm{S}_\lambda}$ cuasi-coherente y de presentación finita $\sh{F}'_\lambda$ para algún $\lambda$ suficientemente grande y que, si otro Módulo sobre $\sh{O}_{\mathrm{S}_\mu}$ cuasi-coherente y de presentación finita $\sh{F}''_\mu$ tiene también a $\sh{F}$ como imagen inversa, entonces $\sh{F}'_\lambda$ y $\sh{F}''_\mu$ tienen la misma imagen inversa en algún $\mathrm{S}_\nu$ adecuado ($\nu\geqslant\lambda$, $\nu\geqslant\mu$).

A partir de este punto interpretaremos desde este punto de vista diversas nociones relacionadas con los Módulos sobre $\sh{O}_\mathrm{S}$ cuasi-coherentes.
\end{env}

\oldpage[IV-3]{23}
\begin{corollary}[8.5.4]
\label{IV.8.5.4}
Supongamos que $\mathrm{S}_0$ es cuasi-compacto y cuasi-separado; entonces, para todo Módulo sobre $\sh{O}_{\mathrm{S}_\lambda}$ cuasi-coherente $\sh{G}_\lambda$, el homomorfismo canónico \sref{IV.8.5.1.4} es biyectivo.
\end{corollary}

\begin{proof}
En efecto, basta aplicar \sref{IV.8.5.2},~(i), al caso $\sh{F}_\lambda=\sh{O}_{\mathrm{S}_\lambda}$, que es de presentación finita.
\end{proof}

\begin{proposition}[8.5.5]
\label{IV.8.5.5}
Supongamos que $\mathrm{S}_0$ es cuasi-compacto y que $\sh{F}_\lambda$ es un Módulo sobre $\sh{O}_{\mathrm{S}_\lambda}$ cuasi-coherente y de presentación finita.
Para que $\sh{F}$ sea localmente libre (resp.\ localmente libre de rango $n$), es necesario y suficiente que exista $\mu\geqslant\lambda$ tal que $\sh{F}_\mu$ lo sea.
\end{proposition}

\begin{proof}
La condición es trivialmente suficiente; demostremos que es necesaria.
Si $\sh{F}$ es localmente libre (resp.\ localmente libre de rango $n$), existe un recubrimiento abierto afín finito $(\mathrm{V}_i)$ de $\mathrm{S}$ tal que $\sh{F}|_{\mathrm{V}_i}$ es isomorfo a $\sh{O}_\mathrm{S}^{n_i}|_{\mathrm{V}_i}$ (resp.\ $\sh{O}_\mathrm{S}^n|_{\mathrm{V}_i}$) para cada $i$.
En virtud de \sref{IV.8.2.11}, existen $\nu\geqslant\lambda$ y, para cada $i$, un abierto cuasi-compacto $\mathrm{U}_{i\nu}$ de $\mathrm{S}_\nu$ tal que $\mathrm{V}_i=u_\nu^{-1}(\mathrm{U}_{i\nu})$.
Como $\mathrm{S}_\nu$ es cuasi-compacto, cada $\mathrm{U}_{i\nu}$ es unión finita de afines abiertos; así nos reducimos al caso en que $\mathrm{S}_0$ es afín y $\sh{F}=\sh{O}_\mathrm{S}^n$; entonces sabemos que existe $\mu\geqslant\lambda$ tal que $\sh{F}_\mu$ es isomorfo a $\sh{O}_{\mathrm{S}_\mu}^n$ \sref{IV.8.5.2.5}.
\end{proof}

\begin{proposition}[8.5.6]
\label{IV.8.5.6}
Supongamos que $\mathrm{S}_0$ es cuasi-compacto y consideremos una sucesión
\[
  \sh{F}_\lambda\to\sh{G}_\lambda\to\sh{H}_\lambda\to 0
\]
de homomorfismos de Módulos sobre $\sh{O}_{\mathrm{S}_\lambda}$ cuasi-coherentes, donde $\sh{F}_\lambda$ y $\sh{G}_\lambda$ son de tipo finito y $\sh{H}_\lambda$ de presentación finita.
Para que la sucesión correspondiente $\sh{F}\to\sh{G}\to\sh{H}\to 0$ sea exacta, es necesario y suficiente que exista $\mu\geqslant\lambda$ tal que la sucesión $\sh{F}_\mu\to\sh{G}_\mu\to\sh{H}_\mu\to 0$ lo sea (en cuyo caso lo mismo vale para $\nu\geqslant\mu$).
\end{proposition}

\begin{proof}
El hecho de que la condición sea suficiente y la última afirmación se deben a que el functor $u_\lambda^*$ (resp.\ $u_{\mu\nu}^*$) es exacto por la derecha.
Para demostrar que la condición es necesaria, observemos que de la hipótesis y de \sref{IV.8.5.2},~(i), se deduce que existe $\nu\geqslant\lambda$ tal que la composición $\sh{F}_\nu\to\sh{G}_\nu\to\sh{H}_\nu$ es nula.
Si ponemos $\sh{H}'_\nu=\operatorname{Coker}(\sh{F}_\nu\to\sh{G}_\nu)$, tenemos por tanto un homomorfismo $f_\nu:\sh{H}'_\nu\to\sh{H}_\nu$; por hipótesis, $u_\nu^*(f_\nu)$ es un isomorfismo, y se deduce de \sref{IV.8.5.2.4} que existe $\mu\geqslant\nu$ tal que $u_{\nu\mu}^*(f_\nu)$ es un isomorfismo, lo que concluye la demostración.
\end{proof}

\begin{corollary}[8.5.7]
\label{IV.8.5.7}
Supongamos que $\mathrm{S}_0$ es cuasi-compacto, $\sh{F}_\lambda$ cuasi-coherente, $\sh{G}_\lambda$ cuasi-coherente de tipo finito, y sea $f_\lambda:\sh{F}_\lambda\to\sh{G}_\lambda$ un homomorfismo.
Para que $f=u_\lambda^*(f_\lambda):\sh{F}\to\sh{G}$ sea sobreyectivo, es necesario y suficiente que exista $\mu\geqslant\lambda$ tal que $f_\mu=u_{\lambda\mu}^*(f_\lambda):\sh{F}_\mu\to\sh{G}_\mu$ lo sea.
\end{corollary}

\begin{proof}
Este es el caso particular de \sref{IV.8.5.6} aplicado a la sucesión $0\to\sh{H}_\lambda\to 0\to 0$, donde $\sh{H}_\lambda=\operatorname{Coker}(f_\lambda)$, que es cuasi-coherente y de tipo finito (teniendo en cuenta que $\sh{H}=\operatorname{Coker}(f)$ y $\sh{H}_\mu=\operatorname{Coker}(f_\mu)$, en virtud de la exactitud por la derecha de los funtores $u_\lambda^*$ y $u_{\lambda\mu}^*$).
\end{proof}

\begin{corollary}[8.5.8]
\label{IV.8.5.8}
Supongamos que $\mathrm{S}_0$ es cuasi-compacto y que los morfismos $u_{\lambda\mu}:\mathrm{S}_\mu\to\mathrm{S}_\lambda$ son planos.
Entonces:
\begin{enumerate}[label=(\roman*)]
  \item Sea $\sh{F}_\lambda\xrightarrow{f_\lambda}\sh{G}_\lambda\xrightarrow{g_\lambda}\sh{H}_\lambda$ una sucesión de homomorfismos de Módulos sobre $\sh{O}_{\mathrm{S}_\lambda}$ cuasi-coherentes, tales que $\operatorname{Im} f_\lambda$ y $\operatorname{Ker} g_\lambda$ son de tipo finito.
Para que la sucesión correspondiente $\sh{F}\xrightarrow{f}\sh{G}\xrightarrow{g}\sh{H}$ sea
\oldpage[IV-3]{24}
exacta, es necesario y suficiente que exista $\mu\geqslant\lambda$ tal que la sucesión $\sh{F}_\mu\xrightarrow{f_\mu}\sh{G}_\mu\xrightarrow{g_\mu}\sh{H}_\mu$ lo sea (en cuyo caso lo mismo vale para $\nu\geqslant\mu$).
  \item Sea $f_\lambda:\sh{F}_\lambda\to\sh{G}_\lambda$ un homomorfismo de Módulos sobre $\sh{O}_{\mathrm{S}_\lambda}$ cuasi-coherentes, y supongamos que $\operatorname{Ker} f_\lambda$ es de tipo finito.
Para que $f=u_\lambda^*(f_\lambda):\sh{F}\to\sh{G}$ sea inyectivo, es necesario y suficiente que exista $\mu\geqslant\lambda$ tal que $f_\mu=u_{\lambda\mu}^*(f_\lambda):\sh{F}_\mu\to\sh{G}_\mu$ lo sea.
\end{enumerate}
\end{corollary}

\begin{proof}
(i) Teniendo en cuenta \sref{IV.8.3.8},~(ii), observemos que, por planitud, $\operatorname{Im} f$ y $\operatorname{Ker} g$ (resp. $\operatorname{Im} f_\mu$ y $\operatorname{Ker} g_\mu$ para $\mu\geqslant\lambda$) son las imágenes inversas de $\operatorname{Im} f_\lambda$ y $\operatorname{Ker} g_\lambda$ \sref[0\textsubscript{I}]{0.6.7.2}. Supongamos que la sucesión $\sh{F}\xrightarrow{f}\sh{G}\xrightarrow{g}\sh{H}$ es exacta. Como $\operatorname{Im} f_\lambda$ es de tipo finito, se deduce de \sref{IV.8.5.2},~(i), que existe $\mu\geqslant\lambda$ tal que la composición $\sh{F}_\mu\xrightarrow{f_\mu}\sh{G}_\mu\xrightarrow{g_\mu}\sh{H}_\mu$ es nula. Cambiando la notación, podemos suponer ya que $g_\lambda\circ f_\lambda=0$. Como $\operatorname{Ker} g_\lambda$ es de tipo finito, se deduce entonces de \sref{IV.8.5.7} que existe $\mu\geqslant\lambda$ tal que el homomorfismo $\sh{F}_\mu\to\operatorname{Ker} g_\mu$ es sobreyectivo, lo que demuestra (i).

(ii) Este es el caso particular de (i) aplicado a la sucesión $0\to\sh{F}_\lambda\to\sh{G}_\lambda$.
\end{proof}

\begin{lemma}[8.5.9]
\label{IV.8.5.9}
Supongamos que $\mathrm{S}_0$ es cuasi-compacto y que $\sh{F}_\lambda$ es cuasi-coherente y de tipo finito. Sean $\sh{G}'_\lambda$ y $\sh{G}''_\lambda$ dos cocientes cuasi-coherentes de $\sh{F}_\lambda$, suponiendo además que $\sh{G}'_\lambda$ es de presentación finita. Para que $\sh{G}''$ sea un cociente de $\sh{G}'$, es necesario y suficiente que exista $\mu\geqslant\lambda$ tal que $\sh{G}''_\mu$ sea un cociente de $\sh{G}'_\mu$.
\end{lemma}

\begin{proof}
Por hipótesis existen dos homomorfismos sobreyectivos $p'_\lambda:\sh{F}_\lambda\to\sh{G}'_\lambda$ y $p''_\lambda:\sh{F}_\lambda\to\sh{G}''_\lambda$. Por la exactitud por la derecha de $u_\lambda^*$ y $u_{\lambda\mu}^*$, los homomorfismos $p'=u_\lambda^*(p'_\lambda)$, $p''=u_\lambda^*(p''_\lambda)$, $p'_\mu=u_{\lambda\mu}^*(p'_\lambda)$ y $p''_\mu=u_{\lambda\mu}^*(p''_\lambda)$ son también sobreyectivos. Además, si existe un homomorfismo $f:\sh{G}'\to\sh{G}''$ (resp. $f_\mu:\sh{G}'_\mu\to\sh{G}''_\mu$) tal que $p''=f\circ p'$ (resp. $p''_\mu=f_\mu\circ p'_\mu$), dicho homomorfismo es necesariamente único; esto muestra que la cuestión es local en $\mathrm{S}_\lambda$ y, en consecuencia, como $\mathrm{S}_0$ es cuasi-compacto y $\mathbf{L}$ es filtrado, podemos suponer que $\mathrm{S}_\lambda$ es afín y por tanto cuasi-separado. La condición del enunciado es claramente suficiente. Recíprocamente, como $\sh{G}'_\lambda$ es de presentación finita y $\mathrm{S}_\lambda$ es cuasi-compacto y cuasi-separado, se deduce de \sref{IV.8.5.2},~(i), que si existe $f:\sh{G}'\to\sh{G}''$ tal que $p''=f\circ p'$, entonces existen $\mu\geqslant\lambda$ y $f_\mu:\sh{G}'_\mu\to\sh{G}''_\mu$ tales que $f=u_\mu^*(f_\mu)$ y $p''_\mu=f_\mu\circ p'_\mu$, lo que demuestra el lema.
\end{proof}

\begin{env}[8.5.10]
\label{IV.8.5.10}
Para el resto de esta sección, para todo Módulo cuasi-coherente $\sh{F}$ sobre un preesquema, denotemos por $\mathfrak{Q}(\sh{F})$ el conjunto de los Módulos cociente de $\sh{F}$ que son de presentación finita. Si $\sh{F}_\lambda$ es cuasi-coherente y $\sh{G}_\lambda\in\mathfrak{Q}(\sh{F}_\lambda)$, se deduce de la exactitud por la derecha de $u_{\lambda\mu}^*$ y $u_\lambda^*$ que $\sh{G}_\mu\in\mathfrak{Q}(\sh{F}_\mu)$ para $\mu\geqslant\lambda$ y $\sh{G}\in\mathfrak{Q}(\sh{F})$. Es claro que $(\mathfrak{Q}(\sh{F}_\lambda),u_{\lambda\mu}^*)$ es un sistema inductivo de conjuntos y que los $u_\lambda^*:\mathfrak{Q}(\sh{F}_\lambda)\to\mathfrak{Q}(\sh{F})$ forman un sistema inductivo de aplicaciones; por tanto, pasando al límite inductivo, existe una aplicación canónica
\[
u_{\mathfrak{Q}}:\varinjlim\mathfrak{Q}(\sh{F}_\lambda)\longrightarrow\mathfrak{Q}(\sh{F}).
\tag{8.5.10.1}\label{IV.8.5.10.1}
\]
Además, si $(\sh{F}'_\lambda)$ es una segunda familia de Módulos sobre $\sh{O}_{\mathrm{S}_\lambda}$ cuasi-coherentes y, para cada $\lambda$, $\sh{F}'_\lambda$ es un cociente de $\sh{F}_\lambda$, entonces $\sh{F}'$ es un cociente de $\sh{F}$, y se tiene un diagrama conmutativo

\oldpage[IV-3]{25}

\[
\begin{tikzcd}
\varinjlim\mathfrak{Q}(\sh{F}_\lambda) \ar[r,"u_{\mathfrak{Q}}"] \ar[d] & \mathfrak{Q}(\sh{F}) \ar[d] \\
\varinjlim\mathfrak{Q}(\sh{F}'_\lambda) \ar[r,"u_{\mathfrak{Q}}"'] & \mathfrak{Q}(\sh{F}')
\end{tikzcd}
\tag{8.5.10.2}\label{IV.8.5.10.2}
\]
\end{env}

\begin{proposition}[8.5.11]
\label{IV.8.5.11}
Supongamos que $\mathrm{S}_0$ es cuasi-compacto (resp. cuasi-compacto y cuasi-separado). Supongamos que $\sh{F}_\lambda$ es cuasi-coherente y de tipo finito (resp. de presentación finita) para cada $\lambda$; entonces la aplicación canónica \eqref{IV.8.5.10.1} es inyectiva (resp. biyectiva).
\end{proposition}

\begin{proof}
La primera afirmación se deduce con mayor precisión del lema \sref{IV.8.5.9}. Para la segunda, consideremos un Módulo cociente $\sh{O}_{\mathrm{S}}$-cuasi-coherente $\sh{G}$ de $\sh{F}$ que sea de presentación finita. Se deduce de \sref{IV.8.5.2},~(ii), que existen $\lambda'\in\mathbf{L}$ y un Módulo sobre $\sh{O}_{\mathrm{S}_{\lambda'}}$ cuasi-coherente y de presentación finita $\sh{G}_{\lambda'}$ tales que $\sh{G}=u_{\lambda'}^*(\sh{G}_{\lambda'})$; como $\mathbf{L}$ es filtrado, podemos reemplazar $\lambda$ y $\lambda'$ por un índice que domine a ambos y suponer que $\lambda'=\lambda$. Consideremos entonces el homomorfismo canónico $p:\sh{F}\to\sh{G}$. Se deduce de \sref{IV.8.5.2},~(i), que existen $\mu\geqslant\lambda$ y un homomorfismo $p_\mu:\sh{F}_\mu\to\sh{G}_\mu$ tales que $p=u_\mu^*(p_\mu)$. En virtud de \sref{IV.8.5.7}, podemos suponer además que $\mu$ es suficientemente grande para que $p_\mu$ sea sobreyectivo, lo que concluye la demostración.
\end{proof}

\subsection{Subpreesquemas de presentación finita de un límite proyectivo de preesquemas}
\label{subsection:IV.8.6}

\begin{env}[8.6.1]
\label{IV.8.6.1}
Dado un preesquema $Y$, denotamos en esta sección por $\mathfrak{Spr}(Y)$ el conjunto ordenado
\sref[I]{I.4.1.10} de los subpreesquemas de $Y$ que son \emph{de presentación finita}
\sref{IV.1.6.1}, y por $\mathfrak{Spro}(Y)$ (resp.\ $\mathfrak{Sprf}(Y)$) el subconjunto de $\mathfrak{Spr}(Y)$ formado por los subpreesquemas inducidos en abiertos (resp.\ en subpreesquemas cerrados) de $Y$, de presentación finita sobre $Y$; equivale a decir que un subpreesquema $Z$ de $Y$ pertenece a $\mathfrak{Spro}(Y)$ (resp.\ $\mathfrak{Sprf}(Y)$) o bien que está inducido en un abierto y que el espacio subyacente es \emph{retrocompacto} en $Y$ (resp.\ que es cerrado y que el ideal $\sh{J}$ de $\sh{O}_Y$ que lo define es \emph{de tipo finito}, lo que también significa que $j_*(\sh{O}_Z)\in\mathfrak{Q}(\sh{O}_Y)$, donde $j:Z\to Y$ es la inyección canónica) \sref{IV.1.6.1} y \sref{IV.1.4.5}.
\end{env}

\begin{env}[8.6.2]
\label{IV.8.6.2}
Conservamos siempre las notaciones de \sref{IV.8.2.2} y suponemos que $S_0$ es uno de los $S_\lambda$.
Dado un subpreesquema $Y_\lambda$ de $S_\lambda$, entonces $Y_\mu=u_{\lambda\mu}^{-1}(Y_\lambda)$ (resp.\ $Y=u_\lambda^{-1}(Y_\lambda)$) es un subpreesquema de $S_\mu$ para $\mu\geqslant\lambda$ (resp.\ de $S$). Es inducido en un abierto (resp.\ es cerrado) si $Y_\lambda$ lo es \sref[I]{I.4.3.2}, y es de presentación finita sobre $S_\mu$ (resp.\ $S$) si $Y_\lambda$ es de presentación finita sobre $S_\lambda$ \sref{IV.1.6.2},~(iii). Se sigue que $(\mathfrak{Spr}(S_\lambda),u_{\lambda\mu}^{-1})$ (resp.\ $(\mathfrak{Spro}(S_\lambda),u_{\lambda\mu}^{-1})$, $(\mathfrak{Sprf}(S_\lambda),u_{\lambda\mu}^{-1})$) es un \emph{sistema inductivo} de conjuntos y que las aplicaciones $u_\lambda^{-1}:\mathfrak{Spr}(S_\lambda)\to\mathfrak{Spr}(S)$ (resp.\ $\mathfrak{Spro}(S_\lambda)\to\mathfrak{Spro}(S)$, $\mathfrak{Sprf}(S_\lambda)\to\mathfrak{Sprf}(S)$) forman un sistema inductivo de aplicaciones; de donde, pasando al límite inductivo, se obtienen aplicaciones canónicas
\end{env}

\begin{align}
  \varinjlim \mathfrak{Spr}(S_\lambda) &\to \mathfrak{Spr}(S) \tag{8.6.2.1}\label{IV.8.6.2.1}\\
  \varinjlim \mathfrak{Spro}(S_\lambda) &\to \mathfrak{Spro}(S) \tag{8.6.2.2}\label{IV.8.6.2.2}\\
  \varinjlim \mathfrak{Sprf}(S_\lambda) &\to \mathfrak{Sprf}(S). \tag{8.6.2.3}\label{IV.8.6.2.3}
\end{align}

\oldpage[IV-3]{26}

Recordemos de \sref[I]{I.4.1.10} que $\mathfrak{Spr}(X)$, para todo preesquema $X$, es un conjunto ordenado por la relación «$Z$ es un subpreesquema del subpreesquema $Y$», que escribimos $Z\leqslant Y$.
Las aplicaciones $u_{\lambda\mu}^{-1}$ y $u_\lambda^{-1}$ son \emph{crecientes} para las relaciones de orden correspondientes en $\mathfrak{Spr}(S_\lambda)$, $\mathfrak{Spr}(S_\mu)$ y $\mathfrak{Spr}(S)$.
Además, escribiendo $\eta\leqslant\eta'$ para dos elementos de este conjunto siempre que existan $\lambda$ y dos elementos $Y_\lambda$, $Y'_\lambda$ de $\mathfrak{Spr}(S_\lambda)$, cuyas imágenes canónicas sean $\eta$ y $\eta'$, y tales que $Y_\lambda\leqslant Y'_\lambda$, se verifica fácilmente que esto no depende de la elección de representantes y que, en efecto, se obtiene una \emph{relación de orden}.
El hecho de que las $u_\lambda^{-1}$ sean crecientes implica inmediatamente que la aplicación canónica \eqref{IV.8.6.2.1} es \emph{creciente}; evidentemente lo mismo vale para \eqref{IV.8.6.2.2} y \eqref{IV.8.6.2.3}.

\begin{proposition}[8.6.3]
\label{IV.8.6.3}
Supongamos que $S_0$ es cuasi-compacto (resp.\ cuasi-compacto y cuasi-separado).
Entonces las aplicaciones \eqref{IV.8.6.2.1}, \eqref{IV.8.6.2.2} y \eqref{IV.8.6.2.3} son inyectivas (resp.\ biyectivas).
\end{proposition}

\begin{proof}
Teniendo en cuenta las observaciones de \sref{IV.8.6.1}, las afirmaciones relativas a \eqref{IV.8.6.2.3} se deducen de \sref{IV.8.5.11} aplicado a $\sh{F}_\lambda=\sh{O}_{S_\lambda}$; asimismo, las afirmaciones relativas a \eqref{IV.8.6.2.2} son casos particulares de \sref{IV.8.3.5} y \sref{IV.8.3.11}, teniendo en cuenta que los $S_\lambda$ y $S$ son cuasi-compactos.
Queda considerar la aplicación \eqref{IV.8.6.2.1}.
Demostremos primero que es sobreyectiva cuando $S_0$ es cuasi-compacto y cuasi-separado.
Sea $Z$ un subpreesquema de $S$, de presentación finita sobre $S$; como $S$ es cuasi-compacto, también lo es $Z$, y por tanto existe un abierto cuasi-compacto $U$ de $S$ tal que $Z$ es un subpreesquema de $U$, cerrado en $U$, de presentación finita sobre $U$.
Existen entonces un índice $\lambda$ y un abierto cuasi-compacto $U_\lambda$ de $S_\lambda$ tales que $U=u_\lambda^{-1}(U_\lambda)$ \sref{IV.8.2.11}; como $S_\lambda$ es cuasi-separado, $U_\lambda$ también lo es \sref[I]{I.1.2.7}, y por tanto podemos restringirnos al caso $U=S$; pero en este caso nos reducimos al hecho de que \eqref{IV.8.6.2.3} es sobreyectiva.

Finalmente, para ver que \eqref{IV.8.6.2.1} es inyectiva cuando $S_0$ es cuasi-compacto, bastará demostrar el resultado más preciso siguiente:
\end{proof}

\begin{corollary}[8.6.3.1]
\label{IV.8.6.3.1}
Supongamos que $S_0$ es cuasi-compacto y sean $Z'_\lambda$, $Z''_\lambda$ dos subpreesquemas de $S_\lambda$, de presentación finita.
Para que $Z'=u_\lambda^{-1}(Z'_\lambda)$ esté dominado por $Z''=u_\lambda^{-1}(Z''_\lambda)$ \sref[I]{I.4.1.10}, es necesario y suficiente que exista $\mu\geqslant\lambda$ tal que $Z'_\mu=u_{\lambda\mu}^{-1}(Z'_\lambda)$ esté dominado por $Z''_\mu=u_{\lambda\mu}^{-1}(Z''_\lambda)$.
\end{corollary}

\begin{proof}
Es trivial que la condición es suficiente.
Para ver que es necesaria, observemos primero que los conjuntos subyacentes $Z'_\lambda$ y $Z''_\lambda$ son localmente constructibles en $S_\lambda$ por hipótesis \sref[I]{I.1.8.4}; por tanto, la hipótesis implica que existe un $\nu\geqslant\lambda$ tal que $Z'_\nu\subset Z''_\nu$ \sref{IV.8.3.5}; reemplazando $\lambda$ por $\nu$, podemos suponer ya que $Z'_\lambda\subset Z''_\lambda$.
Además, por hipótesis \sref[I]{I.1.6.1}, los subespacios $Z'_\lambda$ y $Z''_\lambda$ de $S_\lambda$ son cuasi-compactos; para cada punto $x\in Z'_\lambda$, existe un entorno abierto cuasi-compacto $V(x)$ en $S_\lambda$ tal que $V(x)\cap Z'_\lambda$ y $V(x)\cap Z''_\lambda$ son cerrados en $V(x)$.
Recubriendo $S_\lambda$ por un número finito de tales entornos, vemos que existe un abierto cuasi-compacto $U_\lambda$ de $S_\lambda$ que contiene a $Z'_\lambda$ y tal que $Z'_\lambda$ y $Z''_\lambda\cap U_\lambda$ son cerrados en $U_\lambda$.
Si denotamos por $Y_\lambda$ el subpreesquema de $S_\lambda$ inducido por $Z''_\lambda$ en $U_\lambda\cap Z''_\lambda$, es claro que, con las notaciones usuales, $Y_\mu$ (resp.\ $Y$) es inducido por $Z''_\mu$ en $U_\mu\cap Z''_\mu$ para $\mu\geqslant\lambda$ (resp.\ por $Z''$ en $U\cap Z''$); además, $Z'$ está dominado por $Y$ \sref[I]{I.4.4.1}, y como basta demostrar que $Z'_\mu$ está dominado por $Y_\mu$ para $\mu$
\oldpage[IV-3]{27}
 suficientemente grande, vemos finalmente que (reemplazando $S_\lambda$ por $U_\lambda$) nos reducimos al caso en que $Z'_\lambda$ y $Z''_\lambda$ son subpreesquemas cerrados de $S_\lambda$.
Pero esto ya se ha demostrado, pues \eqref{IV.8.6.2.3} es creciente e inyectiva.
\end{proof}

\begin{corollary}[8.6.4]
\label{IV.8.6.4}
Supongamos que $S_0$ es cuasi-compacto y sea $Z_\lambda$ un subpreesquema de $S_\lambda$, de presentación finita sobre $S_\lambda$.
Para que $Z=u_\lambda^{-1}(Z_\lambda)$ sea el subpreesquema inducido en un abierto de $S$ (resp.\ un subpreesquema cerrado de $S$), es necesario y suficiente que exista $\mu\geqslant\lambda$ tal que $Z_\mu=u_{\lambda\mu}^{-1}(Z_\lambda)$ sea el subpreesquema inducido en un abierto de $S_\mu$ (resp.\ un subpreesquema cerrado de $S_\mu$).
\end{corollary}

\begin{proof}
Sea $(U_\lambda^{(i)})$ un recubrimiento abierto afín finito de $S_\lambda$, y pongamos $U^{(i)}_\mu=u_{\lambda\mu}^{-1}(U_\lambda^{(i)})$ para $\mu\geqslant\lambda$ y $U^{(i)}=u_\lambda^{-1}(U_\lambda^{(i)})$.
Si $Z$ es abierto (resp.\ cerrado) en $S$, $Z\cap U^{(i)}$ es abierto (resp.\ cerrado) en $U^{(i)}$; recíprocamente, si cada $Z_\mu\cap U^{(i)}_\mu$ es abierto (resp.\ cerrado) en $U^{(i)}_\mu$, entonces $Z_\mu$ es abierto (resp.\ cerrado) en $S_\mu$.
Como $L$ es filtrado, basta demostrar el corolario cuando $S_\lambda$ es afín, y por tanto cuasi-separado; el resultado se deduce del hecho de que las aplicaciones \eqref{IV.8.6.2.1}, \eqref{IV.8.6.2.2} y \eqref{IV.8.6.2.3} son biyectivas.
\end{proof}

\subsection{Criterios para que un límite proyectivo de preesquemas sea un preesquema reducido (resp.\ íntegro)}
\label{subsection:IV.8.7}

Conservamos las hipótesis y notaciones de \sref{IV.8.2.2} y suponemos siempre que $S_0$ es uno de los $S_\lambda$.

\begin{proposition}[8.7.1]
\label{IV.8.7.1}
Supongamos que $S$ no es reducido.
Entonces existe $\lambda_0$ tal que, para $\lambda\geqslant\lambda_0$, $S_\lambda$ no es reducido.
\end{proposition}

\begin{proof}
Como la cuestión es local en $S_0$, podemos suponer que $S_0=\operatorname{Spec}(A_0)$ es afín, de donde $S=\operatorname{Spec}(A)$, donde $A=\varinjlim A_\lambda$ es el límite inductivo de un sistema inductivo de $A_0$-álgebras $(A_\lambda)$.
Sabemos entonces por \sref{IV.5.13.2} que el nilradical de $A$ es el límite inductivo de los nilradicales de los $A_\lambda$; si no es nulo, uno de los $A_\lambda$ contiene por tanto un elemento nilpotente $a_\lambda\neq 0$ cuya imagen en $A$ es nilpotente y no nula, y en consecuencia la imagen de $a_\lambda$ en $A_\mu$ es nilpotente y no nula para $\mu\geqslant\lambda$.
\end{proof}

\begin{proposition}[8.7.2]
\label{IV.8.7.2}
Supongamos que se satisface una de las hipótesis siguientes:
\begin{enumerate}[label=\emph{\alph*)}]
\item $S_0$ es cuasi-compacto, el nilradical $\sh{N}_0$ de $\sh{O}_{S_0}$ es un ideal de tipo finito (lo que se satisface, por ejemplo, cuando $S_0$ es noetheriano), y los morfismos $u_{\lambda\mu}:S_\mu\to S_\lambda$ son inmersiones abiertas.
\item Los morfismos $u_{\lambda\mu}:S_\mu\to S_\lambda$ son fielmente planos.
\end{enumerate}
Bajo estas condiciones, para que $S$ sea reducido, es necesario y suficiente que exista $\lambda_0$ tal que $S_\lambda$ sea reducido para $\lambda\geqslant\lambda_0$.

Además, en el caso b), si $S$ es reducido, todos los $S_\lambda$ son reducidos.
\end{proposition}

\begin{proof}
La última afirmación se deduce de que el morfismo $S\to S_\lambda$ es entonces fielmente plano para todo $\lambda$ \sref{IV.8.3.8}, y de \sref{IV.2.1.13}.
Por otro lado, \sref{IV.8.7.1} demuestra que la condición del enunciado es suficiente (sin hipótesis sobre $S_0$ ni sobre los $u_{\lambda\mu}$).
Queda, pues, ver que la condición es necesaria bajo la hipótesis $a)$; entonces, por \sref{IV.8.2.13}, el espacio subyacente de $S$ se identifica con la intersección de los espacios subyacentes de los $S_\lambda$ (identificándose los $S_\lambda$ con subpreesquemas inducidos en abiertos de $S_0$), y el haz estructural $\sh{O}_S$ se identifica con
\oldpage[IV-3]{28}
el haz inducido en $S$ por cada uno de los $\sh{O}_{S_\lambda}$; en particular, para todo $s\in S$, el anillo local $\sh{O}_s$ es el mismo para $S$ y para todos los $S_\lambda$.
Si $\sh{N}_\lambda$ es el nilradical de $\sh{O}_{S_\lambda}$, el nilradical $\sh{N}$ de $\sh{O}_S$ tiene, por tanto, en cada punto de $S$ la misma fibra $\sh{N}_s$ (el nilradical de $\sh{O}_s$) que $u_\lambda^*(\sh{N}_\lambda)$ (inducido en $S$ por $\sh{N}_\lambda$).
La hipótesis de que $S$ es reducido implica entonces $u_\lambda^*(\sh{N}_\lambda)=0$; como se supone que $\sh{N}_0$ es de tipo finito, lo mismo vale para los $\sh{N}_\lambda$, y la conclusión se deduce de \sref{IV.8.5.7}.
\end{proof}

\begin{corollary}[8.7.3]
\label{IV.8.7.3}
Supongamos que se satisface una de las hipótesis siguientes:
\begin{enumerate}[label=\emph{\alph*)}]
\item $S_0$ es un preesquema noetheriano y los morfismos $u_{\lambda\mu}:S_\mu\to S_\lambda$ son inmersiones abiertas.
\item Los morfismos $u_{\lambda\mu}:S_\mu\to S_\lambda$ son fielmente planos.
\end{enumerate}
Entonces, para que $S$ sea íntegro, es necesario y suficiente que exista $\lambda_0$ tal que $S_\lambda$ sea íntegro para $\lambda\geqslant\lambda_0$.
\end{corollary}

\begin{proof}
Decir que un preesquema es íntegro significa que es a la vez reducido e irreducible; el corolario se deduce, pues, de \sref{IV.8.7.2} y de \sref{IV.8.4.3}.
\end{proof}

\begin{remark}[8.7.4]
\label{IV.8.7.4}
Si no se formula ninguna hipótesis sobre los $u_{\lambda\mu}$, puede ocurrir que $S$ sea íntegro aunque ninguno de los $S_\lambda$ sea reducido ni conexo, como muestra el ejemplo \sref{IV.8.4.4}, tomando el anillo $A$ no reducido.
\end{remark}

\subsection{Preesquemas de presentación finita sobre un límite proyectivo de preesquemas}
\label{subsection:IV.8.8}

\begin{env}[8.8.1]
\label{IV.8.8.1}
Conservando siempre las notaciones e hipótesis de \sref{IV.8.2.2}, supondremos dados en esta sección dos $S_\alpha$-preesquemas $X_\alpha$, $Y_\alpha$, lo que define \sref{IV.8.2.5} dos sistemas proyectivos de preesquemas $(X_\lambda,v_{\lambda\mu})$ e $(Y_\lambda,w_{\lambda\mu})$, poniendo $X_\lambda=X_\alpha\times_{S_\alpha}S_\lambda$, $v_{\lambda\mu}=1_{X_\alpha}\times u_{\lambda\mu}$, $w_{\lambda\mu}=1_{Y_\alpha}\times u_{\lambda\mu}$ (para $\alpha\leqslant\lambda\leqslant\mu$), cuyos límites proyectivos son respectivamente $X=X_\alpha\times_{S_\alpha}S$, $Y=Y_\alpha\times_{S_\alpha}S$, con los morfismos canónicos $v_\lambda:X\to X_\lambda$ y $w_\lambda:Y\to Y_\lambda$.
Para $\alpha\leqslant\lambda\leqslant\mu$, existe una aplicación canónica $e_{\mu\lambda}:\operatorname{Hom}_{S_\lambda}(X_\lambda,Y_\lambda)\to\operatorname{Hom}_{S_\mu}(X_\mu,Y_\mu)$, que a todo $S_\lambda$-morfismo $f_\lambda:X_\lambda\to Y_\lambda$ asocia $f_\mu=f_\lambda\times 1_{S_\mu}:X_\lambda\times_{S_\lambda}S_\mu\to Y_\lambda\times_{S_\lambda}S_\mu$, y es claro que $(\operatorname{Hom}_{S_\lambda}(X_\lambda,Y_\lambda),e_{\mu\lambda})$ es un \emph{sistema inductivo} de conjuntos.
Análogamente, existe una aplicación canónica $e_\lambda:\operatorname{Hom}_{S_\lambda}(X_\lambda,Y_\lambda)\to\operatorname{Hom}_S(X,Y)$ que a $f_\lambda$ asocia $f=f_\lambda\times 1_S:X_\lambda\times_{S_\lambda}S\to Y_\lambda\times_{S_\lambda}S$, y $(e_\lambda)$ es un sistema inductivo de aplicaciones; de ahí, pasando al límite inductivo, una aplicación canónica, functorial en $S_\alpha$, $X_\alpha$ e $Y_\alpha$:
\[
  e:\varinjlim\operatorname{Hom}_{S_\lambda}(X_\lambda,Y_\lambda)\to\operatorname{Hom}_S(X,Y).
  \tag{8.8.1.1}\label{IV.8.8.1.1}
\]
\end{env}

\begin{theorem}[8.8.2]
\label{IV.8.8.2}
\begin{enumerate}[label=(\roman*)]
\item Supongamos que $X_\alpha$ es cuasi-compacto (resp.\ cuasi-compacto y cuasi-separado), y que $Y_\alpha$ es localmente de tipo finito (resp.\ localmente de presentación finita) sobre $S_\alpha$.
Entonces la aplicación \eqref{IV.8.8.1.1} es inyectiva (resp.\ biyectiva).
\item Supongamos que $S_0$ es cuasi-compacto y cuasi-separado.
Para todo preesquema $X$ de presentación finita sobre $S$, existen un $\lambda\in L$, un preesquema $X_\lambda$ de presentación finita sobre $S_\lambda$ y un $S$-isomorfismo $X\xrightarrow{\sim}X_\lambda\times_{S_\lambda}S$.
\end{enumerate}
\end{theorem}

\begin{proof}
\oldpage[IV-3]{29}
(i) Consideremos primero el caso en que $S_0=\operatorname{Spec}(A_0)$, $X_\alpha=\operatorname{Spec}(B_\alpha)$ e $Y_\alpha=\operatorname{Spec}(C_\alpha)$ son afines; entonces los $S_\lambda=\operatorname{Spec}(A_\lambda)$ y $S=\operatorname{Spec}(A)$ son también afines, con $A=\varinjlim A_\lambda$, y las afirmaciones de (i) son equivalentes al

\begin{lemma}[8.8.2.1]
\label{IV.8.8.2.1}
Sea $A_0$ un anillo, $(A_\lambda)_{\lambda\in L}$ un sistema inductivo de $A_0$-álgebras, $A=\varinjlim A_\lambda$; sea $B_\alpha$ una $A_\alpha$-álgebra y $C_\alpha$ una $A_\alpha$-álgebra de tipo finito (resp.\ de presentación finita).
Entonces el homomorfismo canónico
\[
  \varinjlim_\lambda\operatorname{Hom}_{A_\lambda\text{-alg.}}(C_\alpha\otimes_{A_\alpha}A_\lambda,B_\alpha\otimes_{A_\alpha}A_\lambda)\to\operatorname{Hom}_{A\text{-alg.}}(C_\alpha\otimes_{A_\alpha}A,B_\alpha\otimes_{A_\alpha}A)
  \tag{8.8.2.2}\label{IV.8.8.2.2}
\]
es inyectivo (resp.\ biyectivo).
\end{lemma}

\begin{proof}
Sabemos que existen isomorfismos canónicos functoriales
\[
  \operatorname{Hom}_{A_\lambda\text{-alg.}}(C_\alpha\otimes_{A_\alpha}A_\lambda,B_\alpha\otimes_{A_\alpha}A_\lambda)\xrightarrow{\sim}\operatorname{Hom}_{A_\alpha\text{-alg.}}(C_\alpha,B_\alpha\otimes_{A_\alpha}A_\lambda)
\]
\[
  \operatorname{Hom}_{A\text{-alg.}}(C_\alpha\otimes_{A_\alpha}A,B_\alpha\otimes_{A_\alpha}A)\xrightarrow{\sim}\operatorname{Hom}_{A_\alpha\text{-alg.}}(C_\alpha,B_\alpha\otimes_{A_\alpha}A)
\]
en virtud de la propiedad universal del producto tensorial de dos álgebras.
Basta, por tanto, demostrar el
\end{proof}

\begin{lemma}[8.8.2.3]
\label{IV.8.8.2.3}
Sea $E$ un anillo, $G$ una $E$-álgebra de tipo finito (resp.\ de presentación finita), $(F_\lambda)$ un sistema inductivo de $E$-álgebras.
Entonces el homomorfismo canónico
\[
  \varinjlim\operatorname{Hom}_{E\text{-alg.}}(G,F_\lambda)\to\operatorname{Hom}_{E\text{-alg.}}(G,\varinjlim F_\lambda)
\]
que a todo sistema inductivo de homomorfismos $\theta_\lambda:G\to F_\lambda$ de $E$-álgebras asocia su límite inductivo, es inyectivo (resp.\ biyectivo).
\end{lemma}

\begin{proof}
Supongamos primero que la $E$-álgebra $G$ es de tipo finito, y sea $(t_i)_{1\leqslant i\leqslant n}$ un sistema de generadores de esta $E$-álgebra; demostremos que, si $(\theta'_\lambda)$, $(\theta''_\lambda)$ son dos sistemas inductivos de homomorfismos $G\to F_\lambda$, con $\varinjlim\theta'_\lambda=\varinjlim\theta''_\lambda$, existe $\mu\geqslant\lambda$ tal que $\theta'_\mu=\theta''_\mu$.
En efecto, si $\varphi_{\mu\lambda}:F_\lambda\to F_\mu$ y $\varphi_\lambda:F_\lambda\to F=\varinjlim F_\lambda$ son los homomorfismos del sistema inductivo $(F_\lambda)$, por hipótesis, para cada índice $i$, existe un índice $\lambda_i$ tal que $\varphi_{\lambda_i}(\theta'_{\lambda_i}(t_i))=\varphi_{\lambda_i}(\theta''_{\lambda_i}(t_i))$, y se puede suponer que todos los $\lambda_i$ son iguales a un único $\lambda$; se deduce del mismo modo la existencia de $\mu\geqslant\lambda$ tal que $\varphi_{\mu\lambda}(\theta'_\lambda(t_i))=\varphi_{\mu\lambda}(\theta''_\lambda(t_i))$ para $1\leqslant i\leqslant n$, es decir, $\theta'_\mu(t_i)=\theta''_\mu(t_i)$ para $1\leqslant i\leqslant n$, de donde $\theta'_\mu=\theta''_\mu$.

Supongamos, en segundo lugar, que $G$ es de presentación finita, de modo que se tiene
\[
  G=E[T_1,\ldots,T_n]/\mathfrak{J},
\]
donde $\mathfrak{J}$ es un ideal de tipo finito, siendo $t_i$ la clase de $T_i$ módulo $\mathfrak{J}$.
Sea $(P_j)_{1\leqslant j\leqslant m}$ un sistema de generadores de $\mathfrak{J}$.
Sea dado un homomorfismo de $E$-álgebras $\theta:G\to F$; pongamos $b_i=\theta(t_i)$; por definición, se tiene entonces $P_j(b_1,b_2,\ldots,b_n)=\theta(P_j(t_1,\ldots,t_n))=0$ para $1\leqslant j\leqslant m$.
Ahora bien, existen $\lambda$ y elementos $x_1,\ldots,x_n$ de $F_\lambda$ tales que $b_i=\varphi_\lambda(x_i)$ para $1\leqslant i\leqslant n$; se tiene, pues, $\varphi_\lambda(P_j(x_1,\ldots,x_n))=P_j(b_1,\ldots,b_n)=0$, y en consecuencia existe $\mu\geqslant\lambda$ tal que $\varphi_{\mu\lambda}(P_j(x_1,\ldots,x_n))=P_j(\varphi_{\mu\lambda}(x_1),\ldots,\varphi_{\mu\lambda}(x_n))=0$ para $1\leqslant j\leqslant m$; se concluye que existe un homomorfismo de $E$-álgebras $\theta_\mu:G\to F_\mu$ tal que $\theta_\mu(t_i)=\varphi_{\mu\lambda}(x_i)$
\oldpage[IV-3]{30}
para $1\leqslant i\leqslant n$; se deduce, por tanto, para todo $\nu\geqslant\mu$, un homomorfismo de $E$-álgebras $\theta_\nu=\varphi_{\nu\mu}\circ\theta_\mu$ de $G$ en $F_\nu$, y es claro que $\theta$ es el límite inductivo de este sistema inductivo de homomorfismos, lo que concluye la demostración del lema.
\end{proof}

Pasemos ahora al caso en que $X_\alpha$ es cuasi-compacto y $Y_\alpha$ localmente de tipo finito sobre $S_\alpha$.
Pongamos $Z_\alpha=X_\alpha\times_{S_\alpha}Y_\alpha$ e introduzcamos el sistema proyectivo correspondiente a $Z_\lambda=X_\lambda\times_{S_\lambda}Y_\lambda$ y su límite $Z=X\times_S Y$; las biyecciones canónicas \sref[I]{I.3.3.14} dan diagramas conmutativos
\[
\xymatrix{
  \operatorname{Hom}_{S_\lambda}(X_\lambda,Y_\lambda) \ar[r]^-{e_{\mu\lambda}} \ar[d]^{\wr} & \operatorname{Hom}_{S_\mu}(X_\mu,Y_\mu) \ar[r]^-{e_\mu} \ar[d]^{\wr} & \operatorname{Hom}_S(X,Y) \ar[d]^{\wr}\\
  \operatorname{Hom}_{X_\lambda}(X_\lambda,Z_\lambda) \ar[r]_-{e_{\mu\lambda}} & \operatorname{Hom}_{X_\mu}(X_\mu,Z_\mu) \ar[r]_-{e_\mu} & \operatorname{Hom}_X(X,Z)
}
\]
y, en consecuencia, nos reducimos a demostrar que \eqref{IV.8.8.1.1} es inyectiva en el caso particular en que $S_\alpha=X_\alpha$ (teniendo en cuenta \sref[I]{I.1.3.4}).
Además, como $X_\alpha$ es cuasi-compacto, y por tanto unión finita de abiertos afines, podemos suponer $X_\alpha$ afín (puesto que $L$ es filtrado). Sean, pues, dos $X_\alpha$-morfismos $f'_\alpha:X_\alpha\to Y_\alpha$, $f''_\alpha:X_\alpha\to Y_\alpha$ tales que $f'$ y $f''$ son $X$-morfismos iguales de $X$ en $Y$; debemos demostrar que $f'_\mu=f''_\mu$ para $\mu\geqslant\alpha$ suficientemente grande. Como $X_\alpha$ es cuasi-compacto, también lo es $f'_\alpha(X_\alpha)\cup f''_\alpha(X_\alpha)$ y, como $Y_\alpha$ es localmente de tipo finito sobre $X_\alpha$, esta unión está contenida en una unión finita de abiertos afines $V_{i\alpha}$ de $Y_\alpha$, de tipo finito sobre $X_\alpha$. Pongamos $U'_{i\alpha}=f_\alpha^{'-1}(V_{i\alpha})$, $U''_{i\alpha}=f_\alpha^{''-1}(V_{i\alpha})$, $U_{i\alpha}=U'_{i\alpha}\cap U''_{i\alpha}$ y $U_\alpha=\bigcup_iU_{i\alpha}$. La hipótesis $f'=f''$ implica $v_\alpha^{-1}(U'_{i\alpha})=v_\alpha^{-1}(U''_{i\alpha})$, pues estos dos conjuntos son respectivamente iguales a $f^{'-1}(w_\alpha^{-1}(V_{i\alpha}))$ y $f^{''-1}(w_\alpha^{-1}(V_{i\alpha}))$. Como los $V_{i\alpha}$ recubren $Y_\alpha$, se tiene $v_\alpha^{-1}(U_\alpha)=f^{'-1}(Y)=X$. Puesto que $X_\alpha$ es cuasi-compacto y todo abierto de $X_\lambda$ es indconstructible, se deduce de \sref{IV.8.3.4} que existe un índice $\mu\geqslant\alpha$ tal que los $U_{i\mu}=v_{\alpha\mu}^{-1}(U_{i\alpha})$ forman un recubrimiento de $X_\mu$. Reemplazando $\alpha$ por $\mu$, podemos suponer, por tanto, que los $U_{i\alpha}$ recubren $X_\alpha$; esto implica que, para todo $x\in X_\alpha$, existe un entorno abierto afín $W(x)$ contenido en uno de los $U_{i\alpha}$, es decir, tal que las restricciones de $f'_\alpha$ y $f''_\alpha$ a $W(x)$ envían $W(x)$ al mismo abierto afín $V_{i\alpha}$. Como $X_\alpha$ es cuasi-compacto, está recubierto por un número finito de abiertos afines $W(x_k)$; en virtud del lema~\sref{IV.8.8.2.1} y de que $L$ es filtrado, existe $\lambda\geqslant\alpha$ tal que las restricciones de $f'_\lambda$ y $f''_\lambda$ a cada uno de los abiertos $v_{\alpha\lambda}^{-1}(W(x_k))$ son iguales, de donde $f'_\lambda=f''_\lambda$.

Supongamos ahora que $X_\alpha$ es cuasi-compacto y cuasi-separado y que $Y_\alpha$ es localmente de presentación finita sobre $S_\alpha$, y demostremos que \eqref{IV.8.8.1.1} es sobreyectiva.
Sea, pues, un $X$-morfismo $f:X\to Y$.
Como $X$ es cuasi-compacto, también lo es
\oldpage[IV-3]{31}
$f(X)$ y, en consecuencia, existe un abierto cuasi-compacto $Y'$ de $Y$ que contiene a $f(X)$; existen entonces $\lambda\geqslant\alpha$ y un abierto cuasi-compacto $Y'_\lambda$ de $Y_\lambda$ tales que $Y'=w_\lambda^{-1}(Y'_\lambda)$ \sref{IV.8.2.11}.
Reemplazando, si es necesario, $\alpha$ por $\lambda$ e $Y_\lambda$ por $Y'_\lambda$, podemos, pues, restringirnos al caso en que $Y_\alpha$ es \emph{cuasi-compacto}, y por tanto también $Y$ y los $Y_\lambda$.
Como $Y$ es cuasi-compacto, es unión finita de abiertos afines $V_i$, y en consecuencia $X$ es la unión de los abiertos $f^{-1}(V_i)$. Puesto que todo punto de $X$ posee un entorno abierto cuasi-compacto contenido en uno de los $f^{-1}(V_i)$ y $X$ es cuasi-compacto, podemos, teniendo en cuenta \sref{IV.8.2.11} y reemplazando $\alpha$, si es necesario, por un índice $\lambda\geqslant\alpha$, suponer que $Y$ es unión finita de abiertos $V_i=w_\alpha^{-1}(V_{i\alpha})$, donde los $V_{i\alpha}$ son abiertos afines de $Y_\alpha$; en consecuencia, $X$ es la unión de los abiertos $f^{-1}(V_i)$. Como todo punto de $X$ posee un entorno abierto cuasi-compacto contenido en uno de los $f^{-1}(V_i)$ y $X$ es cuasi-compacto, podemos recubrir $X$ por un número finito de tales entornos y, repitiendo algunos de los $V_i$ si es necesario, suponer que existe un recubrimiento $(U_i)$ de $X$ por abiertos cuasi-compactos, con el mismo conjunto de índices que $(V_i)$ y tal que $f(U_i)\subset V_i$ para todo $i$. Además, utilizando \sref{IV.8.2.11} (y reemplazando $\alpha$, si es necesario, por un índice $\lambda\geqslant\alpha$), podemos suponer que $U_i=v_\alpha^{-1}(U_{i\alpha})$, donde los $U_{i\alpha}$ son abiertos cuasi-compactos de $X_\alpha$; utilizando asimismo \sref{IV.8.3.4} como antes, podemos suponer que $(U_{i\alpha})$ es un \emph{recubrimiento} de $X_\alpha$.

Bastará entonces demostrar que existen $\lambda\geqslant\alpha$ y, para cada $i$, un morfismo $f_{i\lambda}:U_{i\lambda}\to V_{i\lambda}$ (donde $U_{i\lambda}=v_{\alpha\lambda}^{-1}(U_{i\alpha})$ y $V_{i\lambda}=w_{\alpha\lambda}^{-1}(V_{i\alpha})$) tal que el morfismo correspondiente $f_i=e_\lambda(f_{i\lambda}):U_i\to V_i$ sea igual a la \emph{restricción} de $f$ a $U_i$.
En efecto, si esto es así, como $X_\lambda$ es cuasi-separado \sref[I]{I.1.2.3}, los $U_{i\lambda}\cap U_{j\lambda}$ son \emph{cuasi-compactos}, y el resultado de unicidad demostrado anteriormente (aplicable puesto que $Y_\lambda$ es localmente de tipo finito sobre $S_\lambda$) prueba que existe $\mu\geqslant\lambda$ tal que $f_{i\mu}$ y $f_{j\mu}$ coinciden en $U_{i\mu}\cap U_{j\mu}$ para todo par $(i,j)$ y definen, en consecuencia, un $S_\mu$-morfismo $f_\mu:X_\mu\to Y_\mu$ tal que $e_\mu(f_\mu)=f$.

Nos reducimos así al caso en que $Y_\alpha$ es \emph{afín} y, como además podemos suponer que las imágenes de los $V_{i\alpha}$ en $S_\alpha$ están contenidas en abiertos afines, también podemos suponer que $S_\alpha$ es afín; sean, pues, $S_\alpha=\operatorname{Spec}(A_\alpha)$, $Y_\alpha=\operatorname{Spec}(C_\alpha)$, con $C_\alpha$ una $A_\alpha$-álgebra de presentación finita, y sean $S=\operatorname{Spec}(A)$, $Y=\operatorname{Spec}(C)$, donde $A=\varinjlim A_\lambda$ y $C=\varinjlim(C_\alpha\otimes_{A_\alpha}A_\lambda)=C_\alpha\otimes_{A_\alpha}A$.
Se tiene entonces
\[
  \operatorname{Hom}_S(X,Y)=\operatorname{Hom}_{A\text{-alg.}}(C,\Gamma(X,\sh{O}_X))=\operatorname{Hom}_{A_\alpha\text{-alg.}}(C_\alpha,\Gamma(X,\sh{O}_X))
\]
\sref[I]{I.2.2.4} y, análogamente,
\[
  \operatorname{Hom}_{S_\lambda}(X_\lambda,Y_\lambda)=\operatorname{Hom}_{A_\lambda\text{-alg.}}(C_\alpha\otimes_{A_\alpha}A_\lambda,\Gamma(X_\lambda,\sh{O}_{X_\lambda}))=\operatorname{Hom}_{A_\alpha\text{-alg.}}(C_\alpha,\Gamma(X_\lambda,\sh{O}_{X_\lambda})).
\]
Pero, como $X_\alpha$ es cuasi-compacto y cuasi-separado, se sabe por \sref{IV.8.5.4} que $\varinjlim\Gamma(X_\lambda,\sh{O}_{X_\lambda})=\Gamma(X,\sh{O}_X)$; puesto que $C_\alpha$ es una $A_\alpha$-álgebra de presentación finita, la biyectividad de \eqref{IV.8.8.1.1} se deduce entonces de \sref{IV.8.8.2.3}.

Antes de pasar a la demostración de (ii), observemos los siguientes corolarios de (i):
\end{proof}

\begin{corollary}[8.8.2.4]
\label{IV.8.8.2.4}
Supongamos que $S_0$ es cuasi-compacto, que $X_\alpha$ es de presentación finita sobre $S_\alpha$, y que $Y_\alpha$ es de tipo finito sobre $S_\alpha$ y cuasi-separado sobre $S_\alpha$ (lo que ocurre, por ejemplo, si $Y_\alpha$ es también
\oldpage[IV-3]{32}
de presentación finita sobre $S_\alpha$).
Sea $f_\alpha:X_\alpha\to Y_\alpha$ un $S_\alpha$-morfismo.
Para que $f:X\to Y$ sea un isomorfismo, es necesario y suficiente que exista $\lambda\geqslant\alpha$ tal que $f_\lambda:X_\lambda\to Y_\lambda$ sea un isomorfismo.
\end{corollary}

\begin{proof}
La condición es evidentemente suficiente. Para demostrar que es necesaria, observemos primero que, como la cuestión es local en $S_0$ (puesto que $L$ es filtrado), siempre podemos suponer $S_0$ afín y, por tanto, cuasi-separado. Por hipótesis existe un $S$-morfismo $g:Y\to X$ tal que $g\circ f=1_X$ y $f\circ g=1_Y$. Como $X_\alpha$ es de presentación finita sobre $S_\alpha$, e $Y_\alpha$ es cuasi-compacto y cuasi-separado (porque $S_\alpha$ es cuasi-compacto y cuasi-separado), se deduce de \sref{IV.8.8.2},~(i) que existen $\mu\geqslant\alpha$ y un $S_\mu$-morfismo $g_\mu:Y_\mu\to X_\mu$ tales que $g=g_\mu\times 1_S$. Se deduce además de \sref{IV.8.8.2},~(i), junto con las relaciones $g\circ f=1_X$ y $f\circ g=1_Y$, que existe $\nu\geqslant\mu$ tal que $g_\nu\circ f_\nu=1_{X_\nu}$ y $f_\nu\circ g_\nu=1_{Y_\nu}$, puesto que $X_\alpha$ e $Y_\alpha$ son de tipo finito sobre $S_\alpha$ y cuasi-compactos. Así, $f_\nu:X_\nu\to Y_\nu$ es un isomorfismo, lo que demuestra el corolario.
\end{proof}

\begin{corollary}[8.8.2.5]
\label{IV.8.8.2.5}
Supongamos que $S_0$ es cuasi-compacto y cuasi-separado, y que $X_\alpha$ e $Y_\alpha$ son de presentación finita sobre $S_\alpha$.
Para que $X$ e $Y$ sean $S$-isomorfos, es necesario y suficiente que exista $\lambda\geqslant\alpha$ tal que $X_\lambda$ e $Y_\lambda$ sean $S_\lambda$-isomorfos. Además, para todo $S$-isomorfismo $f:X\xrightarrow{\sim}Y$, existen $\mu\geqslant\lambda$ y un $S_\mu$-isomorfismo $f_\mu:X_\mu\xrightarrow{\sim}Y_\mu$ tales que $f=f_\mu\times 1_S$.
\end{corollary}

\begin{proof}
La condición es evidentemente suficiente. Recíprocamente, si existe un $S$-isomorfismo $f:X\to Y$, se deduce de \sref{IV.8.8.2},~(i) que $f$ es de la forma $f_\mu\times 1_S$ para algún $\mu\geqslant\alpha$ y algún $S_\mu$-morfismo $f_\mu:X_\mu\to Y_\mu$. Como $f$ es un isomorfismo, se deduce de \sref{IV.8.8.2.4} que existe $\nu\geqslant\mu$ tal que $f_\nu:X_\nu\to Y_\nu$ es un isomorfismo.
\end{proof}

\begin{proof}[Demostración de \sref{IV.8.8.2}, (ii)]
Consideremos primero el caso en que $S_0=\operatorname{Spec}(A_0)$ y $X=\operatorname{Spec}(B)$ son afines; entonces la afirmación~(ii) es equivalente al lema~\sref{IV.1.8.4.2}.

Para demostrar (ii) en el caso general, observemos que $S$ es cuasi-compacto y cuasi-separado. Como $X$ es de presentación finita sobre $S$, y $S$ es afín sobre $S_0$, existe un recubrimiento abierto finito $(U_i)$ de $S_0$ y, si $(W_i)$ denota el recubrimiento abierto afín de $S$ formado por las imágenes inversas de los $U_i$ por $S\to S_0$, un recubrimiento finito $(X_r)$ de $X$ por abiertos afines tal que la imagen de cada $X_r$ por $X\to S$ está contenida en algún $W_{i(r)}$. El anillo $A(X_r)$ es entonces un álgebra de presentación finita sobre $A(W_{i(r)})$ \sref{IV.1.4.6}.

Por el lema~\sref{IV.1.8.4.2} y el hecho de que $L$ es filtrado, existen un índice $\lambda\in L$ y, para todo $r$, un esquema afín $Z_{r\lambda}$ de presentación finita sobre la imagen inversa $W_{i(r),\lambda}$ de $U_{i(r)}$ en $S_\lambda$, junto con un $S$-isomorfismo
\[
 g_r:Z_{r\lambda}\times_{S_\lambda}S\xrightarrow{\sim}X_r.
\]
Sea $Z_{rs}$ la imagen inversa por $g_r$ del subpreesquema inducido en el abierto $X_r\cap X_s$ de $X_r$, que es cuasi-compacto porque $X$ es cuasi-separado, y sea
\[
 g'_{rs}:Z_{rs}\xrightarrow{\sim}X_r\cap X_s
\]
la restricción de $g_r$. Por \sref{IV.8.2.11}, existen $\mu\geqslant\lambda$ y, para todo par $(r,s)$, un abierto cuasi-compacto $Z_{rs\mu}$ de $Z_{r\mu}=v_{\lambda\mu}^{-1}(Z_{r\lambda})$ tal que $Z_{rs}$ es la imagen inversa de $Z_{rs\mu}$. Además, como $S_\mu$ es cuasi-separado y $W_{i(r),\mu}$ es un abierto cuasi-compacto de $S_\mu$, cada $Z_{rs\mu}$ es de presentación finita sobre $S_\mu$ \sref[I]{I.1.6.2}.

Para todo par $(r,s)$, consideremos el isomorfismo
\[
 h_{sr}=(g'_{sr})^{-1}\circ g'_{rs}:Z_{rs}\xrightarrow{\sim}Z_{sr}.
\]
Por \sref{IV.8.8.2.4}, existen $\nu\geqslant\mu$ y, para todo par $(r,s)$, un isomorfismo
\[
 h_{sr\nu}:Z_{rs\nu}\xrightarrow{\sim}Z_{sr\nu}
\]
tal que $h_{sr}=h_{sr\nu}\times 1_S$. Finalmente, para toda terna $(r,s,t)$, denotemos por $h'_{sr}$ la restricción de $h_{sr}$ a

\oldpage[IV-3]{33}

$Z_{rs}\cap Z_{rt}$. Se deduce inmediatamente de las definiciones que $h'_{sr}$ es un isomorfismo de $Z_{rs}\cap Z_{rt}$ sobre $Z_{sr}\cap Z_{st}$ y que
\[
 h'_{ts}\circ h'_{sr}=h'_{tr}.
\]
Por \sref{IV.8.8.2},~(i), existe entonces $\rho\geqslant\nu$ tal que, para toda terna $(r,s,t)$,
\[
 h'_{ts\rho}\circ h'_{sr\rho}=h'_{tr\rho}.
\]
En consecuencia, los isomorfismos $h_{sr\rho}$ satisfacen la condición de pegado \sref[0\textsubscript{I}]{0.4.1.7} y definen, por tanto, un preesquema $X_\rho$ tal que $X$ es isomorfo a $X_\rho\times_{S_\rho}S$. Además, los $Z_{r\rho}$ son de presentación finita sobre $S_\rho$ y, al identificarlos con abiertos de $X_\rho$, sus intersecciones $Z_{r\rho}\cap Z_{s\rho}$, que son isomorfas a los $Z_{rs\rho}$, son cuasi-compactas. Así, $X_\rho$ es de presentación finita sobre $S_\rho$ \sref[I]{I.1.6.3}.
\end{proof}

\begin{env}[8.8.3]
\label{IV.8.8.3}
El contenido esencial de \sref{IV.8.8.2} se expresa también diciendo que, si $S_0$ es cuasi-compacto y cuasi-separado, la categoría de los $S$-preesquemas de presentación finita queda determinada, salvo equivalencia, por los datos de las categorías de los $S_\lambda$-preesquemas de presentación finita, los funtores $\rho_{\mu\lambda}:X_\lambda\mapsto X_\lambda\times_{S_\lambda}S_\mu$ (para $\lambda\leqslant\mu$) entre estas categorías y los isomorfismos de transitividad $\rho_{\nu\mu}\circ\rho_{\mu\lambda}\xrightarrow{\sim}\rho_{\nu\lambda}$.
De manera figurada, puede decirse que el dato de un $S$-preesquema de presentación finita $X$ es ``funtorialmente'' equivalente al dato de un $S_\lambda$-preesquema de presentación finita $X'_\lambda$; si un $S_\mu$-preesquema de presentación finita $X''_\mu$ es tal que $X$ es $S$-isomorfo a $X''_\mu\times_{S_\mu}S$, existe $\nu$ tal que $\lambda\leqslant\nu$, $\mu\leqslant\nu$, y $X'_\lambda\times_{S_\lambda}S_\nu$ y $X''_\mu\times_{S_\mu}S_\nu$ son $S_\nu$-isomorfos.
El hecho de que $L$ sea filtrado implica además que, si se da una \emph{familia} finita $(X^{(i)})_{i\in I}$ de $S$-preesquemas de presentación finita y una familia finita $(f^{(j)})_{j\in J}$ de $S$-morfismos entre estos preesquemas (siendo, por tanto, $f^{(j)}$ de la forma $X^{(\sigma(j))}\to X^{(\tau(j))}$, donde $\sigma$ y $\tau$ son dos aplicaciones de $J$ en $I$), entonces existen un índice $\lambda\in L$, una familia $(X^{(i)}_\lambda)_{i\in I}$ de $S_\lambda$-preesquemas de presentación finita y una familia $(f^{(j)}_\lambda)_{j\in J}$ de $S_\lambda$-morfismos tales que $X^{(i)}$ se identifica con $X^{(i)}_\lambda\times_{S_\lambda}S$ y $f^{(j)}$ con $f^{(j)}_\lambda\times 1_S$ para todos $i$ y $j$; además, si se tiene una relación $f^{(j)}=f^{(k)}$, se puede suponer $\lambda$ elegido de modo que $f^{(j)}_\lambda=f^{(k)}_\lambda$.

Consideremos, en particular, tres $S$-preesquemas de presentación finita $X$, $Y$, $Z$ y dos $S$-morfismos $f:X\to Z$, $g:Y\to Z$, de modo que el producto $X\times_Z Y$ relativo a estos morfismos es también un $S$-preesquema de presentación finita \sref[I]{I.1.6.2}.
Se deduce entonces de lo anterior y de \sref[I]{I.3.3.11} que se puede escribir $X\times_Z Y=(X_\lambda\times_{Z_\lambda}Y_\lambda)\times_{S_\lambda}S$ para un $\lambda$ adecuado, siendo $X_\lambda$, $Y_\lambda$, $Z_\lambda$ $S_\lambda$-preesquemas de presentación finita; puede decirse, por tanto, que la determinación de los $S$-preesquemas de presentación finita mediante los datos de los $S_\lambda$-preesquemas de presentación finita es ``\emph{compatible con los productos fibrados}''.
Por otro lado, hemos visto \sref{IV.8.6.3} que, si $g$ es una inmersión, se puede suponer lo mismo para $g_\lambda$; identificando $Y$ (resp.\ $Y_\lambda$) con un subpreesquema de $Z$ (resp.\ $Z_\lambda$), se ve que se puede escribir, para un $\lambda$ adecuado, $f^{-1}(Y)=f_\lambda^{-1}(Y_\lambda)\times_{S_\lambda}S$ \sref[I]{I.4.4.1};
existe, pues, también ``\emph{compatibilidad con la formación de imágenes inversas de subpreesquemas}''.
Más particularmente, si $f_1$, $f_2$ son dos $S$-morfismos de $X$ en $Y$, se llama \emph{núcleo} de este par de morfismos a la imagen inversa $N$ del subpreesquema diagonal de $Y\times_S Y$ por el $S$-morfismo $(f_1,f_2)_S$; se deduce de lo anterior que se tendrá, para un $\lambda$ adecuado, $N=N_\lambda\times_{S_\lambda}S$, donde $N_\lambda$ es el núcleo del par de $S_\lambda$-morfismos $(f_{1\lambda},f_{2\lambda})$ de $X_\lambda$ en $Y_\lambda$.
Estas observaciones se extienden a los productos finitos y a los ``núcleos'' de sistemas finitos arbitrarios de $S$-morfismos de un $S$-preesquema de presentación finita en
\oldpage[IV-3]{34}
otro; se concluye que, de manera general, la formación de $S$-preesquemas de presentación finita mediante los datos de $S_\lambda$-preesquemas de presentación finita es ``\emph{compatible con los límites proyectivos finitos}'', siendo tal límite, por definición, el núcleo de un sistema finito de morfismos de un $S$-preesquema en un producto de $S$-preesquemas \sref[T]{T.1.8}.

En lo que sigue nos permitiremos efectuar libremente las traslaciones implicadas por las propiedades anteriores (así como por \sref{IV.8.3.11}, \sref{IV.8.5.2} y \sref{IV.8.6.3}), sin esforzarnos siempre por justificarlas como se ha hecho anteriormente.
Por ejemplo, el dato de un ``preesquema en grupos'' $G$ de presentación finita sobre $S$ equivale al dato de un preesquema en grupos $G_\lambda$ de presentación finita sobre algún $S_\lambda$ para un $\lambda$ suficientemente grande; en efecto, escribir la condición de asociatividad para el $S$-morfismo ``ley de composición'' $G\times_S G\to G$ del preesquema en grupos equivale a escribir que el núcleo de dos $S$-morfismos de la forma $G\times_S G\times_S G\to G$ es igual a $G\times_S G\times_S G$ \sref[II]{II.8.3.9}, y las demás condiciones que intervienen en la definición de un preesquema en grupos se interpretan de manera análoga.
\end{env}

\subsection{Primeras aplicaciones a la eliminación de las hipótesis noetherianas}
\label{subsection:IV.8.9}

\begin{proposition}[8.9.1]
\label{IV.8.9.1}
Sea $A$ un anillo y sea $X$ un $A$-preesquema.
\begin{enumerate}[label=(\roman*)]
\item Las condiciones siguientes son equivalentes:
\begin{enumerate}[label=\emph{\alph*)}]
\item $X$ es de presentación finita sobre $A$.
\item Existen un anillo noetheriano $A_0$, un preesquema $X_0$ de tipo finito sobre $A_0$, un homomorfismo de anillos $A_0\to A$ y un $A$-isomorfismo $X_0\otimes_{A_0}A\xrightarrow{\sim}X$.
\item Existen un subanillo $A_0$ de $A$, que es una $\mathbf{Z}$-álgebra de tipo finito, un preesquema $X_0$ de tipo finito sobre $A_0$ y un $A$-isomorfismo $X_0\otimes_{A_0}A\xrightarrow{\sim}X$.
\end{enumerate}
\item Si $\sh{F}$ es un $\sh{O}_X$-Módulo cuasicoherente de presentación finita, se puede suponer elegido el subanillo $A_0$ de modo que exista un $\sh{O}_{X_0}$-Módulo coherente $\sh{F}_0$ tal que $\sh{F}$ sea isomorfo a $\sh{F}_0\otimes_{A_0}A$; $\operatorname{Supp}(\sh{F})$ es construible y cerrado en $X$, y existe un subpreesquema cerrado $Z$ de $X$, de espacio subyacente $\operatorname{Supp}(\sh{F})$, tal que la inmersión canónica $Z\to X$ sea de presentación finita.
\item Si $Y$ es otro $A$-preesquema de presentación finita y $f:X\to Y$ un $A$-morfismo, se puede suponer elegido el subanillo $A_0$ de $A$ de modo que existan un preesquema $Y_0$ de tipo finito sobre $A_0$, un $A$-isomorfismo $Y_0\otimes_{A_0}A\xrightarrow{\sim}Y$ y un $A_0$-morfismo $f_0:X_0\to Y_0$ (necesariamente de tipo finito) tal que $f$ se identifique con $f_0\otimes 1_A$.
\end{enumerate}
\end{proposition}

\begin{proof}
\begin{enumerate}[label=(\roman*)]
\item Como $A$ es el límite inductivo de sus subanillos de tipo finito sobre $\mathbf{Z}$, \emph{a)} entraña \emph{c)} en virtud de \sref{IV.8.8.2},~(ii); \emph{c)} entraña \emph{b)}, pues una $\mathbf{Z}$-álgebra de tipo finito es un anillo noetheriano; por último, si $A_0$ es noetheriano, todo $A_0$-preesquema de tipo finito es de presentación finita \sref[I]{I.1.6.1}, de modo que \emph{b)} entraña \emph{a)} en virtud de \sref[I]{I.1.6.2},~(iii).
\end{enumerate}
La aserción~(ii) se sigue inmediatamente de \sref{IV.8.5.2},~(ii), \sref{IV.8.3.11} y \sref[I]{I.1.8.2}, y la aserción~(iii), de \sref{IV.8.8.2},~(i).
\end{proof}

\begin{env}[8.9.2]
\label{IV.8.9.2}
La proposición \sref{IV.8.9.1} y los resultados de \sref{IV.8.5.2} y \sref{IV.8.8.2} permiten extender en numerosos casos a los morfismos de presentación finita $X\to Y$ los resultados demostrados mediante técnicas noetherianas suponiendo $Y$ \emph{localmente noetheriano}.
\oldpage[IV-3]{35}
Veremos numerosos ejemplos en lo que sigue y nos limitaremos aquí a dar algunos resultados de este tipo.
\end{env}

\begin{proposition}[8.9.3]
\label{IV.8.9.3}
Sea $A$ un anillo y sea $M$ un $A$-módulo de presentación finita.
Entonces todo $A$-endomorfismo sobreyectivo de $M$ es biyectivo.
\end{proposition}

\begin{proof}
En efecto, consideremos $A$ como límite inductivo de sus sub-$\mathbf{Z}$-álgebras de tipo finito.
De \sref{IV.8.5.2},~(ii) se sigue que existen una de estas subálgebras $A_0$ y un $A_0$-módulo $M_0$ de presentación finita tales que $M$ es $A$-isomorfo a $M_0\otimes_{A_0}A$; además, si $f:M\to M$ es un $A$-endomorfismo sobreyectivo, se puede suponer \sref{IV.8.5.2},~(i) que existe un $A_0$-endomorfismo $f_0:M_0\to M_0$ tal que $f=f_0\otimes 1_A$; por último, en virtud de \sref{IV.8.5.7}, se puede suponer que $f_0$ es \emph{sobreyectivo}.
Pero, como $A_0$ es noetheriano y $M_0$ un $A_0$-módulo de tipo finito, $M_0$ es un módulo noetheriano; por tanto (Bourbaki, \emph{Alg.}, cap.~VIII, \textsection\,2, n\textsuperscript{o}\,2, lema~3), $f_0$ es biyectivo, y lo mismo ocurre con $f$ por cambio de base.
\end{proof}

\begin{proposition}[8.9.4]
\label{IV.8.9.4}
(``teorema de planitud genérica''). --- Sea $Y$ un preesquema íntegro, sea $u:X\to Y$ un morfismo de tipo finito y localmente de presentación finita, y sea $\sh{F}$ un $\sh{O}_X$-Módulo cuasicoherente de presentación finita.
Entonces existe un abierto no vacío $U$ de $Y$ tal que $\sh{F}|u^{-1}(U)$ es plano sobre $U$.
\end{proposition}

\begin{proof}
El razonamiento del comienzo de \sref{IV.6.9.1} conduce a demostrar el

\begin{lemma}[8.9.4.1]
\label{IV.8.9.4.1}
Sea $A$ un anillo íntegro, sea $B$ una $A$-álgebra de presentación finita y sea $M$ un $B$-módulo de presentación finita.
Entonces existe $f\neq 0$ en $A$ tal que $M_f$ es un $A_f$-módulo libre.
\end{lemma}

\begin{proof}
En efecto, por \sref{IV.8.9.1} existen una $\mathbf{Z}$-subálgebra $A_0$ de $A$ de tipo finito, una $A_0$-álgebra $B_0$ de tipo finito y un $B_0$-módulo $M_0$ de tipo finito tales que $B$ es $A$-isomorfo a $B_0\otimes_{A_0}A$ y $M$ es $B$-isomorfo a $M_0\otimes_{A_0}A$; por \sref{IV.6.9.2}, existe $f_0\neq 0$ en $A_0$ tal que $(M_0)_{f_0}$ es un $(A_0)_{f_0}$-módulo libre.
Como $M_{f_0}=(M_0)_{f_0}\otimes_{A_0}A$ y $A_{f_0}=(A_0)_{f_0}\otimes_{A_0}A$, $M_{f_0}$ es un $A_{f_0}$-módulo libre.
\end{proof}
\end{proof}

\begin{corollary}[8.9.5]
\label{IV.8.9.5}
Sea $S$ un preesquema cuasicompacto y cuasiseparado, sea $u:X\to S$ un morfismo de presentación finita y sea $\sh{F}$ un $\sh{O}_X$-Módulo cuasicoherente de presentación finita.
Entonces existe una partición $(S_i)_{1\leqslant i\leqslant n}$ de $S$ en un número finito de subconjuntos localmente cerrados de $S$ tal que, para $1\leqslant i\leqslant n$, existe un subpreesquema $S'_i$ de $S$, de espacio subyacente $S_i$, que es de presentación finita sobre $S$, y tal que, poniendo $X_i=X\times_S S'_i$, el $\sh{O}_{X_i}$-Módulo $\sh{F}_i=\sh{F}\otimes_{\sh{O}_X}\sh{O}_{X_i}$ es plano sobre $S'_i$.
\end{corollary}

\begin{proof}
Consideremos un recubrimiento finito $(U_j)_{1\leqslant j\leqslant m}$ de $S$ por abiertos afines y definamos, por inducción, $T_1=U_1$ y
\[
T_j=U_j\setminus\bigcup_{h<j}(U_j\cap U_h)\qquad(j\geqslant 2).
\]
Cada $T_j$ es cerrado en el abierto afín $U_j$, y los $T_j$ son disjuntos dos a dos.
Además, los $U_j\cap U_h$ son cuasicompactos porque $S$ es cuasiseparado y, por tanto, como $S$ es también cuasicompacto, son construibles en $S$ \sref{IV.1.8.1}; en consecuencia, cada $T_j$ es asimismo construible.
Bastará, pues, demostrar el corolario para un subpreesquema adecuado $T'_j$ de $S$, de espacio subyacente $T_j$ y de presentación finita sobre $S$, y para el morfismo y el Módulo deducidos respectivamente de $u$ y de $\sh{F}$ por el cambio de base $T'_j\to S$.
Para ello, observemos que, si se dota a $U_j$ de la estructura de preesquema inducida por la de $S$, la inmersión abierta $U_j\to S$ es cuasicompacta porque $S$ es cuasiseparado \sref{IV.1.2.7}, y por tanto es de
\oldpage[IV-3]{36}
presentación finita \sref{IV.1.6.2},~(i).
Basta, pues, con que $T'_j$ sea de presentación finita sobre $U_j$; dicho de otro modo, podemos reducirnos ya al caso en que $U_j=S$ y $T_j=T$ es cerrado y construible en $S$.

Escribamos $S=\operatorname{Spec}(A)$ y consideremos $A$ como límite inductivo de sus $\mathbf{Z}$-subálgebras de tipo finito, de modo que $S=\varprojlim S_\lambda$, donde los $S_\lambda$ son afines y noetherianos.
Por \sref{IV.8.3.11}, existen $\lambda$ y un subconjunto cerrado (necesariamente construible) $T_\lambda$ de $S_\lambda$ tales que $T=u_\lambda^{-1}(T_\lambda)$, donde $u_\lambda:S\to S_\lambda$ es el morfismo canónico.
Dotemos a $T_\lambda$ de una estructura de subpreesquema y tomemos $T'=T_\lambda\times_{S_\lambda}S$.
Como la inmersión canónica $T_\lambda\to S_\lambda$ es de presentación finita \sref{IV.1.6.1}, también lo es la inmersión $T'\to S$.
Puesto que $T'$ es afín, vemos finalmente que basta tratar el caso en que $T'=S$ es afín.

Entonces, con la notación de \sref{IV.8.9.1}, existen $\lambda$, un morfismo de tipo finito $u_\lambda:X_\lambda\to S_\lambda$ y un $\sh{O}_{X_\lambda}$-Módulo coherente $\sh{F}_\lambda$ tales que $X$ es isomorfo a $X_\lambda\times_{S_\lambda}S$ y $\sh{F}$ a $\sh{F}_\lambda\otimes_{\sh{O}_{X_\lambda}}\sh{O}_X$.
Aplicando \sref{IV.6.9.3} a $S_\lambda$, $X_\lambda$ y $\sh{F}_\lambda$, obtenemos subpreesquemas $S'_{\lambda i}$ de $S_\lambda$ cuyos conjuntos subyacentes forman una partición de $S_\lambda$ y tales que, poniendo $X_{\lambda i}=X_\lambda\times_{S_\lambda}S'_{\lambda i}$, el Módulo
\[
\sh{F}_{\lambda i}=\sh{F}_\lambda\otimes_{\sh{O}_{X_\lambda}}\sh{O}_{X_{\lambda i}}
\]
es plano sobre $S'_{\lambda i}$.
Los subpreesquemas $S'_i=S'_{\lambda i}\times_{S_\lambda}S$ de $S$ poseen entonces las propiedades requeridas en virtud de \sref{IV.2.1.4}.
\end{proof}

\subsection{Propiedades de permanencia de los morfismos por paso al límite proyectivo}
\label{subsection:IV.8.10}

En este parágrafo conservamos las hipótesis generales y las notaciones de \sref{IV.8.8.1}.

\begin{proposition}[8.10.1]
\label{IV.8.10.1}
Si existe $\lambda$ tal que, para $\mu\geqslant\lambda$, $f_\mu$ es un morfismo abierto (resp.\ universalmente abierto), entonces $f$ es un morfismo abierto (resp.\ universalmente abierto).
\end{proposition}

\begin{proof}
Como $X=X_\lambda\times_{Y_\lambda}Y$, la aserción relativa a los morfismos universalmente abiertos es un caso particular de \sref{IV.2.4.3},~(iii).
Supongamos, pues, que $f_\mu$ es abierto para $\mu\geqslant\lambda$; basta ver que, para todo abierto \emph{cuasicompacto} $U$ de $X$, $f(U)$ es abierto en $Y$.
Ahora bien, existe $\mu\geqslant\lambda$ tal que $U=v_\mu^{-1}(U_\mu)$, donde $U_\mu$ es un abierto cuasicompacto de $X_\mu$ \sref{IV.2.3.11}; se tiene entonces $f(v_\mu^{-1}(U_\mu))=w_\mu^{-1}(f_\mu(U_\mu))$ \sref[I]{I.3.4.8}, y por tanto $f(U)$ es abierto en $Y$.
\end{proof}

\begin{corollary}[8.10.2]
\label{IV.8.10.2}
Sea $f:X\to Y$ un morfismo.
Para que $f$ sea universalmente abierto, basta que, para todo entero $n>0$, poniendo $Y_n=Y\otimes_\mathbf{Z}\mathbf{Z}[T_1,\ldots,T_n](=\mathbf{V}_Y^n)$, $X_n=X\times_Y Y_n$ y $f_n=f\times 1_{Y_n}:X_n\to Y_n$, la proyección $f_n$ sea un morfismo abierto.
\end{corollary}

\begin{proof}
Para demostrar que $f$ es universalmente abierto, basta demostrar que sigue siéndolo sobre todo abierto afín $U$ de $Y$; por hipótesis, si $U_n=U\otimes_\mathbf{Z}\mathbf{Z}[T_1,\ldots,T_n]$ es la imagen inversa de $U$ en $Y_n$, el morfismo $f_n^{-1}(U_n)\to U_n$, restricción de $f_n$, es abierto, y se ve que podemos reducirnos al caso en que $Y=\operatorname{Spec}(A)$ es afín.
Además, basta evidentemente demostrar que, para todo morfismo $Y'\to Y$, donde $Y'=\operatorname{Spec}(A')$ es también \emph{afín}, $f'=f_{|Y'}$ es abierto.
Supongamos primero que $A'$ es una $A$-álgebra de tipo finito, y por tanto un \emph{cociente} de un álgebra de polinomios $A[T_1,\ldots,T_n]$; entonces $Y'$ es un subpreesquema cerrado de $Y_n$ y $f'$ la restricción
\oldpage[IV-3]{37}
de $f_n$ a $f_n^{-1}(Y')$; pero, para todo abierto $V$ de $X_n$, se tiene $f_n(V\cap f_n^{-1}(Y'))=f_n(V)\cap Y'$, y como por hipótesis $f_n(V)$ es abierto en $Y_n$, esto demuestra que $f'$ es también un morfismo abierto.
Cuando $A'$ es arbitraria, puede considerarse como límite inductivo de sus sub-$A$-álgebras de tipo finito $A'_\lambda$, y el hecho de que $f'$ sea un morfismo abierto se sigue de lo anterior y de \sref{IV.8.10.1}.
\end{proof}

\begin{proposition}[8.10.3]
\label{IV.8.10.3}
Supongamos que existe $\alpha$ tal que: $1^\circ$ $S_\alpha$ es cuasicompacto; $2^\circ$ los morfismos $X_\alpha\to S_\alpha$, $Y_\alpha\to S_\alpha$ son cuasicompactos y el morfismo $Y_\alpha\to S_\alpha$ es cuasiseparado; $3^\circ$ para $\alpha\leqslant\lambda\leqslant\mu$, los morfismos $u_{\lambda\mu}:S_\mu\to S_\lambda$ son planos; $4^\circ$ $\overline{f_\alpha(X_\alpha)}$ es construible en $Y_\alpha$.
Entonces, para que $f$ sea dominante, es necesario y suficiente que exista $\lambda\geqslant\alpha$ tal que $f_\lambda$ sea dominante.
\end{proposition}

\begin{proof}
Las hipótesis implican que $Y_\alpha$ es cuasicompacto y que el morfismo $f_\alpha$ es cuasicompacto \sref{IV.1.2.4}; por tanto, $\overline{f_\alpha(X_\alpha)}=Z_\alpha$ es proconstruible \sref{IV.1.9.5},~(vii) en $Y_\alpha$.
Si se pone $Z_\lambda=w_{\alpha\lambda}^{-1}(Z_\alpha)$ para $\lambda\geqslant\alpha$ y $Z=w_\alpha^{-1}(Z_\alpha)$, se tiene $Z_\lambda=\overline{f_\lambda(X_\lambda)}$ y $Z=\overline{f(X)}$ \sref[I]{I.3.4.8}, y $Z_\lambda$ es proconstruible en $Y_\lambda$ \sref{IV.1.9.5},~(vii).
Basta entonces aplicar \sref{IV.8.3.13}, sustituyendo $S_\lambda$, $Z'_\lambda$ y $Z''_\lambda$ por $Y_\lambda$, $Y_\lambda$ y $Z_\lambda$, respectivamente.
\end{proof}

\begin{proposition}[8.10.4]
\label{IV.8.10.4}
Supongamos que existe $\alpha$ tal que $Y_\alpha$ es cuasicompacto y $f_\alpha$ es de tipo finito y cuasiseparado.
Para que el morfismo $f$ sea separado, es necesario y suficiente que exista $\lambda\geqslant\alpha$ tal que $f_\lambda$ sea separado.
\end{proposition}

\begin{proof}
Como la cuestión es local sobre $Y_\alpha$ (pues $Y_\alpha$ es cuasicompacto y $L$ filtrante), podemos reducirnos al caso en que $Y_\alpha$ es afín, y por tanto cuasiseparado; la hipótesis implica entonces que $X_\alpha$ (y por consiguiente los $X_\lambda$ y $X$) es cuasicompacto y cuasiseparado.
Pongamos $X'_\lambda=X_\lambda\times_{Y_\lambda}X_\lambda$ para $\lambda\geqslant\alpha$ y $X'=X\times_Y X$; se tiene $X'_\lambda=X'_\alpha\times_{Y_\alpha}Y_\lambda$ y $X'=X'_\alpha\times_{Y_\alpha}Y$. La primera proyección $X'_\alpha\to X_\alpha$ es cuasicompacta y cuasiseparada por hipótesis \sref{IV.1.2.3},~(iii), y por tanto $X'_\lambda$ es cuasicompacto y cuasiseparado.
Observemos ahora que, si denotamos por $\Delta_\lambda$ (resp.\ $\Delta$) la diagonal de $X_\lambda\times_{Y_\lambda}X_\lambda$ (resp.\ $X\times_Y X$), se sigue de \sref[I]{I.5.3.4} y \sref[I]{I.3.4.8} que $\Delta_\mu$ (resp.\ $\Delta$) es la imagen inversa de $\Delta_\lambda$ por la aplicación $v'_{\lambda\mu}:X'_\mu\to X'_\lambda$ (resp.\ $v'_\lambda:X'\to X'_\lambda$).
Por otra parte, $\Delta_\lambda$ es \emph{construible} en $X'_\lambda$: en efecto, como $f_\lambda$ es cuasiseparado, la inmersión diagonal $X_\lambda\to X'_\lambda$ es cuasicompacta, y localmente de presentación finita porque $f_\lambda$ es de tipo finito \sref[I]{I.1.4.3} y \sref[I]{I.1.5.4},~(iii); se sigue entonces de \sref{IV.1.8.4.1} que $\Delta_\lambda$ es localmente construible, y por tanto construible, ya que $X'_\lambda$ es cuasicompacto.
Podemos aplicar ahora \sref{IV.8.3.12}, sustituyendo $S_\lambda$ y $Z_\lambda$ por $X'_\lambda$ y $\Delta_\lambda$, respectivamente.
\end{proof}

\begin{theorem}[8.10.5]
\label{IV.8.10.5}
Supongamos $S_0$ cuasicompacto, $X_\alpha$ e $Y_\alpha$ de presentación finita sobre $S_\alpha$, y sea $f_\alpha:X_\alpha\to Y_\alpha$ un $S_\alpha$-morfismo.
Consideremos, para un morfismo, la propiedad de ser:
\begin{enumerate}[label=(\roman*)]
\item un isomorfismo;
\item[\rm(i\,\emph{bis})] un monomorfismo;
\item una inmersión;
\item una inmersión abierta;
\item una inmersión cerrada;
\item separado;
\item sobreyectivo;
\oldpage[IV-3]{38}
\item radicial;
\item afín;
\item cuasi-afín;
\item finito;
\item cuasifinito;
\item propio.
\end{enumerate}
Entonces, si $\mathbf{P}$ denota una de las propiedades anteriores, para que $f$ posea la propiedad $\mathbf{P}$, es necesario y suficiente que exista $\lambda\geqslant\alpha$ tal que $f_\lambda$ posea la propiedad $\mathbf{P}$ (en cuyo caso $f_\mu$ la posee también para $\mu\geqslant\lambda$).

Si se supone además que $S_0$ es cuasiseparado, la misma conclusión vale cuando $\mathbf{P}$ es la propiedad de ser:
\begin{enumerate}[label=(\roman*)]
\setcounter{enumi}{12}
\item proyectivo;
\item cuasiproyectivo.
\end{enumerate}
\end{theorem}

\begin{proof}
El caso en que $\mathbf{P}$ es una de las propiedades~(i) o~(v) se ha incluido en el enunciado solo a título recordatorio; en estos casos, el teorema se sigue de lo demostrado respectivamente en \sref{IV.8.8.2.4} y \sref{IV.8.10.4}.
Además, teniendo en cuenta \sref[I]{I.5.4.1} y \sref[I]{I.5.3.4}, (v) se sigue también de (iv).
El caso (i\,\emph{bis}) se sigue inmediatamente de (i), utilizando \sref[I]{I.5.3.8} y observando (como ya hemos utilizado en \sref{IV.8.10.4}) que la diagonal $\Delta_f$ se deduce de $\Delta_{f_\lambda}$ por el cambio de base $S\to S_\lambda$.

Observaremos, por otra parte, que (vi), (vii) y (xi) son en realidad condiciones sobre las \emph{fibras} $f^{-1}(y)$ de los morfismos considerados, teniendo en cuenta la transitividad de las fibras por cambio de base \sref[I]{I.3.6.4}: la condición~(vi) significa, en efecto, que todas las fibras deben ser no vacías; la condición~(vii), que deben ser radiciales \sref[I]{I.3.5.8}; y la condición~(xi), que deben ser finitas \sref[III]{III.6.2.2}, \sref[III]{III.6.2.3} y \sref[I]{I.6.4.4}, teniendo en cuenta que $f$ y los $f_\lambda$ son morfismos de tipo finito por \sref{IV.1.5.4},~(v).
El teorema en estos tres casos será todavía consecuencia de un resultado general sobre este tipo de propiedades, que se establecerá en \sref{IV.9.3.3}; aplazamos, pues, hasta entonces la demostración del teorema en el caso~(xi) (el lector podrá comprobar, naturalmente, que salvo en los n\textsuperscript{os}~{\hyperlink{ega4.c04.subsec.8-11}{8.11}} y {\hyperlink{ega4.c04.subsec.8-12}{8.12}}, no utilizaremos el teorema en estos casos antes de \sref{IV.9.3.3}, y que \sref{IV.8.11} y \sref{IV.8.12} no se utilizarán antes de \sref{IV.9.3.3}).

Para los casos que quedan por demostrar, podemos limitarnos a probar que la condición del enunciado es necesaria, pues todas las propiedades $\mathbf{P}$ consideradas son invariantes por cambio de base (véanse los caps.~I\textsuperscript{er} y II para las referencias relativas a cada una de estas propiedades).
Podemos suponer además que $S_\alpha=S_0$ y que $Y_\alpha=S_\alpha$, y por tanto $Y_\lambda=S_\lambda$ para todo $\lambda\geqslant\alpha$.
Por último, las propiedades~(i) a~(xii) son locales sobre $S_0$; así, como $S_0$ es unión finita de abiertos afines y $L$ es filtrante, podemos limitarnos a demostrarlas para $S_0=\operatorname{Spec}(A_0)$ \emph{afín} (y por tanto cuasiseparado).
Denotaremos por $A_\lambda$ (resp.\ $A$) el anillo de $S_\lambda$ (resp.\ $S$).

\emph{Casos} (ii), (iii), (iv): Supongamos que $f$ es una inmersión (resp.\ una inmersión abierta, una inmersión cerrada), y sea $X'$ el subpreesquema (resp.\ inducido sobre un abierto, resp.\ cerrado) de $S$ asociado a $f$, que es por tanto un $S$-preesquema de presentación finita.
\oldpage[IV-3]{39}
En virtud de \sref{IV.8.6.3}, existen entonces $\lambda\geqslant\alpha$ y un subpreesquema $X'_\lambda$ de $S_\lambda$, de presentación finita sobre $S_\lambda$, tales que $X'$ es isomorfo a $X'_\lambda\times_{S_\lambda}S$.
Para todo $\mu\geqslant\lambda$, $X'_\mu=X'_\lambda\times_{S_\lambda}S_\mu$ es, pues, un subpreesquema (resp.\ inducido sobre un abierto, resp.\ cerrado) de $S_\mu$, y se sigue entonces de \sref{IV.8.8.2.4} y \sref{IV.8.8.2.5} que existen $\mu\geqslant\lambda$ y un isomorfismo $g_\mu:X_\mu\to X'_\mu$ tales que $g=g_\mu\times 1_S$ es el isomorfismo $X\to X'$ asociado a $f$; de donde la conclusión.

\emph{Casos} (vi) \emph{y} (vii): Sabemos por \sref{IV.1.8.4.1} que $Z_\alpha=f_\alpha(X_\alpha)$ es construible en $S_\alpha$; si ponemos $Z_\lambda=u_{\alpha\lambda}^{-1}(Z_\alpha)$ para $\lambda\geqslant\alpha$ y $Z=u_\alpha^{-1}(Z_\alpha)$, se tiene $Z_\lambda=f_\lambda(X_\lambda)$ y $Z=f(X)$ \sref[I]{I.3.4.8}.
Como $Z_\alpha$ es proconstruible en $S_\alpha$, y en virtud de \sref{IV.8.3.11} la aplicación canónica $\varinjlim\mathfrak{C}(S_\lambda)\to\mathfrak{C}(S)$ es inyectiva, la relación $f(X)=S$ implica la existencia de $\lambda\geqslant\alpha$ tal que $f_\lambda(X_\lambda)=S_\lambda$, lo que demuestra el teorema en el caso~(vi).
Para demostrarlo en el caso~(vii), observemos que el morfismo estructural $X_\alpha\times_{S_\alpha}X_\alpha\to S_\alpha$ es de presentación finita, pues la primera proyección $X_\alpha\times_{S_\alpha}X_\alpha\to X_\alpha$ lo es \sref{IV.1.6.2}. Basta, pues, en virtud de \sref{IV.1.8.7.1}, aplicar el caso~(vi) al morfismo diagonal
\[
  \Delta_{f_\alpha}:X_\alpha\longrightarrow X_\alpha\times_{S_\alpha}X_\alpha,
\]
observando que $\Delta_{f_\lambda}=\Delta_{f_\alpha}\times1_{S_\lambda}$ y $\Delta_f=\Delta_{f_\alpha}\times1_S$ \sref[I]{I.5.3.4} y \sref[I]{I.3.3.11}.

\emph{Casos} (viii) \emph{y} (ix): Como $S=\operatorname{Spec}(A)$ es afín, decir que $f:X\to S$ es afín (resp.\ cuasi-afín) significa que existen un entero $r$ y una inmersión cerrada (resp.\ una inmersión) $j:X\to\mathbf{V}^r_S=\operatorname{Spec}(A[T_1,\ldots,T_r])$ de $S$-preesquemas, pues $f$ es de tipo finito y $S$ cuasicompacto \sref[II]{II.5.1.9}.
Como $\mathbf{V}^r_S=\mathbf{V}^r_{S_\lambda}\times_{S_\lambda}S$, y $\mathbf{V}^r_{S_\lambda}$ es un $S_0$-preesquema de presentación finita, se sigue de \sref{IV.8.8.2},~(i), aplicado al $S$-morfismo $j$, que existen $\lambda$ y un $S_\lambda$-morfismo $j_\lambda:X_\lambda\to\mathbf{V}^r_{S_\lambda}$ tales que $j=j_\lambda\times 1_S$; aplicando (iv) (resp.\ (ii)) a $j$, se deduce que existe $\mu\geqslant\lambda$ tal que $j_\mu$ es una inmersión cerrada (resp.\ una inmersión); por tanto, $f_\mu$ es afín (resp.\ cuasi-afín).

\emph{Caso} (x): Por hipótesis, $X=\operatorname{Spec}(B)$, donde $B$ es una $A$-álgebra que, como $A$-módulo, es de presentación finita \sref{IV.1.4.7}; se sigue, pues, de \sref{IV.8.5.2},~(ii) que existen $\lambda$ y un $A_\lambda$-módulo de presentación finita $B_\lambda$ tales que el $A$-módulo $B$ es isomorfo a $B_\lambda\otimes_{A_\lambda}A$.
La estructura de $A$-álgebra de $B$ está definida por un $A$-homomorfismo $m:B\otimes_A B\to B$; como $B\otimes_A B=(B_\lambda\otimes_{A_\lambda}B_\lambda)\otimes_{A_\lambda} A$, existen, por \sref{IV.8.5.2},~(i), $\mu\geqslant\lambda$ y un $A_\mu$-homomorfismo $m_\mu:B_\mu\otimes_{A_\mu}B_\mu\to B_\mu$ tales que $m=m_\mu\otimes 1$.
Considerando los diagramas habituales que expresan la asociatividad y la conmutatividad de $m$, se ve, aplicando de nuevo \sref{IV.8.5.2},~(i), que podemos suponer $\nu$ suficientemente grande para que $m_\nu$ defina sobre $B_\nu$ una multiplicación asociativa y conmutativa.
De igual modo, podemos suponer $\nu$ suficientemente grande para que $B_\nu$, así definida, admita un elemento unidad.
Si $X'_\nu=\operatorname{Spec}(B_\nu)$, está claro entonces que $X$ es $S$-isomorfo a $X'_\nu\times_{S_\nu}S$; por tanto, en virtud de (i), existe $\rho\geqslant\nu$ tal que $X_\rho$ y $X'_\rho$ son $S_\rho$-isomorfos, lo que termina la demostración en este caso.

Para demostrar el teorema en el caso~(xii), demostraremos primero la proposición siguiente:
\end{proof}

\begin{proposition}[8.10.5.1]
\label{IV.8.10.5.1}
(Lema de Chow para morfismos de presentación finita). --- Sea $\mathrm{A}$ un anillo, sean $X$ e $Y$ dos $\mathrm{A}$-preesquemas de presentación finita y sea $f:X\to Y$ un $\mathrm{A}$-morfismo separado. Entonces existen dos $\mathrm{A}$-preesquemas $X'$ y $P$ de presentación finita, y $\mathrm{A}$-morfismos $p:P\to Y$, $j:X'\to P$ y $g:X'\to X$, tales que el diagrama
\oldpage[IV-3]{40}
\[
\begin{tikzcd}[column sep=large,row sep=large]
X' \arrow[r,"j"] \arrow[d,"g"'] & P \arrow[d,"p"] \\
X \arrow[r,"f"'] & Y
\end{tikzcd}
\]
es conmutativo y tales que: $1^\circ$ $p$ es proyectivo; $2^\circ$ $g$ es proyectivo y sobreyectivo; $3^\circ$ $j$ es una inmersión abierta.
\end{proposition}

\begin{proof}
Elijamos $\mathrm{A}_0\subset\mathrm{A}$, $X_0$, $Y_0$ y $f_0$ como en \sref{IV.8.9.1}, con $Y_0$ noetheriano y $f_0$ de tipo finito; por \sref{IV.8.10.4}, podemos suponer además que $f_0$ es separado. El lema de Chow \sref[II]{II.5.6.1} proporciona entonces morfismos de tipo finito $p_0:P_0\to Y_0$, $g_0:X'_0\to X_0$ y $j_0:X'_0\to P_0$ tales que
\[
\begin{tikzcd}[column sep=large,row sep=large]
X'_0 \arrow[r,"j_0"] \arrow[d,"g_0"'] & P_0 \arrow[d,"p_0"] \\
X_0 \arrow[r,"f_0"'] & Y_0
\end{tikzcd}
\]
es conmutativo, $p_0$ es proyectivo, $g_0$ es proyectivo y sobreyectivo, y $j_0$ es una inmersión abierta. Las aserciones se siguen de la invariancia de estas propiedades por cambio de base \sref[II]{II.5.5.5},~(iii), \sref[I]{I.3.5.2} y \sref{IV.4.3.2}.
\end{proof}

\begin{proof}[Continuación de la demostración de \sref{IV.8.10.5}]
\emph{Caso (xii).} Apliquemos \sref{IV.8.10.5.1} a $f_0:X_0\to S_0$. Obtenemos un diagrama conmutativo
\[
\begin{tikzcd}[column sep=large,row sep=large]
X'_0 \arrow[r,"j_0"] \arrow[d,"g_0"'] & P_0 \arrow[d,"p_0"] \\
X_0 \arrow[r,"f_0"'] & S_0
\end{tikzcd}
\]
donde $p_0$ es proyectivo, $g_0$ es proyectivo y sobreyectivo, y $j_0$ es una inmersión abierta. Para cada $\lambda$ se obtiene el diagrama análogo, en el que $p_\lambda=p_0\times1_{S_\lambda}$, $g_\lambda=g_0\times1_{S_\lambda}$ y $j_\lambda=j_0\times1_{S_\lambda}$ poseen las mismas propiedades respectivas, y lo mismo ocurre con los morfismos límite inverso $p=p_0\times1_S$, $g=g_0\times1_S$ y $j=j_0\times1_S$. Como $g$ es propio \sref[II]{II.5.5.3}, también lo es $p\circ j=f\circ g$ \sref[II]{II.5.4.2}; y como $p$ es separado, $j$ es propio, y por tanto una inmersión cerrada. Apliquemos el caso~(iv) a $j$, observando que $X'_0$ y $P_0$ son $S_0$-preesquemas de presentación finita por \sref{IV.8.10.5.1} y \sref{IV.1.6.2}. Existe $\lambda$ tal que $j_\lambda$ es una inmersión cerrada, y por tanto propia \sref[II]{II.5.4.2}. Así,
\[
  f_\lambda\circ g_\lambda=p_\lambda\circ j_\lambda
\]
es propio \sref[II]{II.5.5.3} y \sref[II]{II.5.4.2}. Como $g_\lambda$ es sobreyectivo y, por la hipótesis sobre $f$ y el caso~(v), se puede suponer que $f_\lambda$ es separado, \sref[II]{II.5.4.3} demuestra que $f_\lambda$ es propio.

\emph{Casos (xiii) y (xiv).} Por el caso~(xii) y \sref[II]{II.5.5.3}, que se aplica porque los $S_\lambda$ son cuasicompactos y cuasiseparados por \sref{IV.1.7.19}, basta considerar el caso en que $f$ es cuasiproyectivo. Supongamos que existe un $\sh{O}_X$-Módulo invertible $\sh{L}$ que es amplio para $f$. Como $S_0$ es cuasicompacto y cuasiseparado, también lo es $X_0$ \sref{IV.1.2.3}; existen, pues, $\lambda$ y un $\sh{O}_{X_\lambda}$-Módulo cuasicoherente de presentación finita $\sh{L}_\lambda$ tales que $\sh{L}=v_\lambda^*(\sh{L}_\lambda)$ \sref{IV.8.5.2},~(ii). Por \sref{IV.8.5.5}, podemos suponer además que $\sh{L}_\lambda$ es invertible. La aserción se sigue ahora del lema más preciso que figura a continuación.

\begin{lemma}[8.10.5.2]
\label{IV.8.10.5.2}
\oldpage[IV-3]{41}
Supongamos que $S_0$ es cuasicompacto, y sea $\sh{L}_\lambda$ un $\sh{O}_{X_\lambda}$-Módulo invertible. Para que $\sh{L}$ sea amplio para $f$ (resp.\ muy amplio para $f$), es necesario y suficiente que exista $\mu\geqslant\lambda$ tal que $\sh{L}_\mu$ sea amplio para $f_\mu$ (resp.\ muy amplio para $f_\mu$).
\end{lemma}

\begin{proof}
La condición es manifiestamente suficiente \sref[II]{II.4.4.10} y \sref[II]{II.4.6.13}. Recíprocamente, como los $S_\lambda$ son cuasicompactos y los $f_\lambda$ son de tipo finito, tras sustituir $\sh{L}$ por una potencia tensorial adecuada, podemos reducirnos al caso en que $\sh{L}$ es muy amplio \sref[II]{II.4.6.11}. La cuestión es local sobre $S_0$ por \sref[II]{II.4.4.5} y por el hecho de que $L$ es filtrante, de modo que podemos suponer que $S_0$, y por tanto $S$, es afín. Por \sref[II]{II.4.4.1},~(ii) y \sref[II]{II.4.1.2}, existe una $S$-inmersión
\[
  j:X\longrightarrow\mathbf{P}^r_S=P
\]
tal que $\sh{L}\simeq j^*(\sh{O}_P(1))$. Por \sref{IV.8.8.2},~(i), el caso~(ii) y \sref[II]{II.4.1.3}, existen $\mu\geqslant\lambda$ y una inmersión $j_\mu:X_\mu\to\mathbf{P}^r_{S_\mu}=P_\mu$ tales que $j=j_\mu\times1_S$. Entonces \sref[II]{II.4.1.3.2} y \sref{IV.8.5.2.5} proporcionan $\nu\geqslant\mu$ tal que
\[
  \sh{L}_\nu\simeq j_\nu^*(\sh{O}_{P_\nu}(1)),
\]
lo que demuestra que $\sh{L}_\nu$ es muy amplio para $f_\nu$ \sref[II]{II.4.4.2}.
\end{proof}
\end{proof}

\subsection{Aplicación a los morfismos cuasifinitos}
\label{subsection:IV.8.11}
\label{IV.8.11}

Nos proponemos demostrar en este parágrafo los dos teoremas siguientes:

\begin{theorem}[8.11.1]
\label{IV.8.11.1}
Sea $f:X\to Y$ un morfismo propio, localmente de presentación finita y cuasifinito.
Entonces el morfismo $f$ es finito.
\end{theorem}

\begin{theorem}[8.11.2]
\label{IV.8.11.2}
Sea $f:X\to Y$ un morfismo localmente de presentación finita, cuasifinito y separado.
Entonces el morfismo $f$ es cuasi-afín y, con mayor razón, cuasiproyectivo.
\end{theorem}

\begin{remark}[8.11.3]
\label{IV.8.11.3}
Veremos más adelante que, para la demostración de \sref{IV.8.11.1} y \sref{IV.8.11.2}, podemos reducirnos al caso en que $Y$ es localmente noetheriano; obsérvese que en este caso se obtiene así otra demostración del teorema de Chevalley \sref[III]{III.4.4.2}.
\end{remark}

\begin{env}[8.11.4]
\label{IV.8.11.4}
Las hipótesis y las conclusiones de \sref{IV.8.11.1} y \sref{IV.8.11.2} son todas locales sobre $Y$.
En efecto, esto se sigue de \sref{IV.1.6.1}, \sref{IV.1.2.6}, \sref[II]{II.5.1.1}, \sref[II]{II.5.4.1} y \sref[II]{II.6.2.2}.
Podemos suponer, por tanto, $Y=\operatorname{Spec}(A)$ afín.
Sabemos que existen entonces un subanillo $A_0$ de $A$, que es una $\bb{Z}$-álgebra de tipo finito, y un morfismo de tipo finito $f_0:X_0\to\operatorname{Spec}(A_0)$ tales que $X$ se identifica con $X_0\otimes_{A_0}A$ y $f$ con $f_0\times 1$ \sref{IV.8.9.1}.
Además, $A$ puede considerarse como límite inductivo de sus subanillos que contienen $A_0$ y que son $\bb{Z}$-álgebras de tipo finito; utilizando el método de \sref{IV.8.1.2},~$c$)) y \sref{IV.8.10.5},~(v),~(xi) y (xii)), se ve que basta demostrar los teoremas \sref{IV.8.11.1} y \sref{IV.8.11.2} para $f_0$.
Supongamos en adelante $Y$ noetheriano; utilizando ahora el método de \sref{IV.8.1.2},~$a$)) y \sref{IV.8.10.5},~(v),~(ix),~(x),~(xi) y (xii)), se puede sustituir $Y$ por $\operatorname{Spec}(\sh{O}_y)$, donde $y$ es un
\oldpage[IV-3]{42}
punto arbitrario de $Y$, de modo que finalmente podemos suponer que $Y=\operatorname{Spec}(A)$, donde $A$ es un \emph{anillo local noetheriano}.
Sea $\mathfrak{m}$ el ideal maximal de $A$, y sea $\hat{A}$ el completado de $A$ para la topología $\mathfrak{m}$-ádica. El anillo $\hat{A}$ es local noetheriano y $\operatorname{Spec}(\hat{A})\to\operatorname{Spec}(A)$ es fielmente plano y cuasicompacto \sref[0]{0.7.3.5}; aplicando \sref{IV.2.7.1},~(i), (vii), (xiv), (xv) y (xvi), podemos reducirnos, por tanto, al caso en que $A$ es \emph{completo}.
Se sigue entonces de \sref[II]{II.6.2.6} que $X=X'\sqcup X''$, donde $X'$ es un $Y$-esquema finito y $X''$ un $Y$-esquema cuasifinito tal que $X''\cap f^{-1}(y)=\varnothing$.
\end{env}

Situémonos primero bajo las hipótesis de \sref{IV.8.11.1}; como $f$ es propio, $X''$, que es cerrado en $X$, es propio sobre $Y$ \sref[II]{II.5.4.10}, de modo que $f(X'')$ es cerrado en $Y$; pero $y$ no pertenece a $f(X'')$ y, además, $y$ es adherente a todo punto de $Y$, por lo que $f(X'')=\varnothing$; en consecuencia, $X''$ es vacío y $f$ es \emph{finito}.

Situémonos ahora bajo las hipótesis de \sref{IV.8.11.2} y, restringiéndonos aún más (como permite lo anterior) al caso en que $Y=\operatorname{Spec}(A)$ es afín y noetheriano de dimensión finita, razonemos además por inducción sobre la dimensión de $Y$.
Reduciéndonos como antes al caso en que $A$ es además local y completo, se tiene $\dim(\sh{O}_y)=\dim(A)=\dim(Y)$ y, para todo $z\neq y$, $\dim(\sh{O}_z)<\dim(\sh{O}_y)$, de modo que $\dim(Y-\{y\})<\dim(Y)$.
Ahora bien, por hipótesis se tiene $f(X'')\subset Y-\{y\}$ y la restricción de $f$ a $X''$ es evidentemente un morfismo cuasifinito y separado; aplicando la hipótesis de inducción a $Y-\{y\}$ y a $X''$, se ve que $X''$ es cuasi-afín sobre $Y-\{y\}$; pero, como el abierto $Y-\{y\}$ es cuasi-afín sobre $Y$ porque $Y$ es noetheriano, $X''$ es también cuasi-afín sobre $Y$ \sref[II]{II.5.1.10},~(ii)); y como además $X'$ es finito (y \emph{a fortiori} afín) sobre $Y$, $X$ es cuasi-afín sobre $Y$ \sref[II]{II.4.6.17} y \sref[II]{II.5.1.2},~$c'$)).

\begin{proposition}[8.11.5]
\label{IV.8.11.5}
Sea $f:X\to Y$ un morfismo de presentación finita.
Las propiedades siguientes son equivalentes:
\begin{enumerate}[label=\emph{\alph*)}]
  \item $f$ es una inmersión cerrada.
  \item $f$ es un monomorfismo propio.
  \item $f$ es propio y, para todo $y\in Y$, $f^{-1}(y)$ es radicial y geométricamente reducido sobre $k(y)$ (es decir, vacío o $k(y)$-isomorfo a $\operatorname{Spec}(k(y))$).
\end{enumerate}
\end{proposition}

Está claro que \emph{a)} entraña \emph{b)}.
Para ver que \emph{b)} entraña \emph{c)}, observemos \sref[I]{I.5.3.8} que, para todo $y\in Y$, el morfismo $f^{-1}(y)\to\operatorname{Spec}(k(y))$ deducido de $f$ por extensión de la base es un monomorfismo, y por tanto inyectivo; en consecuencia, $f^{-1}(y)$ es vacío o se reduce a un punto, y en todo caso es afín.
Además, si $A$ es el anillo de $f^{-1}(y)$, el homomorfismo canónico $A\otimes_k A\to A$ es biyectivo \sref[I]{I.5.3.8}.
Esto implica evidentemente que $A=k$, pues de lo contrario existiría un elemento $a\in A$ que no perteneciera a $k$, y las imágenes de $a\otimes 1$ y de $1\otimes a$ en $A$ serían ambas iguales a $a$, mientras que $a\otimes 1\neq1\otimes a$, ya que $1$ y $a$ forman un sistema libre sobre $k$.

Queda por demostrar que \emph{c)} entraña \emph{a)}.
Se sigue primero de \sref{IV.8.11.1} que $f$ es un morfismo \emph{finito}, de modo que $X=\operatorname{Spec}(\sh{A})$, donde $\sh{A}$ es una $\sh{O}_Y$-Álgebra finita.
Basta, pues, demostrar que el homomorfismo canónico $\sh{O}_Y\to\sh{A}$ es \emph{sobreyectivo} \sref[II]{II.1.4.10}, o, equivalentemente, que para todo $y\in Y$ el homomorfismo $\sh{O}_y\to\sh{A}_y$ es sobreyectivo.
Pero, por hipótesis,
\oldpage[IV-3]{43}
$f^{-1}(y)=\operatorname{Spec}(\sh{A}_y/\mathfrak{m}_y\sh{A}_y)$ \sref[II]{II.1.5.5} es tal que el homomorfismo correspondiente $k(y)=\sh{O}_y/\mathfrak{m}_y\to\sh{A}_y/\mathfrak{m}_y\sh{A}_y$ es biyectivo; como $\sh{A}_y$ es un $\sh{O}_y$-módulo de tipo finito, el lema de Nakayama demuestra que $\sh{O}_y\to\sh{A}_y$ es sobreyectivo, lo que termina la demostración.

\begin{remark}[8.11.5.1]
\label{IV.8.11.5.1}
Obsérvese que el razonamiento anterior demuestra que, si $f$ es un \emph{monomorfismo}, entonces, para todo $y\in Y$, $f^{-1}(y)$ es vacío o $k(y)$-isomorfo a $\operatorname{Spec}(k(y))$.
\end{remark}

\begin{proposition}[8.11.6]
\label{IV.8.11.6}
Si un morfismo $f:X\to Y$ de presentación finita es un homeomorfismo universal, es finito, sobreyectivo y radicial (siendo cierto el recíproco por \sref{IV.2.4.5},~(i))).
\end{proposition}

\begin{proof}
En efecto, como $f$ es de tipo finito, universalmente cerrado y separado en virtud de \sref{IV.2.4.4}, es propio por definición \sref[II]{II.5.4.1}.
Como es evidentemente cuasifinito \sref[II]{II.6.2.3}, es finito por \sref{IV.8.11.1}.
Sabemos además que es radicial \sref{IV.2.4.4}, y evidentemente sobreyectivo.
\end{proof}

\subsection{Nueva demostración y generalización del ``Teorema principal'' de Zariski}
\label{subsection:IV.8.12}
\label{IV.8.12}

\begin{lemma}[8.12.1]
\label{IV.8.12.1}
Sea $f:X\to Y$ un morfismo cuasicompacto y cuasiseparado, sea $\sh{C}$ una $\sh{O}_Y$-Álgebra cuasicoherente y sea $Z=\operatorname{Spec}(\sh{C})$, que es un $Y$-preesquema afín sobre $Y$.
Sea $g:X\to Z$ un $Y$-morfismo y sea $\varphi=\sh{A}(g):\sh{C}\to f_*(\sh{O}_X)$ el $\sh{O}_Y$-homomorfismo correspondiente de $\sh{O}_Y$-Álgebras \sref[II]{II.1.2.7}.
Supongamos que $g$ es una inmersión.
Entonces, para que la imagen cerrada de $X$ por $g$ \sref[I]{I.9.5.3} sea igual a $Z$, es necesario y suficiente que $\varphi$ sea inyectivo; $g$ es entonces una inmersión abierta.
\end{lemma}

\begin{proof}
La hipótesis implica que $f_*(\sh{O}_X)$ es una $\sh{O}_Y$-Álgebra cuasicoherente \sref{IV.1.7.5}; además, como el morfismo canónico $h:Z\to Y$ es afín, y por tanto cuasicompacto y cuasiseparado \sref{IV.1.2.2} y \sref{IV.1.1.2}, $g_*(\sh{O}_X)$ es una $\sh{O}_Z$-Álgebra cuasicoherente \sref{IV.1.7.5}.
Dicho esto, afirmar que la imagen cerrada de $X$ por $g$ es igual a $Z$ significa \sref[I]{I.9.5.2} que el homomorfismo canónico $\theta:\sh{O}_Z\to g_*(\sh{O}_X)$ es inyectivo.
Pero se tiene $h_*(\sh{O}_Z)=\sh{C}$ por definición de $Z$ \sref[II]{II.1.3.1}, y $h_*(g_*(\sh{O}_X))=f_*(\sh{O}_X)$.
Como $Z$ es afín sobre $Y$, equivale a decir que el homomorfismo $\theta:\sh{O}_Z\to g_*(\sh{O}_X)$ es inyectivo o que el homomorfismo correspondiente $\varphi=h_*(\theta):\sh{C}\to f_*(\sh{O}_X)$ es inyectivo \sref[I]{I.1.3.9}.
El hecho de que $g$ sea entonces una inmersión abierta se sigue de \sref[I]{I.9.5.10} y de la hipótesis.
\end{proof}

\begin{lemma}[8.12.2]
\label{IV.8.12.2}
Sea $Y$ un preesquema cuasicompacto y cuasiseparado, sea $f:X\to Y$ un morfismo cuasiseparado de tipo finito, sea $\sh{C}$ una $\sh{O}_Y$-Álgebra cuasicoherente y sea $Z=\operatorname{Spec}(\sh{C})$.
Sea $g:X\to Z$ un $Y$-morfismo y sea $\varphi:\sh{C}\to f_*(\sh{O}_X)$ el $\sh{O}_Y$-homomorfismo correspondiente de $\sh{O}_Y$-Álgebras.
Sea $(\sh{C}_\lambda)$ la familia creciente filtrante de sub-$\sh{O}_Y$-Álgebras cuasicoherentes de tipo finito de $\sh{C}$ (cuya unión es $\sh{C}$ \sref[I]{I.9.6.6} y \sref{IV.1.7.9}); pongamos $Z_\lambda=\operatorname{Spec}(\sh{C}_\lambda)$ y sea $g_\lambda$ el morfismo compuesto $X\xrightarrow{g}Z\to Z_\lambda$.
Entonces las condiciones siguientes son equivalentes:
\begin{enumerate}[label=\emph{\alph*)}]
  \item $g$ es una inmersión.
  \item Existe $\lambda$ tal que $g_\lambda$ es una inmersión.
\end{enumerate}
\oldpage[IV-3]{44}
Además, cuando $g_\lambda$ es una inmersión, lo mismo ocurre con $g_\mu$ para $\mu\geqslant\lambda$.

Basta aplicar \sref[II]{II.3.8.4}, sustituyendo $\sh{L}$ por $\sh{O}_X$ y $\sh{S}$ por $\sh{C}[\mathrm{T}]$, y teniendo en cuenta \sref[II]{II.3.1.7}.
\end{lemma}

\begin{proposition}[8.12.3]
\label{IV.8.12.3}
Sea $Y$ un preesquema cuasicompacto y cuasiseparado, y sea $f:X\to Y$ un morfismo separado de tipo finito.
Sea $\sh{B}=f_*(\sh{O}_X)$, que es una $\sh{O}_Y$-Álgebra cuasicoherente \sref[I]{I.9.2.2}; sea $\sh{C}$ la clausura entera de $\sh{O}_Y$ en $\sh{B}$, que es una $\sh{O}_Y$-Álgebra cuasicoherente \sref[II]{II.6.3.4}; pongamos $Z=\operatorname{Spec}(\sh{C})$, y sea $g:X\to Z$ el $Y$-morfismo correspondiente a la inyección $\sh{C}\to\sh{B}=f_*(\sh{O}_X)$; para todo $\lambda$, sea $g_\lambda:X\to Z_\lambda=\operatorname{Spec}(\sh{C}_\lambda)$ el $Y$-morfismo correspondiente a la inyección $\sh{C}_\lambda\to\sh{B}=f_*(\sh{O}_X)$.
Entonces las condiciones siguientes son equivalentes:
\begin{enumerate}[label=\emph{\alph*)}]
  \item[\emph{a)}] Existe una factorización de $f$ de la forma
\[
X\xrightarrow{f'}Y'\xrightarrow{u}Y
\]
donde $f'$ es una inmersión y $u$ un morfismo finito.
  \item[\emph{a$'$)}] Existe una factorización $X\xrightarrow{f'}Y'\xrightarrow{u}Y$ de $f$, donde $f'$ es una inmersión abierta y $u$ un morfismo finito.
  \item[\emph{b)}] El morfismo $g:X\to Z$ es una inmersión.
  \item[\emph{c)}] Existe $\lambda$ tal que $g_\lambda:X\to Z_\lambda$ es una inmersión.
\end{enumerate}
Además, cuando $g_\lambda$ es una inmersión, $g$ es una inmersión abierta, $g(X)$ es denso en $Z$, y existe $\lambda$ tal que, para $\mu\geqslant\lambda$, $g_\mu$ es una inmersión abierta.
\end{proposition}

Como el homomorfismo $\varphi:\sh{C}\to f_*(\sh{O}_X)$ es inyectivo, se sigue de \sref{IV.8.12.1} que, si $g$ es una inmersión, es una inmersión abierta y $g(X)$ es denso en $Z$; lo mismo ocurre con $g_\lambda$.
El hecho de que \emph{a)} entrañe \emph{a$'$)} se sigue también de \sref{IV.8.12.1}, aplicado sustituyendo $Z$ por $Y'$ y $g$ por $f'$ (pues $Y'$ es finito sobre $Y$, y por tanto afín sobre $Y$): en efecto, si $Y''$ es la imagen cerrada de $X$ por $f'$, $Y''$ es finito sobre $Y$ y $f'$ se factoriza como $X\xrightarrow{f''}Y''\xrightarrow{j}Y'$, donde $j$ es la inyección canónica y $f''$ es una inmersión \sref[I]{I.4.1.10}; se sigue, pues, de \sref{IV.8.12.1} que $f''$ es una inmersión abierta.

La equivalencia de \emph{b)} y \emph{c)} se sigue de \sref{IV.8.12.2}, al igual que el hecho de que $g_\lambda$ sea entonces una inmersión para $\lambda$ suficientemente grande.
Está claro que \emph{c)} entraña \emph{a)}, puesto que $Z_\lambda$ es finito sobre $Y$ \sref[II]{II.6.3.4} y \sref[II]{II.6.1.2}.
Demostremos finalmente que \emph{a)} entraña \emph{c)}.
Hemos visto más arriba que $Y'$ es la imagen cerrada de $X$ por $f'$, y se sigue entonces de \sref{IV.8.12.1} que, poniendo $\sh{B}'=u_*(\sh{O}_{Y'})$, el homomorfismo $\varphi':\sh{B}'\to\sh{B}=f_*(\sh{O}_X)$ es \emph{inyectivo}.
Pero, como por hipótesis $\sh{B}'$ es una $\sh{O}_Y$-álgebra finita, se identifica, por definición de $\sh{C}$, con una de las $\sh{O}_Y$-subálgebras $\sh{C}_\lambda$, lo que demuestra \emph{c)}.

Diremos que un morfismo $f:X\to Y$, donde $Y$ es cuasicompacto y cuasiseparado, es \emph{pseudofinito} si es de tipo finito y satisface la condición \emph{a)} de \sref{IV.8.12.3} (en cuyo caso es necesariamente separado).

\begin{corollary}[8.12.4]
\label{IV.8.12.4}
Sea $Y$ un preesquema cuasicompacto y cuasiseparado, y sea $f:X\to Y$ un morfismo.
\oldpage[IV-3]{45}
\begin{enumerate}[label=(\roman*)]
  \item Supongamos $f$ pseudofinito.
Entonces, para todo morfismo $Y'\to Y$, donde $Y'$ es cuasicompacto y cuasiseparado, $f_{(Y')}:X'=X_{(Y')}\to Y'$ es pseudofinito.
  \item Sea $(U_\lambda)$ un recubrimiento de $Y$ por abiertos cuasicompactos.
Para que $f$ sea pseudofinito, es necesario y suficiente que, para todo $\lambda$, la restricción $f_\lambda:f^{-1}(U_\lambda)\to U_\lambda$ de $f$ sea un morfismo pseudofinito.
  \item Supongamos además que $Y$ es noetheriano y $f$ de tipo finito.
Entonces, para que $f$ sea pseudofinito, es necesario y suficiente que, para todo $y\in Y$, el morfismo $f_y=f\times 1:X_y=X\times_Y\operatorname{Spec}(\sh{O}_y)\to\operatorname{Spec}(\sh{O}_y)$ sea pseudofinito.
\end{enumerate}
\end{corollary}

\begin{proof}
(i) Está claro que $f_{(Y')}$ es de tipo finito \sref{IV.1.5.4}. Además, una factorización $X\xrightarrow{g}Z\xrightarrow{u}Y$, con $g$ una inmersión y $u$ finito, induce una factorización $X'\xrightarrow{g'}Z'\xrightarrow{u'}Y'$ de $f_{(Y')}$, donde $Z'=Z_{(Y')}$, $g'=g_{(Y')}$ y $u'=u_{(Y')}$. El morfismo $g'$ es una inmersión \sref[I]{I.4.3.2} y $u'$ es finito \sref[II]{II.6.1.5}; por tanto, $f_{(Y')}$ es pseudofinito.

(ii) La condición es necesaria en virtud de (i), pues los $U_\lambda$ son cuasiseparados al serlo $Y$.
Para ver que es suficiente, observemos (con la notación de \sref{IV.8.12.3}) que, si se pone $X_\lambda=f^{-1}(U_\lambda)$, entonces $\sh{B}|U_\lambda=(f_\lambda)_*(\sh{O}_{X_\lambda})$, y $\sh{C}|U_\lambda$ es la clausura entera de $\sh{O}_{U_\lambda}$ en $\sh{B}|U_\lambda$. En consecuencia, si $h:Z\to Y$ es el morfismo canónico, $Z_\lambda=\operatorname{Spec}(\sh{C}|U_\lambda)$ se identifica con $h^{-1}(U_\lambda)$. Para que $g:X\to Z$ sea una inmersión, es necesario y suficiente que, para todo $\lambda$, su restricción $g_\lambda:f^{-1}(U_\lambda)\to h^{-1}(U_\lambda)$ sea una inmersión \sref[I]{I.4.2.4}. La conclusión se sigue ahora de \sref{IV.8.12.3}.

(iii) Basta, en virtud de (ii), demostrar que, para todo $y\in Y$, existe un entorno $U$ de $y$ tal que la restricción $f^{-1}(U)\to U$ de $f$ sea pseudofinita.
Denotemos por $(U_\lambda)$ el sistema proyectivo filtrante de entornos abiertos afines de $y$, y apliquemos el método de \sref{IV.8.1.2},~$a$)).
Como $Y$ es noetheriano, las restricciones $f_\lambda:f^{-1}(U_\lambda)\to U_\lambda$ de $f$ son de presentación finita, y lo mismo ocurre con $f_y$.
Por hipótesis, $f_y$ se factoriza como $X_y\xrightarrow{g_y}Z_y\xrightarrow{u_y}\operatorname{Spec}(\sh{O}_y)$, donde $u_y$ es finito y $g_y$ una inmersión.
Como $\sh{O}_y$ es noetheriano, también lo es $Z_y$; y como $Z_y$ es de presentación finita sobre $\operatorname{Spec}(\sh{O}_y)$, existen $\lambda$ y un morfismo de presentación finita $u_\lambda:Z_\lambda\to U_\lambda$ tales que $Z_y$ se identifica con $Z_\lambda\times_{U_\lambda}\operatorname{Spec}(\sh{O}_y)$ y $u_y$ con $u_\lambda\times 1$ \sref{IV.8.8.2},~(ii)). Además, existe un morfismo $g_\lambda:X_\lambda\to Z_\lambda$ tal que $g_y=g_\lambda\times 1$ y $f_\lambda=u_\lambda\circ g_\lambda$ \sref{IV.8.8.2},~(i)). Se puede tomar además $\lambda$ de modo que $g_\lambda$ sea una inmersión y $u_\lambda$ un morfismo finito \sref{IV.8.10.5},~(ii) y (x)), lo que demuestra que $f_\lambda$ es pseudofinito.
\end{proof}

\begin{env}[8.12.5]
\label{IV.8.12.5}
Podemos dar ahora del ``Teorema principal'' de Zariski \sref[III]{III.4.4.3}) una demostración que no utiliza los resultados cohomológicos ``globales'' del cap.~III, sino que recurre a las propiedades más finas de los anillos locales noetherianos; generalizaremos además el enunciado del teorema suprimiendo las hipótesis noetherianas:
\end{env}

\begin{theorem}[8.12.6]
\label{IV.8.12.6}
(``Teorema principal'' de Zariski). --- Sea $Y$ un preesquema cuasicompacto y cuasiseparado.
Si un morfismo $f:X\to Y$ es cuasifinito, separado y de presentación finita, existe una factorización de $f$
\begin{equation}
\label{IV.8.12.6.1}
X\xrightarrow{f'}Y'\xrightarrow{u}Y \tag{8.12.6.1}
\end{equation}
donde $f'$ es una inmersión abierta y $u$ un morfismo finito.
\end{theorem}

\oldpage[IV-3]{46}

\begin{proof}
En virtud de \sref{IV.8.12.4},~(ii)) y del carácter local (sobre $Y$) de las nociones de morfismo cuasifinito, separado y de presentación finita, podemos reducirnos al caso en que $Y=\operatorname{Spec}(A)$ es afín.
Aplicando \sref{IV.8.9.1}), podemos suponer que existen un subanillo $A_0$ de $A$, que es una $\bb{Z}$-álgebra de tipo finito, y un $A$-isomorfismo $X_0\otimes_{A_0}A\xrightarrow{\sim}X$, identificándose $f$ mediante este isomorfismo con $f_0\times 1$, donde $f_0:X_0\to\operatorname{Spec}(A_0)$ es un morfismo \emph{de tipo finito}; además, por \sref{IV.8.10.5},~(v) y (xi)), podemos suponer que $f_0$ es separado y cuasifinito; si se demuestra que $f_0$ es pseudofinito, lo mismo ocurrirá con $f$ por \sref{IV.8.12.4},~(i)).
Como $A_0$ es entonces noetheriano y las nociones de morfismo de tipo finito, separado y cuasifinito se conservan por cambio de base, se sigue de \sref{IV.8.12.4},~(iii)) que incluso podemos suponer que $A$ es un anillo local, \emph{esencialmente de tipo finito sobre} $\bb{Z}$ \sref{IV.1.3.8}).
Pongamos $n=\dim(A)$ y procedamos por inducción sobre $n$; para $n=0$, el teorema es evidente, pues $A$ es un cuerpo y el morfismo $f$ ya es finito \sref[II]{II.6.2.2}).
Pongamos $B=\Gamma(X,\sh{O}_X)$; denotemos por $C$ la clausura entera de $A$ en $B$, pongamos $Z=\operatorname{Spec}(C)$ y sea $g:X\to Z$ el $Y$-morfismo correspondiente al $A$-homomorfismo canónico inyectivo $C\to B$; en virtud de \sref{IV.8.12.3}), basta demostrar que $g$ es una inmersión abierta.
Sea $a$ el punto cerrado de $Y$, y sea $U=Y-\{a\}$; $U$ es noetheriano y todos sus anillos locales son esencialmente de tipo finito sobre $\bb{Z}$ y de dimensión $<n$; teniendo en cuenta la hipótesis de inducción, se ve que la restricción $f^{-1}(U)\to U$ de $f$ es un morfismo pseudofinito.
Se concluye de \sref{IV.8.12.3}) que, si $h:Z\to Y$ es el morfismo estructural, la restricción $f^{-1}(U)\to h^{-1}(U)$ de $g$ es una inmersión abierta.
Pongamos $A'=\hat{A}$, $Y'=\operatorname{Spec}(A')$, $X'=X_{(Y')}$ y $f'=f_{(Y')}:X'\to Y'$.
Como el morfismo canónico $u:Y'\to Y$ es plano, se sigue de \sref{IV.2.3.1}) que $B'=\Gamma(X',\sh{O}_{X'})$ se identifica con la $A'$-álgebra $B\otimes_A A'$.
Por otra parte, como $A$ es un anillo local excelente \sref{IV.7.8.3}), el morfismo $u:Y'\to Y$ es regular y, en consecuencia, \sref{IV.6.14.4}) \emph{la clausura entera de $A'$ en $B'$ es igual a} $C'=C\otimes_A A'$.
Se ve, por tanto, que $Z'=\operatorname{Spec}(C')$ es igual a $Z_{(Y')}$ y que el morfismo $g':X'\to Z'$ procedente de la inyección canónica $C'\to B'$ es igual a $g_{(Y')}$.
Como $u:Y'\to Y$ es fielmente plano y cuasicompacto, para demostrar que $g$ es una inmersión abierta basta demostrar que $g'$ es una inmersión abierta \sref{IV.2.7.1},~(x)).
Ahora bien, $u^{-1}(a)$ consta del punto cerrado $a'$ de $Y'$, y en consecuencia $U'=Y'-\{a'\}=u^{-1}(U)$.
Si $h':Z'\to Y'$ es el morfismo canónico, el hecho de que la restricción $f^{-1}(U)\to h^{-1}(U)$ de $g$ sea una inmersión abierta implica que lo mismo ocurre con la restricción $f'^{-1}(U')\to h'^{-1}(U')$ de $g'$.
El morfismo $f'$ es separado y cuasifinito \sref[II]{II.6.2.4}; como $A'$ es \emph{completo}, se deduce de \sref[II]{II.6.2.6} que $X'$ es $Y'$-isomorfo a una suma $X'_1\sqcup X'_2$, donde la restricción $f'|X'_1=f'_1:X'_1\to Y'$ es un morfismo \emph{finito}, y $X'_2\subset f'^{-1}(U')$.
De ello se sigue que $B'$ es el producto directo de las dos $A'$-álgebras $\Gamma(X'_1,\sh{O}_{X'_1})=B'_1$ y $\Gamma(X'_2,\sh{O}_{X'_2})=B'_2$; se concluye inmediatamente que la clausura entera $C'$ de $A'$ en $B'$ es el producto directo de las clausuras enteras $C'_1$, $C'_2$ de $A'$ en $B'_1$, $B'_2$, respectivamente, de donde $Z'=Z'_1\sqcup Z'_2$, donde $Z'_i=\operatorname{Spec}(C'_i)$ ($i=1,2$); y el morfismo canónico $g':X'\to Z'$ es tal que $g'|X'_i$ es el morfismo canónico $g'_i:X'_i\to Z'_i$ ($i=1,2$).
Pero, como $B'_1$ ya es una $A'$-álgebra finita, se tiene $C'_1=B'_1$, y por tanto $g'_1$ es un isomorfismo.
Por otra parte, como $X'_2\subset f'^{-1}(U')$ y es abierto en $f'^{-1}(U')$, se
\oldpage[IV-3]{47}
sabe ya que $g'_2$ es una inmersión abierta.
Se concluye que $g'$ es, en efecto, una inmersión abierta.
\end{proof}

\begin{remark}[8.12.7]
\label{IV.8.12.7}
Cuando, en \sref{IV.8.12.6}), se supone que $Y$ es un \emph{esquema} afín, la demostración por reducción al caso noetheriano muestra que, en la factorización \sref{IV.8.12.6.1}), los morfismos $f'$ y $u$ son también morfismos \emph{de presentación finita} \sref{IV.1.6.2}).
\end{remark}

\begin{corollary}[8.12.8]
\label{IV.8.12.8}
Sea $Y$ un esquema cuasicompacto tal que existe un $\sh{O}_Y$-Módulo amplio \sref[II]{II.4.5.3}), y sea $f:X\to Y$ un morfismo cuasifinito y cuasiproyectivo.
Entonces existe una factorización de $f$ de la forma
\[
X\xrightarrow{f'}Y'\xrightarrow{u}Y
\]
donde $f'$ es una inmersión abierta y $u$ un morfismo finito.
\end{corollary}

\begin{proof}
La hipótesis implica que $X$ se identifica con un sub-$Y$-esquema cuasicompacto de un $Y$-esquema de la forma $Z=\mathbf{P}^r_Y$ \sref[II]{II.5.3.3}).
Existe, pues, un entorno abierto cuasicompacto $U$ de $X$ en $Z$ tal que $X$ es cerrado en $U$; como $Z$ es un esquema, la inyección canónica $U\to Z$ es un morfismo de presentación finita (\sref{IV.1.2.7} y \sref{IV.1.6.2})), y por tanto el morfismo compuesto $g:U\to Z\to Y$ es también de presentación finita (el hecho de que $\mathbf{P}^r_Y$ sea de presentación finita sobre $Y$ se sigue inmediatamente de la definición \sref[II]{II.4.1.1})).
Sea $\sh{J}$ el Ideal cuasicoherente de $\sh{O}_U$ que define el subpreesquema cerrado $X$; como $U$ es un esquema cuasicompacto, $\sh{J}$ es el límite inductivo filtrante de sus sub-Ideales cuasicoherentes de tipo finito $\sh{J}_\lambda$ \sref[I]{I.9.4.9}).
Si $X_\lambda$ es el subpreesquema cerrado de $U$ definido por $\sh{J}_\lambda$, se tiene entonces $X=\bigcap X_\lambda$.
Para todo $y\in Y$, se tiene por tanto $f^{-1}(y)=\bigcap(X_\lambda\cap g^{-1}(y))$ y, como los conjuntos $X_\lambda\cap g^{-1}(y)$ son cerrados en el espacio noetheriano $g^{-1}(y)$, existe para todo $y$ un índice $\lambda(y)$ tal que $f^{-1}(y)=X_{\lambda(y)}\cap g^{-1}(y)$.
Denotemos por $E_\lambda$ el conjunto de los $y\in Y$ tales que la fibra $X_\lambda\cap g^{-1}(y)$ de la restricción de $g$ a $X_\lambda$ es un $k(y)$-preesquema finito.
La hipótesis de que $f$ es cuasifinito implica, en virtud de lo anterior, que $Y=\bigcup_\lambda E_\lambda$.
Ahora bien, cada uno de los $X_\lambda$ es, por definición, \emph{de presentación finita} sobre $Y$; se sigue, pues, de \sref{IV.9.2.3} y \sref{IV.9.2.6})~($^1$) que los $E_\lambda$ son conjuntos \emph{construibles} en el esquema $Y$; como forman una familia creciente filtrante, existe un índice $\lambda$ tal que $E_\lambda=Y$ \sref{IV.1.9.9}), y para este $\lambda$ el morfismo $f_\lambda:X_\lambda\to Y$, restricción de $g$ a $X_\lambda$, es cuasifinito. Como es de presentación finita y separado, se puede aplicar \sref{IV.8.12.6}), y $f_\lambda$ se factoriza como
\[
X_\lambda\xrightarrow{j_\lambda}Y'_\lambda\xrightarrow{u_\lambda}Y
\]
donde $j_\lambda$ es una inmersión abierta y $u_\lambda$ un morfismo finito.
Como $X$ es un subpreesquema cerrado de $X_\lambda$, hemos demostrado así que $f$ posee la propiedad \sref{IV.8.12.3},~$a$)), de donde se sigue el corolario en virtud de la equivalencia de \sref{IV.8.12.3},~$a$)) y \sref{IV.8.12.3},~$a'$)).
\end{proof}

\footnotetext{($^1$) El lector comprobará que los corolarios \sref{IV.8.12.8} a \sref{IV.8.12.11} no se utilizan en la \textsection~9.}

\begin{corollary}[8.12.9]
\label{IV.8.12.9}
\oldpage[IV-3]{48}
Sea $f:X\to Y$ un morfismo localmente cuasifinito $(\mathbf{Err}_{\text{III}},~20)$.
Para todo $x\in X$, existen un entorno abierto $U$ de $x$ en $X$, un entorno abierto $V$ de $y=f(x)$ en $Y$, tales que $f(U)\subset V$, y una factorización
\[
U\xrightarrow{f'}V'\xrightarrow{u}V
\]
de la restricción de $f$ a $U$, donde $f'$ es una inmersión abierta y $u$ un morfismo finito.
\end{corollary}

\begin{proof}
Basta evidentemente tomar para $V$ un entorno afín de $y$ en $Y$, y para $U$ un entorno afín de $x$ en $X$, contenido en $f^{-1}(V)$ y tal que $f|U$ sea cuasifinito.
Como el morfismo $U\to V$, restricción de $f$, es entonces afín (y por tanto cuasiproyectivo), se puede aplicar \sref{IV.8.12.8}).
\end{proof}

\begin{corollary}[8.12.10]
\label{IV.8.12.10}
Sea $Y$ un preesquema íntegro normal, sea $X$ un preesquema íntegro y sea $f:X\to Y$ un morfismo birracional, localmente cuasifinito $(\mathbf{Err}_{\text{III}},~20)$.
Entonces $f$ es un isomorfismo local; para que $f$ sea una inmersión abierta, es necesario y suficiente que $f$ sea además separado.
\end{corollary}

\begin{proof}
La segunda aserción se sigue inmediatamente de la primera y de \sref[I]{I.8.2.8}).
Para demostrar la primera aserción, podemos suponer $X$ e $Y$ afines y $f$ cuasifinito; consideremos la factorización $f=u\circ f'$ de \sref{IV.8.12.8}, que identifica $X$, mediante $f'$, con el subpreesquema inducido sobre un abierto de $Y'$.
Como $X$ es íntegro, podemos, en virtud de \sref[I]{I.5.2.3}), sustituir $Y'$ por el subpreesquema reducido de $Y'$ cuyo espacio subyacente es $\overline{X}$; podemos, pues, suponer también que $Y'$ es \emph{íntegro}.
Además, como $f$ es birracional, también lo es $u$.
La conclusión se sigue entonces del lema siguiente:
\end{proof}

\begin{lemma}[8.12.10.1]
\label{IV.8.12.10.1}
Sea $Y'$ un preesquema íntegro y sea $Y$ un preesquema íntegro y normal; entonces todo morfismo $u:Y'\to Y$ finito y birracional es un isomorfismo.
\end{lemma}

\begin{proof}
Pongamos $\sh{A}=u_*(\sh{O}_{Y'})$; entonces $\sh{A}$ es una $\sh{O}_Y$-álgebra finita y $Y'$ se identifica con $\operatorname{Spec}(\sh{A})$ \sref[II]{II.1.3.6}.
Si $R(Y)$ es el cuerpo de funciones racionales de $Y$, se tiene, por tanto, para todo $y\in Y$, $\sh{O}_{Y,y}\subset\sh{A}_y\subset R(Y)$; pero, como el anillo $\sh{O}_{Y,y}$ es por hipótesis íntegramente cerrado y $R(Y)$ un cuerpo de fracciones, se tiene necesariamente $\sh{A}_y=\sh{O}_{Y,y}$, de donde el lema.
\end{proof}

\begin{corollary}[8.12.11]
\label{IV.8.12.11}
Sea $Y$ un preesquema íntegro, sea $X$ un preesquema íntegro normal y sea $f:X\to Y$ un morfismo dominante, localmente cuasifinito.
Sean $K$ y $L$ (una extensión de $K$) los cuerpos de funciones racionales de $Y$ y $X$, respectivamente; y sea $Y'$ la clausura entera de $Y$ con respecto a $L$ \sref[II]{II.6.3.4}); entonces $f$ se factoriza de manera única como $f:X\xrightarrow{f'}Y'\xrightarrow{u}Y$, donde $f'$ es birracional y corresponde al automorfismo identidad de $L$; $f'$ es entonces un isomorfismo local, y para que $f'$ sea una inmersión abierta, es necesario y suficiente que $f$ sea separado.
\end{corollary}

\begin{proof}
La existencia y la unicidad de la factorización de $f$ se siguen de \sref[II]{II.6.3.9}).
Se sigue además de \sref[II]{II.6.2.4},~(v)), reduciéndose al caso afín, que $f'$ es localmente cuasifinito; por otra parte, se sigue de \sref[I]{I.5.5.1}) que, para que $f'$ sea separado, es necesario y suficiente que lo sea $f$, ya que $u$ es afín, y por tanto separado; las dos últimas aserciones son, pues, consecuencias de \sref{IV.8.12.10}) aplicado a $f'$.
\end{proof}

\oldpage[IV-3]{49}

\subsection{Traducción en términos de pro-objetos}
\label{subsection:IV.8.13}

La proposición siguiente es esencialmente equivalente a \sref{IV.8.8.2},~(i):

\begin{proposition}[8.13.1]
\label{IV.8.13.1}
Sea $S$ un preesquema y sea $(X_\lambda,v_{\lambda\mu})$ un sistema proyectivo filtrante de $S$-preesquemas; supongamos que existe $\alpha$ tal que $v_{\alpha\lambda}$ es un morfismo afín para todo $\lambda\geqslant\alpha$ (lo que implica por \sref[II]{II.1.6.2} que $v_{\lambda\mu}$ es afín para $\alpha\leqslant\lambda\leqslant\mu$), de modo que el límite proyectivo $X=\varprojlim_\lambda X_\lambda$ existe en la categoría de los $S$-preesquemas \sref{IV.8.2.3}.
Sea $Y$ un $S$-preesquema y, para todo $\lambda\geqslant\alpha$, sea $e_\lambda:\operatorname{Hom}_S(X_\lambda,Y)\to\operatorname{Hom}_S(X,Y)$ la aplicación que a cada $S$-morfismo $f_\lambda:X_\lambda\to Y$ hace corresponder $f=f_\lambda\circ v_\lambda$, donde $v_\lambda:X\to X_\lambda$ es el morfismo canónico.
La familia $(e_\lambda)$ es un sistema inductivo de aplicaciones, que define por tanto una aplicación canónica
\begin{equation}
\label{IV.8.13.1.1}
\varinjlim_\lambda\operatorname{Hom}_S(X_\lambda,Y)\to\operatorname{Hom}_S(X,Y). \tag{8.13.1.1}
\end{equation}

Supongamos que $X_\alpha$ es cuasicompacto (resp. cuasicompacto y cuasiseparado), y que el morfismo estructural $Y\to S$ es localmente de tipo finito (resp. localmente de presentación finita).
Entonces la aplicación \sref{IV.8.13.1.1} es inyectiva (resp. biyectiva).
\end{proposition}

\begin{proof}
Para $\lambda\geqslant\alpha$, pongamos $Z_\lambda=Y\times_S X_\lambda$, de modo que $Z_\lambda=Z_\alpha\times_{X_\alpha}X_\lambda$.
Asimismo, pongamos $Z=Y\times_S X=Z_\alpha\times_{X_\alpha}X$; se sigue entonces de \sref{IV.8.2.5} que, si se pone además $w_{\lambda\mu}=1\times v_{\lambda\mu}:Z_\mu\to Z_\lambda$ para $\alpha\leqslant\lambda\leqslant\mu$ y $w_\lambda=1\times v_\lambda:Z\to Z_\lambda$ para $\alpha\leqslant\lambda$, entonces $Z$ es el límite proyectivo del sistema $(Z_\lambda,w_{\lambda\mu})$, y los $w_\lambda$ son los morfismos canónicos correspondientes.
Además, el morfismo $Z_\alpha\to X_\alpha$ es localmente de tipo finito (resp. localmente de presentación finita) por \sref{IV.1.3.4} y \sref{IV.1.4.3}.
Por último, por \sref[I]{I.3.3.14}, se tiene
\[
\operatorname{Hom}_S(X_\lambda,Y)=\operatorname{Hom}_{X_\lambda}(X_\lambda,Z_\lambda) \quad\text{y}\quad \operatorname{Hom}_S(X,Y)=\operatorname{Hom}_X(X,Z)
\]
Basta ahora aplicar \sref{IV.8.8.2},~(i), tomando $X_\lambda=S_\lambda$ y sustituyendo $Y_\lambda$ por $Z_\lambda$.
\end{proof}

\begin{corollary}[8.13.2]
\label{IV.8.13.2}
Con la notación de \sref{IV.8.13.1}, supongamos que $X_\alpha$ es cuasicompacto y cuasiseparado, y que los $v_{\alpha\lambda}$ son afines para $\alpha\leqslant\lambda$; supongamos además que $Y=\varprojlim_\rho Y_\rho$, donde $(Y_\rho,t_{\rho\sigma})$ es un sistema proyectivo filtrante de $S$-preesquemas tal que, para cada $\rho$, el morfismo estructural $Y_\rho\to S$ es localmente de presentación finita.
Entonces se tiene una biyección canónica
\begin{equation}
\label{IV.8.13.2.1}
\operatorname{Hom}_S(X,Y)\xrightarrow{\sim}\varprojlim_\rho\left(\varinjlim_\lambda\operatorname{Hom}_S(X_\lambda,Y_\rho)\right). \tag{8.13.2.1}
\end{equation}
\end{corollary}

\begin{proof}
En efecto, como $Y$ es el límite proyectivo de los $Y_\rho$, la aplicación canónica $\operatorname{Hom}_S(X,Y)\to\varprojlim_\rho\operatorname{Hom}_S(X,Y_\rho)$ es biyectiva; por otra parte, las hipótesis proporcionan, para cada $\rho$, una biyección canónica $\operatorname{Hom}_S(X,Y_\rho)\xrightarrow{\sim}\varinjlim_\lambda\operatorname{Hom}_S(X_\lambda,Y_\rho)$ por \sref{IV.8.13.1}; de donde la conclusión.
\end{proof}

\begin{env}[8.13.3]
\label{IV.8.13.3}
Los resultados anteriores permiten interpretar en la teoría de preesquemas las nociones de ``provariedad'' y ``proesquema'' que aparecen en ciertas aplicaciones (por ejemplo, en la teoría local de cuerpos de clases siguiendo las ideas de Serre~[39], o en la teoría de N\'eron de la reducción de
\oldpage[IV-3]{50}
variedades abelianas~[32]).
Recordemos aquí brevemente la noción de \emph{pro-objeto} de una categoría, remitiendo al cap.~V para desarrollos más amplios (por lo demás, no utilizaremos antes del cap.~V la interpretación que sigue, y el lector puede por tanto omitirla hasta el final de este número).
Dada una categoría $\mathcal{C}$, la categoría $\mathbf{Pro}(\mathcal{C})$ de los \emph{pro-objetos} de $\mathcal{C}$ tiene por objetos los \emph{sistemas proyectivos} (en el universo considerado) $\mathbf{X}=(X_\mu)_{\mu\in M}$ de objetos de $\mathcal{C}$ cuyos conjuntos de índices (que dependen del sistema proyectivo) se suponen conjuntos preordenados filtrantes.
Dados dos pro-objetos de este tipo $\mathbf{X}=(X_\mu)_{\mu\in M}$ y $\mathbf{X}'=(X'_{\mu'})_{\mu'\in M'}$, los \emph{morfismos} de $\mathbf{X}$ en $\mathbf{X}'$ son, por definición, los elementos del conjunto $\varprojlim_{\mu'}(\varinjlim_\mu\operatorname{Hom}(X_\mu,X'_{\mu'}))$; es inmediato que estos conjuntos pueden servir como conjuntos de morfismos. La composición de sistemas de morfismos $u_{\mu'}^\mu:X_\mu\to X'_{\mu'}$ y $u_{\mu''}^{\prime\mu'}:X'_{\mu'}\to X''_{\mu''}$, inductivos respecto del índice superior y proyectivos respecto del índice inferior, se efectúa ``argumento por argumento'', es decir, considerando el sistema $u_{\mu''}^{\prime\mu'}\circ u_{\mu'}^\mu$.
\end{env}

\begin{env}[8.13.4]
\label{IV.8.13.4}
Sea $S$ un preesquema cuasicompacto y cuasiseparado, y denotemos por $\mathcal{C}$ la subcategoría plena de la categoría $(\mathbf{Sch})_{/S}$ de los $S$-preesquemas, formada por los $S$-preesquemas $X$ que poseen la propiedad siguiente: el morfismo estructural $X\to S$ se factoriza como $X\xrightarrow{g}X_0\xrightarrow{f}S$, donde $g:X\to X_0$ es afín y $f:X_0\to S$ es \emph{de presentación finita}; diremos abreviadamente que los preesquemas de $\mathcal{C}$ son \emph{esencialmente afines sobre} $S$.
Consideremos, por otra parte, la subcategoría plena $\mathcal{C}'_0$ de $(\mathbf{Sch})_{/S}$ formada por los $S$-preesquemas \emph{de presentación finita}, y la categoría $\mathbf{Pro}(\mathcal{C}'_0)$ de pro-objetos de $\mathcal{C}'_0$.
Diremos que un objeto $\mathbf{X}=(X_\mu)_{\mu\in M}$ de $\mathbf{Pro}(\mathcal{C}'_0)$ es \emph{esencialmente afín} si existe $\gamma\in M$ tal que, para $\mu\geqslant\gamma$, el morfismo de transición $v_{\gamma\mu}:X_\mu\to X_\gamma$ es afín (lo que implica que, para $\gamma\leqslant\nu\leqslant\mu$, $v_{\nu\mu}:X_\mu\to X_\nu$ es afín).
Obsérvese que un objeto de $\mathbf{Pro}(\mathcal{C}'_0)$ isomorfo a un objeto esencialmente afín \emph{no es} necesariamente él mismo esencialmente afín.
Denotaremos por $\mathcal{C}'$ la subcategoría plena de $\mathbf{Pro}(\mathcal{C}'_0)$ formada por los pro-objetos esencialmente afines.

Dicho esto, se sigue de \sref{IV.8.2.2} y \sref{IV.8.2.3} que, para todo objeto $\mathbf{X}=(X_\mu)_{\mu\in M}$ de $\mathcal{C}'$, existe el $S$-preesquema $X=\varprojlim_\mu X_\mu$; además, como el morfismo $X\to X_\mu$ es afín para todo $\mu$ suficientemente grande \sref{IV.8.2.2}, $X$ es, por definición, \emph{esencialmente afín sobre} $S$.
Pongamos $X=L(\mathbf{X})$; demostremos que se ha definido así un \emph{funtor} canónico
\begin{equation}
\label{IV.8.13.4.1}
L:\mathcal{C}'\to\mathcal{C} \tag{8.13.4.1}
\end{equation}

Para dos objetos $\mathbf{X}=(X_\mu)$ y $\mathbf{X}'=(X'_{\mu'})$ de $\mathcal{C}'$, se tiene, para cada $\mu'$, una aplicación canónica
\[
\varinjlim_\mu\operatorname{Hom}_S(X_\mu,X'_{\mu'})\to\operatorname{Hom}_S(\varprojlim_\mu X_\mu,X'_{\mu'})
\]
definida en \sref{IV.8.13.1.1}; por otra parte, por definición del límite proyectivo, existe una biyección canónica
\[
\varprojlim_{\mu'}\operatorname{Hom}_S(\varprojlim_\mu X_\mu,X'_{\mu'})\xrightarrow{\sim}\operatorname{Hom}_S(\varprojlim_\mu X_\mu,\varprojlim_{\mu'}X'_{\mu'})
\]
\oldpage[IV-3]{51}

de donde se obtiene una aplicación canónica
\begin{equation}
\label{IV.8.13.4.2}
\varprojlim_{\mu'}\varinjlim_\mu\operatorname{Hom}_S(X_\mu,X'_{\mu'})\to\operatorname{Hom}_S(\varprojlim_\mu X_\mu,\varprojlim_{\mu'}X'_{\mu'}) \tag{8.13.4.2}
\end{equation}
evidentemente funtorial en $\mathbf{X}$ y $\mathbf{X}'$, y que completa la definición del funtor $L$.
\end{env}

\begin{proposition}[8.13.5]
\label{IV.8.13.5}
Bajo las hipótesis y con la notación de \sref{IV.8.13.4}, el funtor $L$ es plenamente fiel.
Si además $S$ es un preesquema noetheriano (lo que implica ya que $S$ es cuasicompacto y cuasiseparado \sref{IV.1.2.8}), $L$ es \emph{una equivalencia de categorías}.
\end{proposition}

\begin{proof}
Decir que $L$ es plenamente fiel significa que la aplicación \sref{IV.8.13.4.2} es biyectiva, lo que es un caso particular de \sref{IV.8.13.2}.
En efecto, como los morfismos estructurales $X_\mu\to S$ son de presentación finita, son en particular cuasicompactos y cuasiseparados; por tanto, los $X_\mu$ son cuasicompactos y cuasiseparados.

Para demostrar que, cuando $S$ es noetheriano, $L$ es una equivalencia de categorías, basta, puesto que ya sabemos que $L$ es plenamente fiel, probar que todo $S$-preesquema esencialmente afín $X$ es $S$-\emph{isomorfo} a un objeto de la forma $L(\mathbf{X})$, donde $\mathbf{X}\in\mathcal{C}'$ $(\mathbf{0}_{\text{III}},~8.1.5)$.
Ahora bien, por hipótesis existe una factorización $X\xrightarrow{g}X_0\xrightarrow{f}S$ del morfismo estructural, siendo $f$ de presentación finita y $g$ afín.
Podemos escribir, por tanto, $X=\operatorname{Spec}(\sh{A})$, donde $\sh{A}$ es una $\sh{O}_{X_0}$-Álgebra cuasicoherente \sref[II]{II.1.3.1}.
Ahora bien, como $X_0$ es noetheriano (pues es de tipo finito sobre $S$ y $S$ es noetheriano), $\sh{A}$ es el límite inductivo filtrante de la familia $(\sh{A}_\lambda)$ de sus sub-$\sh{O}_{X_0}$-Álgebras cuasicoherentes de tipo finito \sref[I]{I.9.6.6}.
Pongamos $X_\lambda=\operatorname{Spec}(\sh{A}_\lambda)$; los morfismos $X_\lambda\to X_0$ son de tipo finito, y por tanto de presentación finita puesto que $X_0$ es noetheriano; en consecuencia, lo mismo ocurre con los morfismos compuestos $X_\lambda\to X_0\to S$ \sref{IV.1.6.2}; dicho de otro modo, los $X_\lambda$ pertenecen a $\mathcal{C}'_0$, y como los morfismos $X_\lambda\to X_0$ son afines, $\mathbf{X}=(X_\lambda)$ es un objeto de $\mathcal{C}'$ cuyo límite proyectivo existe y es $S$-isomorfo a $X$ por \sref{IV.8.2.2}.
Esto termina la demostración.
\end{proof}

\begin{remark}[8.13.6]
\label{IV.8.13.6}
Se sigue de \sref{IV.1.6.2} y \sref[II]{II.1.6.2} que, si $X$ e $Y$ son esencialmente afines sobre $S$, lo mismo ocurre con $X\times_S Y$.
Se concluye, por ejemplo, $(\mathbf{0}_{\text{III}},~8.2.5)$ que un $\mathcal{C}$-grupo no es otra cosa que un $(\mathbf{Sch})_{/S}$-grupo que es un preesquema esencialmente afín sobre $S$.
Por otra parte, los productos finitos existen en la categoría $\mathcal{C}'$: en efecto, si $\mathbf{X}=(X_\mu)_{\mu\in M}$, $\mathbf{Y}=(Y_\rho)_{\rho\in R}$ son dos objetos de $\mathcal{C}'$, los productos $X_\mu\times_S Y_\rho$ son $S$-preesquemas de presentación finita y, tomando como morfismos de transición $X_\nu\times_S Y_\sigma\to X_\mu\times_S Y_\rho$ los productos de los morfismos de transición $X_\nu\to X_\mu$ e $Y_\sigma\to Y_\rho$, se ve inmediatamente que $(X_\mu\times_S Y_\rho)$ es un producto de $\mathbf{X}$ y $\mathbf{Y}$ en $\mathbf{Pro}(\mathcal{C}'_0)$; además, por \sref[II]{II.1.6.2}, los morfismos de transición así definidos son afines para $\mu$ y $\rho$ suficientemente grandes, y por tanto el producto $\mathbf{X}\times\mathbf{Y}$ pertenece también a $\mathcal{C}'$.
Se concluye entonces, como antes, que un $\mathcal{C}'$-grupo es un $\mathbf{Pro}(\mathcal{C}'_0)$-grupo que es esencialmente afín.
Se deduce de \sref{IV.8.13.5} que las categorías de los $\mathcal{C}'$-grupos y de los $\mathcal{C}$-grupos son \emph{equivalentes} cuando $S$ es \emph{noetheriano}.
Parece plausible que, cuando $S$ es el espectro de un cuerpo $k$, la categoría de los $\mathcal{C}$-grupos sea \emph{equivalente} a la de los $\boldsymbol{K}$-\emph{grupos}, donde $\boldsymbol{K}$ es la categoría de los $S$-preesquemas \emph{cuasicompactos}; dicho de otro modo, todo preesquema en grupos cuasicompacto sobre $k$ sería esencialmente afín.
Por otra parte, si se denota
\oldpage[IV-3]{52}
por $\mathcal{C}'_0$-$\boldsymbol{gr}$ la categoría de los $\mathcal{C}'_0$-grupos, es muy probable que la categoría de los $\mathcal{C}'$-grupos sea equivalente a la subcategoría plena de $\mathbf{Pro}(\mathcal{C}'_0\text{-}\boldsymbol{gr})$ formada por los ``grupos proalgebraicos esencialmente afines'', es decir, los sistemas proyectivos $\mathbf{G}=(G_\mu)_{\mu\in M}$, donde los $G_\mu$ son grupos algebraicos sobre $k$ y los morfismos de transición $G_\nu\to G_\mu$ son afines para $\mu$ suficientemente grande (lo que también puede expresarse diciendo que $\mathbf{G}$ es una extensión de un grupo algebraico por un grupo proalgebraico afín).
La conjunción de estas dos conjeturas es además equivalente a la afirmación siguiente: todo preesquema en grupos cuasicompacto sobre $k$ es una extensión de un ``grupo algebraico'' (es decir, un preesquema en grupos de tipo finito sobre $k$) por un preesquema en grupos \emph{afín} sobre $k$.
Como los únicos grupos proalgebraicos encontrados hasta ahora en la práctica son de hecho esencialmente afines, sería ventajoso sustituir el estudio de los grupos proalgebraicos generales (introducidos y estudiados por Serre~[40]) por el de los esquemas en grupos cuasicompactos sobre $k$, cuya definición es conceptualmente mucho más sencilla.
\end{remark}

\subsection{Caracterización de un preesquema localmente de presentación finita sobre otro mediante el funtor que representa}
\label{subsection:IV.8.14}

\begin{env}[8.14.1]
\label{IV.8.14.1}
Dado un preesquema $S$, diremos, como en \sref{IV.8.13.4}, que un sistema proyectivo filtrante $(X_\lambda,v_{\lambda\mu})$ de $S$-preesquemas es \emph{esencialmente afín} si existe $\alpha$ tal que $v_{\alpha\lambda}$ es un morfismo afín para $\lambda\geqslant\alpha$.

El enunciado siguiente, que será especialmente útil en el cap.~V, precisa \sref{IV.8.8.2},~(i) proporcionando un recíproco:
\end{env}

\begin{proposition}[8.14.2]
\label{IV.8.14.2}
Sea $S$ un preesquema y sea $f:X\to S$ un morfismo.
Para todo $S$-preesquema $T$, pongamos
\[
h_X(T)=\operatorname{Hom}_S(T,X)
\]
de modo que $h_X$ es un funtor contravariante de la categoría $(\mathbf{Sch})_{/S}$ de los $S$-preesquemas a la categoría $\mathbf{Ens}$ de los conjuntos $(\mathbf{0}_{\text{III}},~8.1.1)$, y $X$ es un objeto que representa este funtor $(\mathbf{0}_{\text{III}},~8.1.8)$.
Las condiciones siguientes son equivalentes:
\begin{enumerate}[label=\emph{\alph*)}]
  \item $f$ es localmente de presentación finita.
  \item Para todo sistema proyectivo filtrante $(Z_\lambda)$ de $S$-preesquemas, esencialmente afín \sref{IV.8.13.4} y formado por preesquemas cuasicompactos y cuasiseparados, la aplicación canónica \sref{IV.8.13.1.1}
\begin{equation}
\label{IV.8.14.2.1}
\varinjlim_\lambda h_X(Z_\lambda)\to h_X(\varprojlim_\lambda Z_\lambda) \tag{8.14.2.1}
\end{equation}
es biyectiva.
  \item Para todo sistema proyectivo filtrante $(Z_\lambda)$ de $S$-preesquemas tal que los $Z_\lambda$ sean esquemas afines, la aplicación \sref{IV.8.14.2.1} es biyectiva.
  \item[\rm c$'$)] Para todo abierto afín $U$ de $S$ y todo sistema proyectivo filtrante $(Z_\lambda)$ de $U$-preesquemas tal que los $Z_\lambda$ sean esquemas afines, la aplicación \sref{IV.8.14.2.1} es biyectiva.
\end{enumerate}
\end{proposition}

\begin{proof}
El hecho de que \emph{a)} entrañe \emph{b)} es precisamente \sref{IV.8.13.1}; es trivial que \emph{b)} entraña \emph{c)} y que \emph{c)} entraña \emph{c$'$)}.
Queda ver que \emph{c$'$)} entraña \emph{a)} y, como la propiedad \emph{a)} es local sobre $S$, podemos reducirnos al caso en que $S$ es \emph{afín}.
\oldpage[IV-3]{53}

Supongamos primero que $X$ es también afín; la aserción que hay que demostrar es entonces equivalente a la siguiente:

\begin{corollary}[8.14.2.2]
\label{IV.8.14.2.2}
Sea $A$ un anillo y sea $B$ una $A$-álgebra.
Para que, para todo sistema inductivo filtrante $(C_\lambda)$ de $A$-álgebras, la aplicación canónica
\begin{equation}
\label{IV.8.14.2.3}
\varinjlim_\lambda\operatorname{Hom}_{A\text{-}\mathrm{alg.}}(B,C_\lambda)\to\operatorname{Hom}_{A\text{-}\mathrm{alg.}}(B,\varinjlim_\lambda C_\lambda) \tag{8.14.2.3}
\end{equation}
sea biyectiva, es necesario y suficiente que $B$ sea una $A$-álgebra de presentación finita.
\end{corollary}

Solo queda demostrar que la condición es necesaria.
Tomemos primero para $(C_\lambda)$ el sistema inductivo filtrante de las sub-$A$-álgebras de tipo finito de $B$, de modo que $\varinjlim_\lambda C_\lambda=B$.
El hecho de que \sref{IV.8.14.2.3} sea biyectiva implica, en particular, que la aplicación identidad $1_B:B\to B$ se factoriza como $B\to C_\lambda\to B$ para un $\lambda$ adecuado, lo que implica que $C_\lambda=B$; por tanto, $B$ es una $A$-álgebra \emph{de tipo finito}.
Escribamos $B=C/\mathfrak{J}$, donde $C=A[T_1,\ldots,T_n]$ y $\mathfrak{J}$ es un ideal de $C$. El ideal $\mathfrak{J}$ es el límite inductivo filtrante de sus ideales de tipo finito $\mathfrak{J}_\lambda\subset\mathfrak{J}$; poniendo $C_\lambda=C/\mathfrak{J}_\lambda$ y utilizando la exactitud de los límites inductivos filtrantes, se ve que $B$ es nuevamente isomorfo al límite inductivo del sistema inductivo filtrante $(C_\lambda)$.
Existen, por tanto, $\lambda$ y un $A$-homomorfismo $u:B\to C_\lambda$ tales que el compuesto $B\xrightarrow{u}C_\lambda\xrightarrow{p_\lambda}B$ (donde $p_\lambda$ es el homomorfismo canónico) es la identidad.
Sea $q_\lambda:C\to C_\lambda$ el homomorfismo canónico, y pongamos
$t_i=p_\lambda(q_\lambda(T_i))$.
Entonces $p_\lambda(u(t_i))=p_\lambda(q_\lambda(T_i))$, es decir,
$u(t_i)-q_\lambda(T_i)\in\mathfrak{J}/\mathfrak{J}_\lambda$.
Existe, pues, $\mu\geqslant\lambda$ tal que estos $n$ elementos pertenecen a
$\mathfrak{J}_\mu/\mathfrak{J}_\lambda$ ($1\leqslant i\leqslant n$).
Si $p_{\mu\lambda}:C_\lambda\to C_\mu$ es el homomorfismo canónico, entonces
$p_{\mu\lambda}(u(t_i))=p_{\mu\lambda}(q_\lambda(T_i))=q_\mu(T_i)$.
Sustituyendo $\lambda$ por $\mu$ y $u$ por $p_{\mu\lambda}\circ u$, podemos por tanto
suponer que $u(t_i)=q_\lambda(T_i)$ para todo $i$; y si
$r=p_\lambda\circ q_\lambda$ es el homomorfismo canónico
$C\to C/\mathfrak{J}=B$, podemos escribir $u(r(T_i))=q_\lambda(T_i)$ para todo $i$,
de donde $q_\lambda=u\circ r$.
Esto implica necesariamente que $\mathfrak{J}=\mathfrak{J}_\lambda$: si $z\in\mathfrak{J}$, entonces $r(z)=0$, y por tanto $q_\lambda(z)=u(r(z))=0$ y $z\in\mathfrak{J}_\lambda$. Así, $B=C_\lambda$.

Pasemos ahora al caso en que $S$ es afín y $X$ arbitrario; lo que hay que demostrar es que todo abierto afín $V$ de $X$ es de presentación finita sobre $S$, y en virtud de lo que acabamos de demostrar, basta probar que, para todo sistema proyectivo filtrante $(Z_\lambda)$ de $S$-preesquemas afines, la aplicación
\begin{equation}
\label{IV.8.14.2.4}
\varinjlim_\lambda\operatorname{Hom}_S(Z_\lambda,V)\to\operatorname{Hom}_S(\varprojlim_\lambda Z_\lambda,V) \tag{8.14.2.4}
\end{equation}
es biyectiva.
Es inmediato que esta aplicación es inyectiva, pues si $(v_\lambda)$, $(v'_\lambda)$ son dos sistemas inductivos de $S$-homomorfismos $v_\lambda:Z_\lambda\to V$, $v'_\lambda:Z_\lambda\to V$ tales que los morfismos correspondientes
\[
Z\xrightarrow{u_\lambda}Z_\lambda\xrightarrow{v_\lambda}V,\qquad
Z\xrightarrow{u_\lambda}Z_\lambda\xrightarrow{v'_\lambda}V
\]
son iguales (siendo $u_\lambda$ el morfismo canónico), entonces los morfismos
\[
Z\xrightarrow{u_\lambda}Z_\lambda\xrightarrow{v_\lambda}V\xrightarrow{j}X,\qquad Z\xrightarrow{u_\lambda}Z_\lambda\xrightarrow{v'_\lambda}V\xrightarrow{j}X
\]
(donde $j$ es la inyección canónica) son iguales, lo que implica por hipótesis $j\circ v_\lambda=j\circ v'_\lambda$ para un $\lambda$ adecuado, y por tanto $v_\lambda=v'_\lambda$.

\oldpage[IV-3]{54}

Queda demostrar que \sref{IV.8.14.2.4} es sobreyectiva.
Sea $v:Z\to V$ un $S$-morfismo.
Por hipótesis, existen un índice $\lambda$ y un $S$-morfismo
$w_\lambda:Z_\lambda\to X$ tales que $j\circ v$ se factoriza como
\[
 Z\xrightarrow{u_\lambda}Z_\lambda\xrightarrow{w_\lambda}X.
\]
Basta demostrar que existe $\mu\geqslant\lambda$ tal que
\[
 Z_\mu\xrightarrow{u_{\lambda\mu}}Z_\lambda\xrightarrow{w_\lambda}X
\]
(donde $u_{\lambda\mu}$ es el morfismo de transición) se factoriza como
\[
 Z_\mu\xrightarrow{v_\mu}V\xrightarrow{j}X.
\]
Para todo $\mu\geqslant\lambda$, pongamos
$U_\mu=u_{\lambda\mu}^{-1}(w_\lambda^{-1}(V))$.
Entonces
\[
 u_\mu^{-1}(U_\mu)=u_\lambda^{-1}(U_\lambda)
 =u_\lambda^{-1}(w_\lambda^{-1}(V))
 =v^{-1}(V)=Z.
\]
Los $Z_\lambda$ son cuasicompactos y los $U_\mu$, por ser abiertos, son
ind-construibles \sref{IV.1.9.6}; por tanto, \sref{IV.8.3.4} proporciona un
$\mu\geqslant\lambda$ para el cual $U_\mu=Z_\mu$.
\end{proof}

\begin{remark}[8.14.3]
\label{IV.8.14.3}
El hecho de que la aplicación \sref{IV.8.14.2.1} sea inyectiva cuando $f$ es localmente de tipo
finito \sref{IV.8.8.2},~(i) conduce naturalmente a preguntarse si este resultado admite un recíproco.
No es así, ni siquiera cuando $S$ y $X$ son afines, pues existen monomorfismos
$X\to S$ que no son de tipo finito \sref[I]{I.2.4.2}, y estos refutan tal
conjetura.
\end{remark}
